\documentclass[11pt,a4paper]{article}
\usepackage[utf8]{inputenc}
\usepackage[vietnamese,english]{babel}
\usepackage{amsmath,amssymb}
\usepackage{geometry}
\geometry{top=2cm,bottom=2cm,left=2cm,right=2cm}
\usepackage{tcolorbox}
\usepackage{tikz}
\usetikzlibrary{shapes,arrows,positioning}
\usepackage{enumitem}
\usepackage{hyperref}
\usepackage{booktabs}
\usepackage{tabularx}

\hypersetup{
    colorlinks=true,
    linkcolor=blue,
    urlcolor=blue,
    pdftitle={CSD203 - Cấu trúc dữ liệu và giải thuật - Ultra Edition},
    pdfauthor={CSD203 Team - FPT University}
}

% Định nghĩa môi trường cho câu hỏi
\newcounter{questioncount}
\newenvironment{question}[1]{%
  \par\vspace{1em}\noindent\textbf{\large Câu #1}\par\vspace{0.4em}
}{%
  \par\vspace{1em}\hrule\vspace{0.6em}
}

% Định nghĩa môi trường choices
\newenvironment{choices}{%
  \begin{enumerate}[label=\textbf{\Alph*.},leftmargin=2.2em,itemsep=2pt,topsep=3pt]
}{%
  \end{enumerate}
}
\newcommand{\choice}{\item}

% Định nghĩa môi trường codebox
\newenvironment{codebox}{%
  \begin{tcolorbox}[colback=gray!8!white,colframe=gray!50!black,arc=1.5mm,boxrule=0.6pt,left=3mm,right=3mm,top=2.5mm,bottom=2.5mm]
  \ttfamily\small
}{%
  \end{tcolorbox}
}

% Định nghĩa môi trường tóm tắt chuyên đề
\newenvironment{summarybox}[1]{%
  \begin{tcolorbox}[colback=blue!4!white,colframe=blue!70!black,arc=2mm,boxrule=0.8pt,title=\textbf{#1},fonttitle=\normalsize\bfseries,left=3.5mm,right=3.5mm,top=3mm,bottom=3mm]
}{%
  \end{tcolorbox}
}

\begin{document}

\begin{center}
  {\Huge \textbf{CSD203: DATA STRUCTURES \& ALGORITHMS}}\\[0.5em]
  {\LARGE \textbf{ULTRA EDITION -- BỘ NGÂN HÀNG CÂU HỎI TOÀN DIỆN}}\\[0.5em]
  {\Large \textbf{Phân loại chuyên sâu theo 7 chuyên đề trọng tâm \& Giải thuật bổ trợ}}\\[0.6em]
  {\normalsize \textbf{Nguồn tài liệu chuẩn hóa: 276 câu Ngân hàng CSD203 + 100 câu Đề thi FE/RE SU26}}\\[1em]
  \rule{\textwidth}{1.5pt}
\end{center}

\vspace{0.5em}
\noindent\textbf{BẢNG TỔNG QUAN PHÂN LOẠI NỘI DUNG (376 CÂU HỎI TOÀN KHÓA):}
\begin{center}
\small
\begin{tabular}{|c|l|c|p{7.8cm}|}
\hline
\textbf{STT} & \textbf{Chuyên đề nội dung} & \textbf{Số câu} & \textbf{Trọng tâm kiến thức chuẩn thi} \\
\hline
1 & \textbf{Mảng (Arrays \& Lists)} & 25 & Truy xuất chỉ mục ($O(1)$), Slicing âm/dương, Phép gán, Nối mảng, Mảng đa chiều. \\
\hline
2 & \textbf{Danh sách liên kết (Linked Lists)} & 30 & Singly/Doubly/Circular List, Thao tác Thêm/Xóa/Duyệt nút, Con trỏ head, tail, next, prev. \\
\hline
3 & \textbf{Ngăn xếp (Stacks)} & 19 & Nguyên lý LIFO, push/pop, LifoQueue, Đảo chuỗi, Biểu thức tiền tố (Prefix) \& hậu tố (Postfix). \\
\hline
4 & \textbf{Hàng đợi (Queues)} & 16 & Nguyên lý FIFO, enqueue/dequeue, Hàng đợi mảng vòng (Circular Queue), Deque, Priority Queue. \\
\hline
5 & \textbf{Bảng băm (Hash Tables)} & 54 & Hàm băm, Xung đột, Chaining, Địa chỉ mở (Linear/Quadratic Probing, Double Hashing), Load Factor. \\
\hline
6 & \textbf{Cây tìm kiếm nhị phân \& Đồ thị} & 112 & \textbf{6.1 Trees \& BST (62 câu):} Duyệt cây, BST, AVL Tree, Heap.\\
  & \textbf{(Binary Search Trees \& Graphs)} & & \textbf{6.2 Graphs (50 câu):} BFS/DFS, Dijkstra, Prim, Kruskal, Topological Sort. \\
\hline
7 & \textbf{Thuật toán sắp xếp (Sorting)} & 29 & Bubble, Selection, Insertion, Quick Sort, Merge Sort, Radix Sort, Cận dưới so sánh \& Độ ổn định. \\
\hline
8 & \textbf{Các giải thuật bổ trợ} & 91 & Tìm kiếm nhị phân/tuyến tính, Đệ quy \& Đệ quy đuôi, Phân tích Big-O, Mã hóa Huffman, Thuật toán KMP. \\
\hline
\multicolumn{2}{|c|}{\textbf{TỔNG CỘNG}} & \textbf{376} & \textbf{100\% câu hỏi có đáp án chi tiết và giải thích chuyên sâu chuẩn xác} \\
\hline
\end{tabular}
\end{center}

\bigskip
\tableofcontents
\bigskip
\rule{\textwidth}{0.8pt}

\newpage
\section*{Chương 1: Mảng (Arrays, Lists, Slicing \& Operations)}
\addcontentsline{toc}{section}{Chương 1: Mảng (Arrays, Lists, Slicing \& Operations)}

\begin{summarybox}{Tóm tắt lý thuyết trọng tâm: Mảng \& Danh sách Python}
\begin{itemize}[noitemsep,topsep=2pt]
  \item \textbf{Cơ chế lưu trữ:} Mảng lưu các phần tử liên tiếp trong bộ nhớ. Truy xuất theo chỉ mục có độ phức tạp $O(1)$.
  \item \textbf{Cắt lát (Slicing):} Cú pháp \texttt{list[start:stop:step]}. Nếu chỉ số âm, $-1$ là phần tử cuối cùng. Lưu ý slicing tạo ra bản sao nông (shallow copy), phép gán trên bản sao không thay đổi danh sách gốc trừ khi dùng slice assignment (\texttt{A[i:j] = ...}).
  \item \textbf{Phương thức danh sách:} \texttt{append(x)} (thêm 1 phần tử), \texttt{extend(iterable)} (nối mảng), \texttt{pop(index)} (xóa và trả về phần tử tại vị trí index, mặc định là cuối), \texttt{remove(value)} (xóa giá trị đầu tiên khớp), \texttt{len(A)} ($O(1)$).
\end{itemize}
\end{summarybox}

\begin{question}{1.1}
% [ID gốc: 1.1]
% [Nguồn: Ngân hàng CSD203 (Câu #8)]
What is the output of the following list operation?
\begin{codebox}
aList = [10, 20, 30, 40, 50, 60, 70, 80]
print(aList[2:5])
\end{codebox}
\begin{choices}
\choice \texttt{[30, 40, 50, 60]}
\choice \texttt{[20, 30, 40]}
\choice \texttt{[30, 40, 50]}
\choice \texttt{[20, 30, 40, 50]}
\end{choices}
\textbf{Đáp án:} C
\textbf{Giải thích:} Cắt lát \texttt{aList[2:5]} trích xuất các phần tử tại vị trí 2, 3, 4 trong danh sách \texttt{[10, 20, 30, 40, 50, 60, 70, 80]}, cho ra \texttt{[30, 40, 50]}.
\end{question}

\begin{question}{1.2}
% [ID gốc: 1.3]
% [Nguồn: Ngân hàng CSD203 (Câu #19)]
What is the output of the following program?
\begin{codebox}
a=[1,2,3,4,5]
a[0]=6;
print(len(a))
\end{codebox}
\begin{choices}
\choice 0
\choice 5
\choice 6
\choice 1
\end{choices}
\textbf{Đáp án:} B
\textbf{Giải thích:} Lệnh gán \texttt{a[0] = 6} chỉ thay thế giá trị phần tử đầu tiên mà không thay đổi kích thước danh sách, nên \texttt{len(a)} vẫn là 5.
\end{question}

\begin{question}{1.3}
% [ID gốc: 1.4]
% [Nguồn: Ngân hàng CSD203 (Câu #26)]
What will be the output of the following code?
\begin{codebox}
a=[1, 2, 3, 4, 5]
b=a[2:4]
print(b)
\end{codebox}
\begin{choices}
\choice \texttt{[3, 4, 5]}
\choice \texttt{[2, 4]}
\choice \texttt{[2, 3, 4]}
\choice \texttt{[3, 4]}
\end{choices}
\textbf{Đáp án:} D
\textbf{Giải thích:} Phép cắt lát \texttt{b = a[2:4]} trên mảng \texttt{[1, 2, 3, 4, 5]} trích xuất phần tử tại chỉ số 2 và 3, kết quả là \texttt{[3, 4]}.
\end{question}

\begin{question}{1.4}
% [ID gốc: 1.7]
% [Nguồn: Ngân hàng CSD203 (Câu #40)]
List A is defined as follows: \texttt{A = [1, 2, 3, 4, 5]} Select all of the following statements that remove the middle element 3 from A so that it equals \texttt{[1, 2, 4, 5]} :
\begin{choices}
\choice \texttt{A[2:2] = []}
\choice \texttt{A[2].remove()}
\choice \texttt{del A[2]}
\choice \texttt{A[2] = []}
\end{choices}
\textbf{Đáp án:} C
\textbf{Giải thích:} Để xóa phần tử ở giữa danh sách Python \texttt{A = [1, 2, 3, 4, 5]}, có thể sử dụng lệnh \texttt{del A[2]} hoặc \texttt{A.pop(2)}.
\end{question}

\begin{question}{1.5}
% [ID gốc: 1.8]
% [Nguồn: Ngân hàng CSD203 (Câu #42)]
What will be the output of the following code?
\begin{codebox}
b=[1] *5
print(b)
\end{codebox}
\begin{choices}
\choice \texttt{[1, 1, 1, 1, 1]}
\choice 5
\choice 0
\choice Error
\end{choices}
\textbf{Đáp án:} A
\textbf{Giải thích:} Biểu thức \texttt{b = [1] * 5} tạo ra một danh sách chứa 5 phần tử đều có giá trị 1: \texttt{[1, 1, 1, 1, 1]}.
\end{question}

\begin{question}{1.6}
% [ID gốc: 1.11]
% [Nguồn: Ngân hàng CSD203 (Câu #48)]
What will be the output of the following code?
\begin{codebox}
S = []
S.append(1)
S.append(2)
S.append(3)
print(S.peak())
\end{codebox}
\begin{choices}
\choice \texttt{[1,2,3]}
\choice 1
\choice 3
\choice Error
\end{choices}
\textbf{Đáp án:} D
\textbf{Giải thích:} Thuật toán KMP (Knuth-Morris-Pratt) tối ưu hóa việc tìm kiếm chuỗi mẫu nhờ việc tính toán trước bảng tiền tố (mảng LPS) để tránh lùi con trỏ văn bản.
\end{question}

\begin{question}{1.7}
% [ID gốc: 1.15]
% [Nguồn: Ngân hàng CSD203 (Câu #65)]
What is the output of the following program?
\begin{codebox}
A = [1, 2, 3, 4, 5]
A.pop(1)
print (A)
\end{codebox}
\begin{choices}
\choice 1
\choice \texttt{[1, 3, 4, 5]}
\choice \texttt{[2, 3, 4, 5]}
\choice \texttt{[1, 2, 3, 4, 5]}
\end{choices}
\textbf{Đáp án:} B
\textbf{Giải thích:} Phương thức \texttt{A.pop(1)} xóa và trả về phần tử tại chỉ mục 1 của danh sách (giá trị 2). Sau khi xóa, danh sách \texttt{A} còn lại \texttt{[1, 3, 4, 5]}.
\end{question}

\begin{question}{1.8}
% [ID gốc: 1.16]
% [Nguồn: Ngân hàng CSD203 (Câu #72)]
List a is defined as follows: \texttt{A = ['a', 'b', 'c']} Which of the following statements adds 'd' and 'e' to the end of A, so that it then equals \texttt{['a', 'b', 'c', 'd', 'e']}:
\begin{choices}
\choice \texttt{A.extend(['d', 'e'])}
\choice \texttt{A.append(['d', 'e'])}
\choice \texttt{A[-1:] = ['d', 'e']}
\choice \texttt{A.append('d', 'e')}
\end{choices}
\textbf{Đáp án:} A
\textbf{Giải thích:} Để thêm nhiều phần tử vào cuối danh sách Python, ta dùng phương thức \texttt{extend(['d', 'e'])} hoặc toán tử \texttt{+ ['d', 'e']}.
\end{question}

\begin{question}{1.9}
% [ID gốc: 1.23]
% [Nguồn: Ngân hàng CSD203 (Câu #93)]
What will be the output of the following code?
\begin{codebox}
b=[1, 2, 3, 4, 5]
print(b[-3:-1])
\end{codebox}
\begin{choices}
\choice \texttt{[1]}
\choice \texttt{[2, 3]}
\choice \texttt{[3, 4]}
\choice Error
\end{choices}
\textbf{Đáp án:} C
\textbf{Giải thích:} Cắt lát âm \texttt{b[-3:-1]} trên \texttt{[1, 2, 3, 4, 5]} lấy các phần tử từ chỉ số -3 (giá trị 3) đến trước chỉ số -1 (giá trị 5), cho kết quả \texttt{[3, 4]}.
\end{question}

\begin{question}{1.10}
% [ID gốc: 1.24]
% [Nguồn: Ngân hàng CSD203 (Câu #94)]
What is the output of the following code?
\begin{codebox}
a = "Hello, World!"
print(a[1])
\end{codebox}
\begin{choices}
\choice H
\choice e
\choice Hello
\choice World
\end{choices}
\textbf{Đáp án:} B
\textbf{Giải thích:} Chỉ mục chuỗi \texttt{a[1]} trên chuỗi \texttt{'Hello, World!'} lấy ký tự thứ hai là \texttt{'e'}.
\end{question}

\begin{question}{1.11}
% [ID gốc: 1.25]
% [Nguồn: Ngân hàng CSD203 (Câu #96)]
What is the output of the following program?
\begin{codebox}
a=[5, 4, 7, 2, 9]
print(a.index(4))
\end{codebox}
\begin{choices}
\choice 0
\choice 1
\choice 2
\choice 9
\end{choices}
\textbf{Đáp án:} B
\textbf{Giải thích:} Phương thức \texttt{a.index(4)} trong mảng \texttt{[5, 4, 7, 2, 9]} trả về vị trí xuất hiện đầu tiên của giá trị 4, tức chỉ mục 1.
\end{question}

\begin{question}{1.12}
% [ID gốc: 1.27]
% [Nguồn: Ngân hàng CSD203 (Câu #98)]
What will be the output of the following code?
\begin{codebox}
b=[1] + [5, 8]
print(b)
\end{codebox}
\begin{choices}
\choice \texttt{[1]}
\choice \texttt{[5, 8]}
\choice \texttt{[1, 5, 8]}
\choice Error
\end{choices}
\textbf{Đáp án:} C
\textbf{Giải thích:} Phép nối danh sách \texttt{[1] + [5, 8]} tạo ra danh sách mới kết hợp cả hai: \texttt{[1, 5, 8]}.
\end{question}

\begin{question}{1.13}
% [ID gốc: 1.28]
% [Nguồn: Ngân hàng CSD203 (Câu #99)]
What is the output of the following list assignment?
\begin{codebox}
aList = [4, 8, 12, 16]
aList[1:4] = [20, 24, 28]
print(aList)
\end{codebox}
\begin{choices}
\choice \texttt{[4, 20, 24, 28, 8, 12, 16]}
\choice \texttt{[4, 20, 24, 28]}
\choice \texttt{[4, 8, 12, 16]}
\choice \texttt{[20, 24, 28]}
\end{choices}
\textbf{Đáp án:} B
\textbf{Giải thích:} Gán lát cắt \texttt{aList[1:4] = [20, 24, 28]} thay thế 3 phần tử tại chỉ số 1, 2, 3 bằng các giá trị mới tương ứng.
\end{question}

\begin{question}{1.14}
% [ID gốc: 1.30]
% [Nguồn: Ngân hàng CSD203 (Câu #111)]
What will be the output of the following code?
\begin{codebox}
S=[]
for i in range(1,5):
    if (i%2==1):
        S.append(i)
while (len(S) > 0):
    print(S.pop(0))
\end{codebox}
\begin{choices}
\choice 1 3
\choice 1 3 5
\choice 3 1
\choice 5 3 1
\end{choices}
\textbf{Đáp án:} A
\textbf{Giải thích:} Phép cắt lát (slicing) trong Python tạo ra một bản sao nông (shallow copy) của danh sách ban đầu. Việc gán lại phần tử \texttt{b[0] = 99} chỉ thay đổi trên danh sách \texttt{b} mà không làm thay đổi danh sách \texttt{a} ban đầu. Do đó \texttt{a} vẫn giữ nguyên là \texttt{[10, 20, 30, 40]}.
\end{question}

\begin{question}{1.15}
% [ID gốc: 1.31]
% [Nguồn: Ngân hàng CSD203 (Câu #130)]
What is the output of the following program?
\begin{codebox}
txt = "1546a"
x = txt.isascii()
print(x)
\end{codebox}
\begin{choices}
\choice Error
\choice 1
\choice True
\choice False
\end{choices}
\textbf{Đáp án:} C
\textbf{Giải thích:} Phương thức \texttt{txt.isascii()} kiểm tra xem tất cả các ký tự trong chuỗi có thuộc bảng mã ASCII hay không, '1546a' toàn ký tự ASCII nên trả về True.
\end{question}

\begin{question}{1.16}
% [ID gốc: 1.33]
% [Nguồn: Ngân hàng CSD203 (Câu #133)]
What will be the output of the following code?
\begin{codebox}
x = hex(12);
print(x)
\end{codebox}
\begin{choices}
\choice 0xa
\choice 0xb
\choice 0xc
\choice 0xd
\end{choices}
\textbf{Đáp án:} C
\textbf{Giải thích:} Hàm \texttt{hex(12)} trong Python chuyển số nguyên 12 sang dạng chuỗi hexa tiền tố '0x', kết quả là \texttt{'0xc'}.
\end{question}

\begin{question}{1.17}
% [ID gốc: 1.35]
% [Nguồn: Ngân hàng CSD203 (Câu #140)]
What does the function \texttt{S.discard(e)} do?
\begin{choices}
\choice Remove element e from the set, if present.
\choice Add element e to the set.
\choice Return True if the set contains element e
\choice Generate an iteration of all elements of the set
\end{choices}
\textbf{Đáp án:} A
\textbf{Giải thích:} Phương thức \texttt{S.discard(e)} trên tập hợp (Set) trong Python xóa phần tử \texttt{e} nếu có mặt mà không báo lỗi nếu \texttt{e} không tồn tại.
\end{question}

\begin{question}{1.18}
% [ID gốc: 1.50]
% [Nguồn: Ngân hàng CSD203 (Câu #207)]
What is the output of the following program?
\begin{codebox}
txt = "acbd"
x = txt.istitle()
print(x)
\end{codebox}
\begin{choices}
\choice Error
\choice 1
\choice True
\choice False
\end{choices}
\textbf{Đáp án:} D
\textbf{Giải thích:} Phương thức \texttt{txt.istitle()} kiểm tra chuỗi có viết hoa chữ cái đầu của mỗi từ hay không; 'acbd' toàn chữ thường nên trả về False.
\end{question}

\begin{question}{1.19}
% [ID gốc: 1.51]
% [Nguồn: Ngân hàng CSD203 (Câu #208)]
What is the output of the following code?
\begin{codebox}
my_list = ["Hello", "Python"]
print("-".join(my_list))
\end{codebox}
\begin{choices}
\choice HelloPython-
\choice Hello-Python
\choice -HelloPython
\choice HelloPython
\end{choices}
\textbf{Đáp án:} B
\textbf{Giải thích:} Phương thức \texttt{'-'.join(['Hello', 'Python'])} nối các chuỗi trong danh sách bằng dấu gạch nối, kết quả là \texttt{'Hello-Python'}.
\end{question}

\begin{question}{1.20}
% [ID gốc: 1.60]
% [Nguồn: Ngân hàng CSD203 (Câu #248)]
What will be the output of the following code?
\begin{codebox}
list = [1, 2, 3]
list.append([4, 5])
print(len(list))
\end{codebox}
\begin{choices}
\choice 5
\choice 4
\choice 3
\choice 2
\end{choices}
\textbf{Đáp án:} B
\textbf{Giải thích:} Lệnh \texttt{list.append([4, 5])} thêm cả danh sách \texttt{[4, 5]} làm một phần tử đơn lẻ vào cuối \texttt{[1, 2, 3]}, độ dài danh sách mới là 4.
\end{question}

\begin{question}{1.21}
% [ID gốc: 1.61]
% [Nguồn: Ngân hàng CSD203 (Câu #249)]
What is the output of the following list operation?
\begin{codebox}
list = [1, 3, 5, 2, 0, 7, 4]
print(list[:2])
\end{codebox}
\begin{choices}
\choice \texttt{[1, 3]}
\choice \texttt{[1, 3, 5]}
\choice \texttt{[1]}
\choice \texttt{[5, 2, 0, 7, 4]}
\end{choices}
\textbf{Đáp án:} A
\textbf{Giải thích:} Cắt lát \texttt{list[:2]} trên \texttt{[1, 3, 5, 2, 0, 7, 4]} lấy các phần tử từ đầu đến chỉ số trước 2, kết quả là \texttt{[1, 3]}.
\end{question}

\begin{question}{1.22}
% [ID gốc: 1.63]
% [Nguồn: Ngân hàng CSD203 (Câu #251)]
What is the output of the following program?
\begin{codebox}
txt = "50800a"
x = txt.isdigit()
print(x)
\end{codebox}
\begin{choices}
\choice Error
\choice 1
\choice True
\choice False
\end{choices}
\textbf{Đáp án:} D
\textbf{Giải thích:} Phương thức \texttt{txt.isdigit()} kiểm tra chuỗi có chứa toàn chữ số hay không; chuỗi '50800a' chứa ký tự 'a' nên trả về False.
\end{question}

\begin{question}{1.23}
% [ID gốc: 1.82]
% [Nguồn: Đề thi FE/RE CSD203 (Câu #68)]
% Quiz: CSD203_SU26_RE - Question #18 (ID: d366e605-d83e-4df1-9bd4-4054a221c8a8)
What will be the output of the following code?
\begin{codebox}
x = [1, 2, 2, 3, 4, 3, 9]
y = x.count(3)
print(y)
\end{codebox}
\begin{choices}
\choice 4
\choice 3
\choice 2
\choice 1
\end{choices}
\textbf{Đáp án:} C
\textbf{Giải thích:} Phương thức count(3) đếm số lần xuất hiện của giá trị 3 trong danh sách x. Trong [1, 2, 2, 3, 4, 3, 9], số 3 xuất hiện 2 lần.
\end{question}

\begin{question}{1.24}
% [ID gốc: 1.85]
% [Nguồn: Đề thi FE/RE CSD203 (Câu #87)]
% Quiz: CSD203_SU26_RE - Question #37 (ID: a9396772-d5fb-4fe3-8315-b3a630fd9c71)
What will be the output of the following code snippet?
\begin{codebox}
my_list = [1, 2, 3]
my_list.append([4, 5])
print(len(my_list))
\end{codebox}
\begin{choices}
\choice 3
\choice 4
\choice 5
\choice 6
\end{choices}
\textbf{Đáp án:} B
\textbf{Giải thích:} Phương thức append([4, 5]) chèn cả danh sách [4, 5] làm một phần tử duy nhất vào cuối my_list. Danh sách mới là [1, 2, 3, [4, 5]] có độ dài 4.
\end{question}

\begin{question}{1.25}
% [ID gốc: 1.89]
% [Nguồn: Đề thi FE/RE CSD203 (Câu #97)]
% Quiz: CSD203_SU26_RE - Question #47 (ID: 89a430cc-496c-4173-93d0-990943e17843)
What will be the output of the following code?
\begin{codebox}
b = [1, 2, 3, 4, 5]
print(b[-3:-1])
\end{codebox}
\begin{choices}
\choice [1]
\choice [2, 3]
\choice [3, 4]
\choice Error
\end{choices}
\textbf{Đáp án:} C
\textbf{Giải thích:} Cắt lát (slicing) với chỉ số âm b[-3:-1]: chỉ số -3 ứng với phần tử giá trị 3, chỉ số -1 ứng với phần tử giá trị 5 nhưng không được bao gồm. Kết quả trích xuất là [3, 4].
\end{question}

\newpage
\section*{Chương 2: Danh sách liên kết (Linked Lists)}
\addcontentsline{toc}{section}{Chương 2: Danh sách liên kết (Linked Lists)}

\begin{summarybox}{Tóm tắt lý thuyết trọng tâm: Danh sách liên kết (Singly, Doubly, Circular)}
\begin{itemize}[noitemsep,topsep=2pt]
  \item \textbf{Cấu trúc Node:} Mỗi nút gồm vùng chứa dữ liệu (\texttt{data}) và một hoặc hai con trỏ liên kết (\texttt{next}, \texttt{prev}). Bộ nhớ cấp phát động, không cần kích thước cố định ban đầu.
  \item \textbf{Singly Linked List (Đơn):} Mỗi nút chỉ trỏ tới nút kế tiếp. Thêm/xóa ở đầu: $O(1)$. Thêm/xóa ở cuối khi chỉ có con trỏ \texttt{head}: $O(n)$ do phải duyệt từ đầu.
  \item \textbf{Doubly Linked List (Đôi):} Mỗi nút có con trỏ \texttt{next} và \texttt{prev}. Nút đầu có \texttt{head.prev = None}, nút cuối có \texttt{tail.next = None}. Cho phép duyệt 2 chiều linh hoạt.
  \item \textbf{Circularly Linked List (Vòng):} Con trỏ \texttt{next} của nút cuối cùng trỏ ngược lại nút \texttt{head} (không có nút nào trỏ tới \texttt{None}).
\end{itemize}
\end{summarybox}

\begin{question}{2.1}
% [ID gốc: 1.5]
% [Nguồn: Ngân hàng CSD203 (Câu #29)]
What does the following code do?
\begin{codebox}
def Operation(self, newdata):
    NewNode = Node(newdata)
    if self.head is None:
        self.head = NewNode
        return
    laste = self.head
    while(laste.next):
        laste = laste.next
    laste.next=NewNode
\end{codebox}
\begin{choices}
\choice Insert a new node at the tail of the list
\choice Remove the last node
\choice Traverse the list
\choice Insert a new node at the beginning of the list
\end{choices}
\textbf{Đáp án:} A
\textbf{Giải thích:} Đoạn mã duyệt đến cuối danh sách liên kết qua vòng lặp \texttt{while laste.next} rồi gán \texttt{laste.next = NewNode}, tức chèn nút mới vào cuối danh sách.
\end{question}

\begin{question}{2.2}
% [ID gốc: 1.6]
% [Nguồn: Ngân hàng CSD203 (Câu #32)]
Which is the correct code to access data of the second node in a linked list?
\begin{choices}
\choice \texttt{head.next.data}
\choice \texttt{head.next.next.data}
\choice \texttt{head.head.data}
\choice \texttt{head.data}
\end{choices}
\textbf{Đáp án:} A
\textbf{Giải thích:} Trong danh sách liên kết đơn có con trỏ \texttt{head} trỏ tới nút đầu tiên, con trỏ \texttt{head.next} trỏ tới nút thứ hai. Do đó, cú pháp chính xác để truy cập dữ liệu của nút thứ hai là \texttt{head.next.data}.
\end{question}

\begin{question}{2.3}
% [ID gốc: 1.10]
% [Nguồn: Ngân hàng CSD203 (Câu #46)]
In Linked List implementation, a node carries information regarding\_\_\_\_\_\_\_\_
\begin{choices}
\choice Data and Link
\choice Node
\choice Link
\choice Data
\end{choices}
\textbf{Đáp án:} A
\textbf{Giải thích:} Trong cài đặt Danh sách liên kết, mỗi nút (node) gồm hai thành phần cơ bản: phần dữ liệu (\texttt{data}) chứa thông tin/giá trị lưu trữ và phần liên kết (\texttt{link} / \texttt{next}) chứa con trỏ tham chiếu đến nút kế tiếp.
\end{question}

\begin{question}{2.4}
% [ID gốc: 1.12]
% [Nguồn: Ngân hàng CSD203 (Câu #50)]
Suppose a singly linked list of integers is given below: head ->1->2->3->4->5->None What will be the output of the following code?
\begin{codebox}
def listprint(self):
    printval = self.head
    while (printval):
        print (printval.data)
        printval = printval.next
\end{codebox}
\begin{choices}
\choice 1 2 3 4 5
\choice 2 3 4 5
\choice 1 2 3 4
\choice 5 4 3 2 1
\end{choices}
\textbf{Đáp án:} A
\textbf{Giải thích:} Khởi đầu \texttt{p} trỏ nút 1. Vòng 1: \texttt{p.next} (2) và \texttt{p.next.next} (3) khác None $\to$ \texttt{p} nhảy 2 bước đến nút 3. Vòng 2: \texttt{p.next} (4) và \texttt{p.next.next} (5) khác None $\to$ \texttt{p} nhảy 2 bước đến nút 5. Đến nút 5 thì \texttt{p.next} là None nên vòng lặp dừng. Giá trị \texttt{p.data} tại thời điểm dừng bằng 5.
\end{question}

\begin{question}{2.5}
% [ID gốc: 1.17]
% [Nguồn: Ngân hàng CSD203 (Câu #74)]
You are given pointers to first node of a singly linked list, which of the following operations are dependent on the length of the linked list?
\begin{choices}
\choice Delete the first element
\choice Insert a new element as a first element
\choice Add a new element at the end of the list
\choice None of these
\end{choices}
\textbf{Đáp án:} C
\textbf{Giải thích:} Trong danh sách liên kết đơn chỉ có con trỏ trỏ đến nút đầu tiên (\texttt{head}), để xóa nút cuối cùng (\texttt{delete last node}) ta phải duyệt qua toàn bộ $n$ nút từ đầu đến nút áp chót ($O(n)$ thao tác). Do đó thao tác này phụ thuộc trực tiếp vào độ dài của danh sách liên kết.
\end{question}

\begin{question}{2.6}
% [ID gốc: 1.18]
% [Nguồn: Ngân hàng CSD203 (Câu #81)]
Which is the correct code to acess data of node p in a linked list?
\begin{choices}
\choice \texttt{p.data}
\choice \texttt{p.next}
\choice \texttt{p.next.data}
\choice \texttt{head.p.data}
\end{choices}
\textbf{Đáp án:} A
\textbf{Giải thích:} Trong cấu trúc nút của danh sách liên kết, trường lưu trữ dữ liệu thông thường được đặt tên là \texttt{data} (hoặc \texttt{val}). Vì vậy để truy cập dữ liệu của nút \texttt{p}, cú pháp chuẩn là \texttt{p.data}.
\end{question}

\begin{question}{2.7}
% [ID gốc: 1.21]
% [Nguồn: Ngân hàng CSD203 (Câu #91)]
How to access to the value part of a singly linked list node x?
\begin{choices}
\choice \texttt{x.value}
\choice \texttt{x.next}
\choice \texttt{x.next.value}
\choice \texttt{x.value.next}
\end{choices}
\textbf{Đáp án:} A
\textbf{Giải thích:} Trong các bài toán và hiện thực danh sách liên kết phổ biến (như trên LeetCode / Python), phần giá trị của nút \texttt{x} được truy xuất thông qua thuộc tính \texttt{x.val}.
\end{question}

\begin{question}{2.8}
% [ID gốc: 1.22]
% [Nguồn: Ngân hàng CSD203 (Câu #92)]
What does the following code do?
\begin{codebox}
x = 0
p = head
while (p != none):
    x = x + 1
    p = p.next
\end{codebox}
\begin{choices}
\choice Insert a new node to the front of a linked list
\choice Calculate the size of a linked list
\choice Traverse the linked list
\choice Insert a new node to the end of the linked list
\end{choices}
\textbf{Đáp án:} B
\textbf{Giải thích:} Hàm bắt đầu từ nút \texttt{p}, mỗi lần bước qua nút kế tiếp \texttt{p = p.next} thì tăng biến đếm \texttt{c} lên 1 cho đến khi \texttt{p is None}. Đây là giải thuật đếm tổng số lượng nút (kích thước) của danh sách liên kết.
\end{question}

\begin{question}{2.9}
% [ID gốc: 1.26]
% [Nguồn: Ngân hàng CSD203 (Câu #97)]
Given a pointer h points to the first node of a singly linked list. Which one is used to remove the list?
\begin{choices}
\choice \texttt{h = h.next}
\choice \texttt{h = null}
\choice \texttt{h.remove}
\choice \texttt{h.delete}
\end{choices}
\textbf{Đáp án:} B
\textbf{Giải thích:} Trong danh sách liên kết đơn, con trỏ \texttt{h} trỏ tới nút đầu tiên (\texttt{head}). Danh sách rỗng khi và chỉ khi con trỏ đầu trỏ tới \texttt{None} (tức \texttt{h == None}).
\end{question}

\begin{question}{2.10}
% [ID gốc: 1.32]
% [Nguồn: Ngân hàng CSD203 (Câu #131)]
What does the following code do?
\begin{codebox}
class Node:
    def __init__(self, data):
        self.data = data
        self.next = None
        self.prev = None
\end{codebox}
\begin{choices}
\choice Create a node of singly linked list
\choice Create a node of circular linked list
\choice Create a node of doubly linked list
\choice None of these
\end{choices}
\textbf{Đáp án:} C
\textbf{Giải thích:} Hàm duyệt tuần tự từ \texttt{head} đến nút cuối cùng qua con trỏ \texttt{p.next} và in ra giá trị \texttt{p.data} của từng nút. Đây là hàm duyệt (traversal) hiển thị danh sách liên kết.
\end{question}

\begin{question}{2.11}
% [ID gốc: 1.37]
% [Nguồn: Ngân hàng CSD203 (Câu #159)]
Given a pointer tail points to the last node of a non-empty singly linked list. Which one is used to add new node x to the end of the list?
\begin{choices}
\choice \texttt{tail = x;}
\choice \texttt{tail.next = x; tail = x}
\choice \texttt{tail.next = x}
\choice \texttt{x= tail.next; tail = x}
\end{choices}
\textbf{Đáp án:} B
\textbf{Giải thích:} Cấu trúc dữ liệu Union-Find (Disjoint Set Union) hỗ trợ hai thao tác \texttt{find} và \texttt{union} với thời gian gần như hằng số nhờ kỹ thuật nén đường đi.
\end{question}

\begin{question}{2.12}
% [ID gốc: 1.38]
% [Nguồn: Ngân hàng CSD203 (Câu #164)]
How to delete the first element of a non-empty linked list that is pointed by pointer Head?
\begin{choices}
\choice \texttt{Head = None;}
\choice \texttt{Head = tail;}
\choice \texttt{Head = Head.next}
\choice \texttt{Head = tail.next}
\end{choices}
\textbf{Đáp án:} C
\textbf{Giải thích:} Ma trận khoảng cách trong thuật toán Floyd-Warshall được khởi tạo bằng trọng số của các cạnh trực tiếp nối giữa các đỉnh.
\end{question}

\begin{question}{2.13}
% [ID gốc: 1.39]
% [Nguồn: Ngân hàng CSD203 (Câu #165)]
What is the primary benefit of using a circularly linked list?
\begin{choices}
\choice It allows circular traversal without reaching the end of the list.
\choice It allows easier removal of the first element.
\choice It eliminates the need for a tail pointer.
\choice It supports multiple pointers to the same node.
\end{choices}
\textbf{Đáp án:} A
\textbf{Giải thích:} Phép duyệt Postorder trên cây nhị phân thường được sử dụng để xóa toàn bộ các nút trong cây hoặc tính toán kích thước thư mục.
\end{question}

\begin{question}{2.14}
% [ID gốc: 1.41]
% [Nguồn: Ngân hàng CSD203 (Câu #170)]
What does the following algorithm do?
\begin{codebox}
Algorithm add_last(L,e):
newest = Node(e)
newest.next = None
L.tail.next = newest
L.tail = newest
\end{codebox}
\begin{choices}
\choice Inserting a new node at the end of a singly linked list.
\choice Inserting a new node at the head of a singly linked list.
\choice Removing a node at the end of a singly linked list.
\choice Removing a node at the head of a singly linked list.
\end{choices}
\textbf{Đáp án:} A
\textbf{Giải thích:} Đoạn mã tạo nút mới và gán \texttt{L.tail.next = newest; L.tail = newest}, đây chính là thao tác chèn nút mới vào cuối danh sách liên kết đơn có con trỏ tail.
\end{question}

\begin{question}{2.15}
% [ID gốc: 1.49]
% [Nguồn: Ngân hàng CSD203 (Câu #201)]
A linked list is a series of connected
\begin{choices}
\choice ADTs
\choice Nodes
\choice vectors
\choice head
\end{choices}
\textbf{Đáp án:} B
\textbf{Giải thích:} Danh sách liên kết (Linked List) là một cấu trúc dữ liệu tuyến tính gồm một chuỗi các nút (Nodes) được liên kết với nhau bằng các con trỏ.
\end{question}

\begin{question}{2.16}
% [ID gốc: 1.52]
% [Nguồn: Ngân hàng CSD203 (Câu #217)]
In a circularly linked list, what does the next pointer of the last node point to?
\begin{choices}
\choice NULL
\choice The head node
\choice The previous node
\choice Itself
\end{choices}
\textbf{Đáp án:} B
\textbf{Giải thích:} Trong danh sách liên kết vòng (Circularly Linked List), con trỏ \texttt{next} của nút cuối cùng (tail) luôn trỏ ngược lại nút đầu tiên (head) của danh sách.
\end{question}

\begin{question}{2.17}
% [ID gốc: 1.53]
% [Nguồn: Ngân hàng CSD203 (Câu #218)]
What is the main advantage of a doubly linked list over a singly linked list?
\begin{choices}
\choice Easier insertion at the beginning of the list
\choice Ability to traverse the list in both directions
\choice Reduced memory usage
\choice Faster access to elements by index
\end{choices}
\textbf{Đáp án:} B
\textbf{Giải thích:} Ưu điểm lớn nhất của danh sách liên kết đôi (Doubly Linked List) so với danh sách liên kết đơn là khả năng duyệt danh sách theo cả hai chiều (tiến và lùi) nhờ có hai con trỏ \texttt{next} và \texttt{prev}.
\end{question}

\begin{question}{2.18}
% [ID gốc: 1.54]
% [Nguồn: Ngân hàng CSD203 (Câu #220)]
In a doubly linked list, what does the prev pointer of the head node point to?
\begin{choices}
\choice The next node in the list
\choice The previous node in the list
\choice It points to NULL
\choice It points to the last node in the list
\end{choices}
\textbf{Đáp án:} C
\textbf{Giải thích:} Trong danh sách liên kết đôi, nút đầu danh sách (head) không có nút đứng trước, do đó con trỏ \texttt{prev} của nó luôn trỏ tới \texttt{NULL} (hoặc \texttt{None}).
\end{question}

\begin{question}{2.19}
% [ID gốc: 1.58]
% [Nguồn: Ngân hàng CSD203 (Câu #246)]
How is memory utilization in a Singly Linked List compared to an array?
\begin{choices}
\choice Singly Linked Lists use more memory due to additional pointers.
\choice Singly Linked Lists use less memory as elements are stored in contiguous memory locations.
\choice Singly Linked Lists and arrays use the same amount of memory.
\choice Memory utilization depends on the specific implementation.
\end{choices}
\textbf{Đáp án:} A
\textbf{Giải thích:} Danh sách liên kết đơn tiêu tốn nhiều bộ nhớ hơn mảng vì ngoài dữ liệu, mỗi nút còn phải lưu trữ thêm con trỏ tham chiếu đến nút kế tiếp.
\end{question}

\begin{question}{2.20}
% [ID gốc: 1.59]
% [Nguồn: Ngân hàng CSD203 (Câu #247)]
What is the correct sequence of steps to add a new node to the end of a doubly linked list in Python?
\begin{choices}
\choice Traverse the list to the end, update the last node's next pointer, and set the new node's previous pointer.
\choice Create a new node, set its data and pointers, and update the head pointer.
\choice Create a new node, set its data and pointers, and update the tail pointer.
\choice Traverse the list to the end, update the last node's previous pointer, and set the new node's previous pointer.
\choice None of the others.
\end{choices}
\textbf{Đáp án:} C
\textbf{Giải thích:} Để thêm nút mới vào cuối danh sách liên kết đôi: tạo nút mới, thiết lập dữ liệu và con trỏ liên kết, sau đó cập nhật con trỏ \texttt{tail} trỏ tới nút mới.
\end{question}

\begin{question}{2.21}
% [ID gốc: 1.62]
% [Nguồn: Ngân hàng CSD203 (Câu #250)]
Suppose a singly linked list of integers is given below: head ->1->2->3->4->5->None What will be the output of the following code?
\begin{codebox}
def listprint(self):
    printval = self.head
    while printval.next is not None:
        print (printval.data)
        printval = printval.next
\end{codebox}
\begin{choices}
\choice 1 2 3 4 5
\choice 2 3 4 5
\choice 1 2 3 4
\choice 5 4 3 2 1
\end{choices}
\textbf{Đáp án:} C
\textbf{Giải thích:} Vòng lặp \texttt{while printval.next is not None:} dừng lại khi con trỏ tới nút cuối cùng (nút 5) mà không in giá trị của nút 5, do đó chỉ in ra các giá trị 1, 2, 3, 4.
\end{question}

\begin{question}{2.22}
% [ID gốc: 1.71]
% [Nguồn: Đề thi FE/RE CSD203 (Câu #14)]
% Quiz: CSD203_SU26_FE - Question #14 (ID: 38db38d1-07a0-43a0-bb15-9bb3eeee3bab)
In Linked List implementation, a node carries information regarding:
\begin{choices}
\choice Data
\choice Link
\choice Data and Link
\choice Node
\end{choices}
\textbf{Đáp án:} C
\textbf{Giải thích:} Một nút trong danh sách liên kết bao gồm hai thành phần cơ bản: phần lưu trữ dữ liệu (Data) và phần lưu địa chỉ liên kết tới nút tiếp theo (Link).
\end{question}

\begin{question}{2.23}
% [ID gốc: 1.72]
% [Nguồn: Đề thi FE/RE CSD203 (Câu #18)]
% Quiz: CSD203_SU26_FE - Question #18 (ID: ab645e55-1d12-4e1a-aed6-a90019adccdb)
How to remove an element x after element y of a non-empty singly linked list?
\begin{choices}
\choice y.next = x.next
\choice x.next = y.next
\choice head.next = x
\choice head.next = y
\end{choices}
\textbf{Đáp án:} A
\textbf{Giải thích:} Để loại bỏ phần tử x nằm ngay sau phần tử y trong danh sách liên kết đơn, ta cần cập nhật con trỏ next của y trỏ thẳng đến phần tử đứng sau x, tức là y.next = x.next.
\end{question}

\begin{question}{2.24}
% [ID gốc: 1.73]
% [Nguồn: Đề thi FE/RE CSD203 (Câu #21)]
% Quiz: CSD203_SU26_FE - Question #21 (ID: b6cc5478-436e-4510-ac7a-f0e524b77a85)
You are given a pointer to the first node of a singly linked list. Which of the following operations is dependent on the length of the linked list?
\begin{choices}
\choice Delete the first element
\choice Insert a new element as a first element
\choice Add a new element at the end of the list
\choice None of these
\end{choices}
\textbf{Đáp án:} C
\textbf{Giải thích:} Thêm một phần tử vào cuối danh sách liên kết đơn (khi chỉ có con trỏ đầu) đòi hỏi phải duyệt qua toàn bộ danh sách, do đó thời gian phụ thuộc vào độ dài n.
\end{question}

\begin{question}{2.25}
% [ID gốc: 1.74]
% [Nguồn: Đề thi FE/RE CSD203 (Câu #22)]
% Quiz: CSD203_SU26_FE - Question #22 (ID: 3920a17a-a97a-4a21-b490-5e51c0ef7e06)
What additional memory does a doubly linked list require per node compared to a singly linked list?
\begin{choices}
\choice 1 pointer
\choice 2 pointers
\choice 3 pointers
\choice 0 pointer
\end{choices}
\textbf{Đáp án:} A
\textbf{Giải thích:} Danh sách liên kết đôi yêu cầu thêm một con trỏ trỏ về nút trước đó (prev) trên mỗi nút so với danh sách liên kết đơn.
\end{question}

\begin{question}{2.26}
% [ID gốc: 1.75]
% [Nguồn: Đề thi FE/RE CSD203 (Câu #23)]
% Quiz: CSD203_SU26_FE - Question #23 (ID: b3e710f6-735f-4a2a-ab45-9ae965e94881)
In a circularly linked list, what does the next pointer of the last node point to?
\begin{choices}
\choice NULL
\choice The head node
\choice The previous node
\choice Itself
\end{choices}
\textbf{Đáp án:} B
\textbf{Giải thích:} Trong danh sách liên kết vòng (circularly linked list), nút cuối cùng chứa con trỏ next trỏ ngược lại về nút đầu tiên (head node).
\end{question}

\begin{question}{2.27}
% [ID gốc: 1.77]
% [Nguồn: Đề thi FE/RE CSD203 (Câu #34)]
% Quiz: CSD203_SU26_FE - Question #34 (ID: c8c1d3b7-b365-477b-b5c7-7067f11292c4)
In a circularly linked list with n nodes, how can you efficiently check if the list is circular when given access to the head node?
\begin{choices}
\choice Traverse the entire list and check if any node's next pointer points back to the head
\choice Use two pointers, where one moves faster than the other, and check if they meet
\choice Count the number of nodes in the list and compare it to a predefined value
\choice Check the tail node directly to see if its next pointer points to the head
\end{choices}
\textbf{Đáp án:} D
\textbf{Giải thích:} Trong trường hợp danh sách vòng chuẩn, việc kiểm tra tính chất vòng có thể xác định trực tiếp bằng cách nhìn vào liên kết của nút cuối cùng (tail node) trỏ về head, hoặc qua việc duy trì con trỏ tail.
\end{question}

\begin{question}{2.28}
% [ID gốc: 1.80]
% [Nguồn: Đề thi FE/RE CSD203 (Câu #63)]
% Quiz: CSD203_SU26_RE - Question #13 (ID: c8696ee4-d285-4400-b877-696835611e79)
What does the following code do?
\begin{codebox}
def Operation(self, newdata):
    NewNode = Node(newdata)
    if self.head is None:
        self.head = NewNode
        return
    laste = self.head
    while laste.next:
        laste = laste.next
    laste.next = NewNode
\end{codebox}
\begin{choices}
\choice Insert a new node at the tail of the list
\choice Insert a new node at the beginning of the list
\choice Traverse the list
\choice Remove the last node
\end{choices}
\textbf{Đáp án:} A
\textbf{Giải thích:} Đoạn mã duyệt đến nút cuối cùng của danh sách liên kết bằng vòng while laste.next, sau đó gán laste.next = NewNode, tức chèn nút mới vào cuối (tail) danh sách.
\end{question}

\begin{question}{2.29}
% [ID gốc: 1.81]
% [Nguồn: Đề thi FE/RE CSD203 (Câu #66)]
% Quiz: CSD203_SU26_RE - Question #16 (ID: beacb838-dc78-4425-8dbb-bc7e90fe0a43)
Suppose a singly linked list of integers is given below:
$$\text{head} \to 1 \to 2 \to 3 \to 4 \to 5 \to \text{None}$$
What will be the output of the following code?
\begin{codebox}
def listprint(self):
    printval = self.head
    while printval:
        print(printval.data)
        printval = printval.next.next
\end{codebox}
\begin{choices}
\choice 1 3 5
\choice 2 3 4 5
\choice 1 2 3 4
\choice 5 4 3 2 1
\end{choices}
\textbf{Đáp án:} A
\textbf{Giải thích:} Vòng lặp bắt đầu từ head (giá trị 1), mỗi bước nhảy qua 2 nút bằng printval = printval.next.next. Các giá trị được in ra lần lượt là 1, 3, 5.
\end{question}

\begin{question}{2.30}
% [ID gốc: 1.83]
% [Nguồn: Đề thi FE/RE CSD203 (Câu #71)]
% Quiz: CSD203_SU26_RE - Question #21 (ID: bcbbe06b-66eb-4c86-87ab-09b6a4d95ad0)
What is the key characteristic of a circularly linked list?
\begin{choices}
\choice The last node points to the first node.
\choice It has a fixed size.
\choice The last node points to NULL.
\choice It allows for random access of elements.
\end{choices}
\textbf{Đáp án:} A
\textbf{Giải thích:} Đặc điểm chính của danh sách liên kết vòng là con trỏ next của nút cuối cùng trỏ ngược lại nút đầu tiên thay vì trỏ NULL.
\end{question}

\newpage
\section*{Chương 3: Ngăn xếp (Stacks)}
\addcontentsline{toc}{section}{Chương 3: Ngăn xếp (Stacks)}

\begin{summarybox}{Tóm tắt lý thuyết trọng tâm: Cấu trúc Ngăn xếp (Stack)}
\begin{itemize}[noitemsep,topsep=2pt]
  \item \textbf{Nguyên lý hoạt động:} LIFO (Last In, First Out) -- phần tử vào sau cùng sẽ được lấy ra đầu tiên.
  \item \textbf{Các thao tác cơ bản:} \texttt{push(x)} (thêm vào đỉnh), \texttt{pop()} (lấy ra từ đỉnh), \texttt{peek() / top()} (xem phần tử ở đỉnh), \texttt{isEmpty()} (kiểm tra rỗng), \texttt{isFull()} (khi cài đặt bằng mảng cố định). Mọi thao tác đều có chi phí $O(1)$.
  \item \textbf{Hiện thực trong Python:} Có thể dùng danh sách chuẩn (\texttt{append} / \texttt{pop}), lớp \texttt{collections.deque}, hoặc module \texttt{queue.LifoQueue} (\texttt{put} / \texttt{get}).
  \item \textbf{Ứng dụng kinh điển:} Khử đệ quy, quản lý Call Stack, kiểm tra dấu ngoặc hợp lệ, đảo ngược chuỗi, chuyển đổi và tính toán biểu thức tiền tố (Prefix / Polish) và hậu tố (Postfix / Reverse Polish Notation).
\end{itemize}
\end{summarybox}

\begin{question}{3.1}
% [ID gốc: 1.13]
% [Nguồn: Ngân hàng CSD203 (Câu #51)]
What would be the contents of an empty stack S1 after the following operations are executed (visit from bottom to top)? S1.push(3); S1.push(5); S1.push(9); S1.pop(); S1.push(11);
\begin{choices}
\choice 3 5 9 11
\choice 5 9 11
\choice 3 5 9
\choice 3 5 11
\end{choices}
\textbf{Đáp án:} D
\textbf{Giải thích:} Ngăn xếp hoạt động theo cơ chế LIFO: \texttt{push(3)} $\to$ [3]; \texttt{push(2)} $\to$ [3, 2]; \texttt{push(1)} $\to$ [3, 2, 1]; \texttt{pop()} lấy phần tử 1 ra $\to$ [3, 2]; \texttt{push(5)} đưa 5 vào $\to$ nội dung ngăn xếp là 3, 2, 5 (với 5 ở đỉnh).
\end{question}

\begin{question}{3.2}
% [ID gốc: 1.14]
% [Nguồn: Ngân hàng CSD203 (Câu #59)]
What will be the output of the following code?
\begin{codebox}
from collections import deque
S = deque()
S.append(1)
S.append(2)
S.pop()
S.append(3)
S.pop()
print(S)
\end{codebox}
\begin{choices}
\choice \texttt{deque([1])}
\choice \texttt{deque([1, 2, 3])}
\choice 3
\choice 1
\end{choices}
\textbf{Đáp án:} A
\textbf{Giải thích:} Cấu trúc \texttt{deque} hoạt động với thao tác \texttt{append} và \texttt{pop}: thêm 1 $\to$ [1]; thêm 2 $\to$ [1, 2]; \texttt{pop()} loại bỏ 2 $\to$ [1]; thêm 3 $\to$ [1, 3]; \texttt{pop()} loại bỏ 3 $\to$ [1]. Kết quả in ra là \texttt{deque([1])}.
\end{question}

\begin{question}{3.3}
% [ID gốc: 1.20]
% [Nguồn: Ngân hàng CSD203 (Câu #84)]
What will be the output of the following code?
\begin{codebox}
S = []
S.append(1)
S.append(2)
S.append(3)
S.pop()
print(S)
\end{codebox}
\begin{choices}
\choice \texttt{[1, 2, 3]}
\choice \texttt{[1, 2]}
\choice 1
\choice 3
\end{choices}
\textbf{Đáp án:} B
\textbf{Giải thích:} Danh sách Python hoạt động như ngăn xếp với \texttt{append} (push) và \texttt{pop()} (lấy phần tử ở đỉnh). Sau khi thêm 1, 2, 3 rồi gọi \texttt{S.pop()} thì phần tử 3 bị xóa, danh sách còn lại \texttt{[1, 2]}.
\end{question}

\begin{question}{3.4}
% [ID gốc: 1.34]
% [Nguồn: Ngân hàng CSD203 (Câu #139)]
Process of inserting an element in stack is called
\begin{choices}
\choice Add
\choice Insert
\choice Push
\choice Pop
\end{choices}
\textbf{Đáp án:} C
\textbf{Giải thích:} Trong ngăn xếp (stack), thao tác chèn (thêm) một phần tử vào đỉnh ngăn xếp được gọi chuẩn hóa là thao tác \texttt{Push} (ngược lại thao tác lấy ra là \texttt{Pop}).
\end{question}

\begin{question}{3.5}
% [ID gốc: 1.36]
% [Nguồn: Ngân hàng CSD203 (Câu #143)]
What will be the output of the following code?
\begin{codebox}
S = []
S.append(1)
S.append(2)
S.append(3)
S.append(4)
S.pop()
print(S)
\end{codebox}
\begin{choices}
\choice \texttt{[1, 2, 3]}
\choice \texttt{[2, 3, 4]}
\choice \texttt{[4]}
\choice \texttt{[1]}
\end{choices}
\textbf{Đáp án:} A
\textbf{Giải thích:} Danh sách hoạt động như ngăn xếp: sau 4 lần \texttt{append} danh sách có \texttt{[1, 2, 3, 4]}; lệnh \texttt{S.pop()} lấy ra phần tử ở đỉnh (giá trị 4). Danh sách còn lại \texttt{[1, 2, 3]}.
\end{question}

\begin{question}{3.6}
% [ID gốc: 1.40]
% [Nguồn: Ngân hàng CSD203 (Câu #169)]
What will be the output of the following code?
\begin{codebox}
from queue import LifoQueue
S = LifoQueue(maxsize=5)
S.put('A')
S.put('B')
S.get()
S.put('C')
print(S.qsize())
\end{codebox}
\begin{choices}
\choice 5
\choice 1
\choice 2
\choice 3
\end{choices}
\textbf{Đáp án:} C
\textbf{Giải thích:} LifoQueue là ngăn xếp: đưa vào A, B $\to$ [A, B]; gọi \texttt{get()} lấy B $\to$ [A]; đưa vào C $\to$ [A, C]. Kích thước \texttt{qsize()} bằng 2.
\end{question}

\begin{question}{3.7}
% [ID gốc: 1.42]
% [Nguồn: Ngân hàng CSD203 (Câu #171)]
What will be the output of the following code?
\begin{codebox}
from queue import LifoQueue
S = LifoQueue(maxsize=3)
S.put('A')
S.put('B')
S.put('C')
print(S.get())
\end{codebox}
\begin{choices}
\choice \texttt{['A', 'B', 'C']}
\choice \texttt{['C', 'B', 'A']}
\choice A
\choice C
\end{choices}
\textbf{Đáp án:} D
\textbf{Giải thích:} LifoQueue hoạt động theo cơ chế LIFO: đưa vào 'A', 'B', 'C'. Lệnh \texttt{get()} lấy phần tử ở đỉnh ngăn xếp là 'C'.
\end{question}

\begin{question}{3.8}
% [ID gốc: 1.43]
% [Nguồn: Ngân hàng CSD203 (Câu #173)]
Prefix notation is also known as
\begin{choices}
\choice Reverse Polish Notation
\choice Reverse Notation
\choice Polish Notation
\choice Polish Reverse Notation
\end{choices}
\textbf{Đáp án:} C
\textbf{Giải thích:} Ký pháp tiền tố (Prefix notation) trong đó toán tử đứng trước các toán hạng còn được gọi là ký pháp Ba Lan (Polish Notation).
\end{question}

\begin{question}{3.9}
% [ID gốc: 1.44]
% [Nguồn: Ngân hàng CSD203 (Câu #175)]
Consider the following operation performed on a stack of size 5.
\begin{codebox}
Push(1);
Push(2);
Push(3);
Pop();
Push(4);
Pop();
Pop();
Push(5);
\end{codebox}
After the completion of all operations, the number of elements present on stack are
\begin{choices}
\choice 1
\choice 2
\choice 3
\choice 4
\end{choices}
\textbf{Đáp án:} B
\textbf{Giải thích:} Ngăn xếp kích thước 5: Push(1), Push(2), Push(3) $\to$ [1,2,3]; Pop() $\to$ [1,2]; Push(4) $\to$ [1,2,4]; Pop() $\to$ [1,2]; Pop() $\to$ [1]; Push(5) $\to$ [1,5]. Số phần tử còn lại là 2.
\end{question}

\begin{question}{3.10}
% [ID gốc: 1.45]
% [Nguồn: Ngân hàng CSD203 (Câu #178)]
What is the result of the following Python code for reversing a string using a stack?
\begin{codebox}
def reverse_string(s):
    stack = []
    for i, char in enumerate(s):
        if i % 2 != 0:
            stack.append(char)
    reversed_string = ""
    while stack:
        reversed_string += stack.pop()
    return reversed_string

print(reverse_string("FPT University"))
\end{codebox}
\begin{choices}
\choice yhn
\choice Ptn
\choice ytn
\choice oht
\choice yirvn P
\end{choices}
\textbf{Đáp án:} E
\textbf{Giải thích:} Chuỗi 'FPT University' với các ký tự ở vị trí chỉ số lẻ: 1('P'), 3(' '), 5('n'), 7('v'), 9('r'), 11('i'), 13('y'). Đưa vào stack rồi pop ra tạo thành 'yirvn P'.
\end{question}

\begin{question}{3.11}
% [ID gốc: 1.46]
% [Nguồn: Ngân hàng CSD203 (Câu #192)]
Given an empty stack S and a sequence of the following operations: Push(F); Push(P); Pop(); Push(T); Pop(); Push(U); What does the stack look like? (visit from bottom to top)
\begin{choices}
\choice FU
\choice FPTU
\choice TU
\choice FP
\end{choices}
\textbf{Đáp án:} A
\textbf{Giải thích:} Trạng thái stack: Push(F) $\to$ [F]; Push(P) $\to$ [F,P]; Pop() $\to$ [F]; Push(T) $\to$ [F,T]; Pop() $\to$ [F]; Push(U) $\to$ [F,U]. Từ đáy lên đỉnh là FU.
\end{question}

\begin{question}{3.12}
% [ID gốc: 1.47]
% [Nguồn: Ngân hàng CSD203 (Câu #194)]
Consider the following function:
\begin{codebox}
def fn(expression):
    stack = []
    for char in expression:
        if char.isdigit():
            stack.append(int(char))
        else:
            b = stack.pop()
            a = stack.pop()
            if char == '+':
                stack.append(a + b)
            elif char == '-':
                stack.append(a - b)
            elif char == '*':
                stack.append(a * b)
            elif char == '/':
                stack.append(a // b)
    return stack.pop()
\end{codebox}
What is the output of the function for the input string ``231*+9-''?
\begin{choices}
\choice -4
\choice 4
\choice -12
\choice 12
\end{choices}
\textbf{Đáp án:} A
\textbf{Giải thích:} Tính biểu thức hậu tố '231*+9-': đưa 2, 3, 1 vào stack; gặp '*' tính $3 \times 1 = 3$; gặp '+' tính $2 + 3 = 5$; đưa 9 vào stack; gặp '-' tính $5 - 9 = -4$.
\end{question}

\begin{question}{3.13}
% [ID gốc: 1.55]
% [Nguồn: Ngân hàng CSD203 (Câu #230)]
In a stack implementation, what is the primary difference between the "pop" and "peek" operations?
\begin{choices}
\choice "Pop" removes the top element, while "peek" only returns its value.
\choice "Peek" removes the top element, while "pop" only returns its value.
\choice Both "pop" and "peek" remove the top element but have different return values.
\choice Both "pop" and "peek" return the top element but have different removal behaviors.
\end{choices}
\textbf{Đáp án:} A
\textbf{Giải thích:} Thao tác \texttt{pop} lấy ra và đồng thời xóa phần tử khỏi đỉnh ngăn xếp, trong khi thao tác \texttt{peek} chỉ xem giá trị của phần tử ở đỉnh mà không xóa nó.
\end{question}

\begin{question}{3.14}
% [ID gốc: 1.57]
% [Nguồn: Ngân hàng CSD203 (Câu #232)]
What will be the output of the following code?
\begin{codebox}
stack = []
stack.append(10)
stack.append(20)
stack.pop()
stack.append(30)
print(stack[-1])
\end{codebox}
\begin{choices}
\choice 10
\choice 20
\choice 30
\choice IndexError
\end{choices}
\textbf{Đáp án:} C
\textbf{Giải thích:} Thao tác ngăn xếp: \texttt{append(10)}, \texttt{append(20)}, \texttt{pop()} (xóa 20), \texttt{append(30)}. Phần tử ở đỉnh hiện tại \texttt{stack[-1]} là 30.
\end{question}

\begin{question}{3.15}
% [ID gốc: 1.65]
% [Nguồn: Ngân hàng CSD203 (Câu #263)]
What happens when you attempt to pop an element from an empty Python list used as a stack?
\begin{choices}
\choice An IndexError is raised
\choice The operation fails silently
\choice A None value is returned
\choice The stack resizes to accommodate the operation
\end{choices}
\textbf{Đáp án:} A
\textbf{Giải thích:} Khi thực hiện thao tác \texttt{pop()} trên một danh sách Python rỗng, chương trình sẽ ném ra ngoại lệ \texttt{IndexError: pop from empty list}.
\end{question}

\begin{question}{3.16}
% [ID gốc: 1.68]
% [Nguồn: Đề thi FE/RE CSD203 (Câu #4)]
% Quiz: CSD203_SU26_FE - Question #4 (ID: ccf40fa1-3f08-4723-b6f4-fc024edfc8f1)
In a stack implemented using a linked list, where is the top of the stack typically maintained?
\begin{choices}
\choice At the head of the list
\choice At a middle node
\choice At the tail of the list
\choice At a random position
\end{choices}
\textbf{Đáp án:} A
\textbf{Giải thích:} Trong cấu trúc ngăn xếp (stack) cài đặt bằng danh sách liên kết, đỉnh ngăn xếp thường được duy trì ở phần đầu danh sách (head) để đảm bảo thao tác thêm và xóa phần tử diễn ra với độ phức tạp $O(1)$.
\end{question}

\begin{question}{3.17}
% [ID gốc: 1.69]
% [Nguồn: Đề thi FE/RE CSD203 (Câu #11)]
% Quiz: CSD203_SU26_FE - Question #11 (ID: 37f14749-33e8-4227-a6bf-18aa072babc7)
Which two operations are supported by a stack? Choose 2 correct answers.
\begin{choices}
\choice enqueue
\choice push
\choice dequeue
\choice pop
\end{choices}
\textbf{Đáp án:} BD
\textbf{Giải thích:} Ngăn xếp (stack) hoạt động theo cơ chế LIFO, hỗ trợ hai thao tác chính là đẩy phần tử vào (push) và lấy phần tử ra (pop).
\end{question}

\begin{question}{3.18}
% [ID gốc: 1.76]
% [Nguồn: Đề thi FE/RE CSD203 (Câu #24)]
% Quiz: CSD203_SU26_FE - Question #24 (ID: ff58ad7f-05fc-4661-9e7c-fb6d32b9dfea)
What will be the output of the following code?
\begin{codebox}
S = []
S.append(1)
S.append(2)
S.append(3)
S.pop()
print(S.pop())
\end{codebox}
\begin{choices}
\choice [1, 2, 3]
\choice [1, 2]
\choice 2
\choice 3
\end{choices}
\textbf{Đáp án:} C
\textbf{Giải thích:} Lần gọi S.pop() đầu tiên lấy ra và loại bỏ phần tử 3 ở đỉnh ngăn xếp. Lần gọi S.pop() thứ hai tiếp tục lấy ra và in giá trị 2.
\end{question}

\begin{question}{3.19}
% [ID gốc: 1.88]
% [Nguồn: Đề thi FE/RE CSD203 (Câu #96)]
% Quiz: CSD203_SU26_RE - Question #46 (ID: 417840c5-c0d0-4109-acc2-6cd4094ffe47)
Suppose S is an empty stack. What is the state of the stack S after the following operations sequentially?
\begin{codebox}
S.push(5)
S.push(3)
S.pop()
S.push(7)
S.push(9)
S.pop()
\end{codebox}
\begin{choices}
\choice [5, 3, 7]
\choice [7, 9]
\choice [5, 7]
\choice [7]
\end{choices}
\textbf{Đáp án:} C
\textbf{Giải thích:} S.push(5) → [5]; S.push(3) → [5, 3]; S.pop() loại 3 → [5]; S.push(7) → [5, 7]; S.push(9) → [5, 7, 9]; S.pop() loại 9 → [5, 7]. Trạng thái cuối cùng của ngăn xếp là [5, 7].
\end{question}

\newpage
\section*{Chương 4: Hàng đợi (Queues)}
\addcontentsline{toc}{section}{Chương 4: Hàng đợi (Queues)}

\begin{summarybox}{Tóm tắt lý thuyết trọng tâm: Cấu trúc Hàng đợi (Queue)}
\begin{itemize}[noitemsep,topsep=2pt]
  \item \textbf{Nguyên lý hoạt động:} FIFO (First In, First Out) -- phần tử vào trước sẽ được phục vụ và lấy ra trước.
  \item \textbf{Các thao tác cơ bản:} \texttt{enqueue(x)} (thêm vào đuôi hàng đợi), \texttt{dequeue()} (lấy ra từ đầu hàng đợi), \texttt{front() / peek()} (xem đầu hàng đợi). Chi phí lý tưởng: $O(1)$.
  \item \textbf{Circular Queue (Hàng đợi vòng):} Cài đặt trên mảng một chiều bằng phép toán chia dư (modulo): $\text{rear} = (\text{rear} + 1) \pmod N$, $\text{front} = (\text{front} + 1) \pmod N$ nhằm khắc phục hiện tượng lãng phí bộ nhớ ở đầu mảng.
  \item \textbf{Deque (Double-Ended Queue):} Cho phép thêm và xóa ở cả hai đầu với chi phí $O(1)$ (\texttt{append}, \texttt{appendleft}, \texttt{pop}, \texttt{popleft}).
  \item \textbf{Priority Queue (Hàng đợi ưu tiên):} Mỗi phần tử gắn kèm một độ ưu tiên, phần tử có độ ưu tiên cao nhất luôn được lấy ra trước (thường hiện thực bằng Binary Heap với chi phí chèn/xóa là $O(\log n)$).
\end{itemize}
\end{summarybox}

\begin{question}{4.1}
% [ID gốc: 1.2]
% [Nguồn: Ngân hàng CSD203 (Câu #13)]
What method is used to remove an element from a Queue?
\begin{choices}
\choice \texttt{dequeue()}
\choice \texttt{enqueue()}
\choice \texttt{push()}
\choice \texttt{pop()}
\end{choices}
\textbf{Đáp án:} A
\textbf{Giải thích:} Phương thức \texttt{dequeue()} trong hàng đợi (Queue) lấy ra và xóa phần tử ở đầu hàng đợi theo nguyên lý FIFO.
\end{question}

\begin{question}{4.2}
% [ID gốc: 1.9]
% [Nguồn: Ngân hàng CSD203 (Câu #43)]
The following circular queue can accommodate a maximum six elements with the following data front = 2 rear = 4 queue = \_\_\_; L ; M; N; \_\_\_; \_\_\_ What are the values of front and rear after add O operation takes place? queue = \_\_\_; L ; M; N; O; \_\_\_
\begin{choices}
\choice front = 3, rear = 5
\choice front = 1, rear = 4
\choice front = 2, rear = 5
\choice front = 3, rear = 4
\end{choices}
\textbf{Đáp án:} C
\textbf{Giải thích:} Hàng đợi ưu tiên (Priority Queue) là hàng đợi mà mỗi phần tử có một độ ưu tiên, phần tử có độ ưu tiên cao nhất luôn được lấy ra trước.
\end{question}

\begin{question}{4.3}
% [ID gốc: 1.19]
% [Nguồn: Ngân hàng CSD203 (Câu #82)]
What will be the output of the following code?
\begin{codebox}
S=[]
for i in range(1,5):
    if (i%2 != 1):
        S.append(i)
while (len(S) > 0):
    print(S.pop(0))
\end{codebox}
\begin{choices}
\choice 1 3
\choice 2 4
\choice 3 1
\choice 4 2
\end{choices}
\textbf{Đáp án:} B
\textbf{Giải thích:} Vòng lặp \texttt{for} duyệt qua $i \in \{1, 2, 3, 4\}$, các số chẵn (chia 2 dư 0) là 2 và 4 được đưa vào \texttt{S}. Vòng lặp \texttt{while} lấy lần lượt phần tử đầu \texttt{S.pop(0)} theo cơ chế FIFO và in ra: 2 rồi đến 4.
\end{question}

\begin{question}{4.4}
% [ID gốc: 1.29]
% [Nguồn: Ngân hàng CSD203 (Câu #108)]
Given an empty queue Q and a sequence of the following operations: Enqueue(5) Enqueue(15) Enqueue(25) Dequeue(); Dequeue() The value of the front element is:
\begin{choices}
\choice 0
\choice 5
\choice 15
\choice 25
\end{choices}
\textbf{Đáp án:} D
\textbf{Giải thích:} Hàng đợi hoạt động theo nguyên lý FIFO: \texttt{Enqueue(5)} $\to$ [5]; \texttt{Enqueue(7)} $\to$ [5, 7]; \texttt{Dequeue()} lấy 5 ra $\to$ [7]; \texttt{Enqueue(12)} $\to$ [7, 12]; \texttt{Dequeue()} lấy 7 ra $\to$ phần tử còn lại trong hàng đợi là 12.
\end{question}

\begin{question}{4.5}
% [ID gốc: 1.48]
% [Nguồn: Ngân hàng CSD203 (Câu #200)]
What method is used to remove an element from a Queue?
\begin{choices}
\choice dequeue()
\choice enqueue()
\choice push()
\choice pop()
\end{choices}
\textbf{Đáp án:} A
\textbf{Giải thích:} Phương thức chuẩn dùng để lấy ra một phần tử khỏi hàng đợi là \texttt{dequeue()}.
\end{question}

\begin{question}{4.6}
% [ID gốc: 1.56]
% [Nguồn: Ngân hàng CSD203 (Câu #231)]
What is the output of the following function for the input [1, 2, 3, 4, 5, 6]?
\begin{codebox}
def fn(queue):
    q = deque(queue)
    size = len(q)
    stack = []

    for _ in range(size // 2):
        stack.append(q.popleft())

    while stack:
        q.append(stack.pop())

    for _ in range(size // 2):
        q.append(q.popleft())
        q.append(q.popleft())

    return list(q)
\end{codebox}
\begin{choices}
\choice \texttt{[[1, 5], [2, 6], [3, 4]]}
\choice \texttt{[4, 5, 6, 3, 2, 1]}
\choice \texttt{[1, 3, 5, 2, 4, 6]}
\choice The code will raise an error.
\end{choices}
\textbf{Đáp án:} B
\textbf{Giải thích:} Hàm đưa nửa đầu danh sách \texttt{[1, 2, 3]} vào stack rồi pop ngược lại vào cuối queue thành \texttt{[4, 5, 6, 3, 2, 1]}, sau đó các thao tác xoay queue tiếp theo giữ nguyên trật tự này, kết quả là \texttt{[4, 5, 6, 3, 2, 1]}.
\end{question}

\begin{question}{4.7}
% [ID gốc: 1.64]
% [Nguồn: Ngân hàng CSD203 (Câu #261)]
Consider the following code. What is the output?
\begin{codebox}
from queue import Queue
queue = Queue()
queue.put(5)
queue.put(15)

while not queue.empty():
    queue.get()

print(queue.empty())
\end{codebox}
\begin{choices}
\choice true
\choice false
\choice 0
\choice An exception is thrown
\end{choices}
\textbf{Đáp án:} A
\textbf{Giải thích:} Hàng đợi \texttt{Queue} đưa vào 5 và 15, vòng lặp \texttt{while not queue.empty(): queue.get()} lấy hết tất cả các phần tử ra ngoài, do đó lệnh \texttt{print(queue.empty())} in ra \texttt{True}.
\end{question}

\begin{question}{4.8}
% [ID gốc: 1.66]
% [Nguồn: Ngân hàng CSD203 (Câu #264)]
What does the following function do when given an array of integers as input?
\begin{codebox}
def fn(arr, k):
    queue = deque(arr)
    for _ in range(k):
        queue.append(queue.popleft())
    return list(queue)
\end{codebox}
If the input is [10, 20, 30, 40, 50] and k = 2, what will be the output?
\begin{choices}
\choice [40, 50, 10, 20, 30]
\choice [30, 40, 50, 10, 20]
\choice [20, 30, 40, 50, 10]
\choice [50, 10, 20, 30, 40]
\end{choices}
\textbf{Đáp án:} B
\textbf{Giải thích:} Hàm xoay vòng $k = 2$ lần: lấy phần tử đầu đưa về cuối mảng 2 lần, biến đổi \texttt{[10, 20, 30, 40, 50]} thành \texttt{[30, 40, 50, 10, 20]}.
\end{question}

\begin{question}{4.9}
% [ID gốc: 1.67]
% [Nguồn: Ngân hàng CSD203 (Câu #266)]
What will be the output of the following code?
\begin{codebox}
from collections import deque
Q = deque()
Q.append('A')
Q.append('B')
Q.append('C')
Q.pop()
print(Q.popleft())
\end{codebox}
\begin{choices}
\choice A
\choice B
\choice [A, B, C]
\choice Error
\end{choices}
\textbf{Đáp án:} A
\textbf{Giải thích:} Cấu trúc deque: thêm 'A', 'B', 'C'; lệnh \texttt{Q.pop()} xóa phần tử cuối 'C'; lệnh \texttt{Q.popleft()} lấy ra phần tử ở đầu hàng đợi là 'A'.
\end{question}

\begin{question}{4.10}
% [ID gốc: 1.70]
% [Nguồn: Đề thi FE/RE CSD203 (Câu #13)]
% Quiz: CSD203_SU26_FE - Question #13 (ID: 11125bd5-ce6d-49e1-95a3-f946e875105e)
Which method is not a part of the Queue module in Python?
\begin{choices}
\choice put()
\choice get()
\choice push()
\choice empty()
\end{choices}
\textbf{Đáp án:} C
\textbf{Giải thích:} Module queue trong Python sử dụng phương thức put() để thêm phần tử và get() để lấy phần tử, trong khi push() không phải là phương thức chuẩn của module này.
\end{question}

\begin{question}{4.11}
% [ID gốc: 1.78]
% [Nguồn: Đề thi FE/RE CSD203 (Câu #39)]
% Quiz: CSD203_SU26_FE - Question #39 (ID: b04b9800-06d6-4f91-8732-c1e68d3f0584)
If popleft() is used on an empty deque, what happens?
\begin{choices}
\choice Raises error
\choice Appends None
\choice Returns null
\choice Removes nothing
\end{choices}
\textbf{Đáp án:} A
\textbf{Giải thích:} Khi gọi phương thức xóa phần tử bên trái trên một hàng đợi hai đầu (collections.deque) rỗng, Python phát sinh ngoại lệ IndexError.
\end{question}

\begin{question}{4.12}
% [ID gốc: 1.79]
% [Nguồn: Đề thi FE/RE CSD203 (Câu #43)]
% Quiz: CSD203_SU26_FE - Question #43 (ID: 83700be3-72b5-4760-861b-2bb2ed8be01a)
What will be the output of the following code?
\begin{codebox}
from queue import Queue
q = Queue(maxsize=3)
q.put(10)
q.put(20)
q.put(30)
print(q.full())
q.get()
print(q.full())
\end{codebox}
\begin{choices}
\choice False True
\choice True True
\choice True False
\choice False False
\end{choices}
\textbf{Đáp án:} C
\textbf{Giải thích:} Đưa 3 phần tử vào hàng đợi có maxsize=3 làm q.full() trả về True. Sau khi gọi q.get() lấy ra 1 phần tử, hàng đợi còn trống chỗ nên q.full() chuyển thành False.
\end{question}

\begin{question}{4.13}
% [ID gốc: 1.84]
% [Nguồn: Đề thi FE/RE CSD203 (Câu #85)]
% Quiz: CSD203_SU26_RE - Question #35 (ID: 22a2c6e1-0e80-4a70-a546-b7c4976f645c)
What does the following function do when given an array of integers as input?
\begin{codebox}
from collections import deque

def fn(arr, k):
    queue = deque(arr)
    for _ in range(k):
        queue.append(queue.popleft())
    return list(queue)
\end{codebox}
If the input is [10, 20, 30, 40, 50] and $k = 2$, what will be the output?
\begin{choices}
\choice [40, 50, 10, 20, 30]
\choice [30, 40, 50, 10, 20]
\choice [20, 30, 40, 50, 10]
\choice [50, 10, 20, 30, 40]
\end{choices}
\textbf{Đáp án:} B
\textbf{Giải thích:} Hàm thực hiện xoay trái mảng k lần (chuyển phần tử đầu xuống cuối). Sau lần 1: [20, 30, 40, 50, 10]. Sau lần 2: [30, 40, 50, 10, 20].
\end{question}

\begin{question}{4.14}
% [ID gốc: 1.86]
% [Nguồn: Đề thi FE/RE CSD203 (Câu #89)]
% Quiz: CSD203_SU26_RE - Question #39 (ID: 837ca48a-e372-40b5-a530-0c6b21bf2fca)
Given an empty queue Q and a sequence of the following operations:
\begin{codebox}
Enqueue(5)
Enqueue(15)
Enqueue(25)
Dequeue()
Dequeue()
\end{codebox}
The value of the front element is \_\_\_\_\_.
\begin{choices}
\choice 0
\choice 5
\choice 15
\choice 25
\end{choices}
\textbf{Đáp án:} D
\textbf{Giải thích:} Sau ba lần enqueue, hàng đợi là [5, 15, 25] với đầu hàng đợi là 5. Dequeue() lấy 5 ra còn [15, 25]; Dequeue() lấy 15 ra còn [25]. Phần tử ở đầu (front) hàng đợi lúc này là 25.
\end{question}

\begin{question}{4.15}
% [ID gốc: 1.87]
% [Nguồn: Đề thi FE/RE CSD203 (Câu #92)]
% Quiz: CSD203_SU26_RE - Question #42 (ID: 3fc97385-79e8-4cf4-a31c-5e4265b961bc)
Which of the following is an advantage of a linked list-based queue?
\begin{choices}
\choice It can grow and shrink dynamically.
\choice It is more efficient for enqueuing elements.
\choice It is more memory-efficient than an array-based queue.
\choice It is more efficient for dequeuing elements.
\end{choices}
\textbf{Đáp án:} A
\textbf{Giải thích:} Ưu điểm lớn nhất của hàng đợi cài đặt bằng danh sách liên kết so với mảng là khả năng mở rộng và thu nhỏ kích thước linh hoạt mà không cần khai báo kích thước cố định trước.
\end{question}

\begin{question}{4.16}
% [ID gốc: 1.90]
% [Nguồn: Đề thi FE/RE CSD203 (Câu #99)]
% Quiz: CSD203_SU26_RE - Question #49 (ID: d8e9a888-56d4-4980-9daa-0e8ac6c1a0cc)
What values are returned during the following sequence of queue operations, if executed on an initially empty queue?
\begin{codebox}
enqueue(4)
enqueue(6)
dequeue()
enqueue(7)
dequeue()
enqueue(5)
dequeue()
\end{codebox}
\begin{choices}
\choice 4, 6, 7
\choice 4, 6, 7, 5
\choice 7, 6, 4
\choice 6, 7, 4
\end{choices}
\textbf{Đáp án:} A
\textbf{Giải thích:} enqueue(4), enqueue(6) → [4, 6]; dequeue() trả về 4 → [6]; enqueue(7) → [6, 7]; dequeue() trả về 6 → [7]; enqueue(5) → [7, 5]; dequeue() trả về 7 → [5]. Các giá trị trả về theo thứ tự là 4, 6, 7.
\end{question}

\newpage
\section*{Chương 5: Bảng băm (Hash Tables)}
\addcontentsline{toc}{section}{Chương 5: Bảng băm (Hash Tables)}

\begin{summarybox}{Tóm tắt lý thuyết trọng tâm: Bảng băm (Hash Table) \& Xử lý xung đột}
\begin{itemize}[noitemsep,topsep=2pt]
  \item \textbf{Khái niệm \& Hàm băm:} Bảng băm ánh xạ khóa $k$ thành chỉ mục trong mảng thông qua hàm băm $h(k)$ (phổ biến là hàm chia $h(k) = k \pmod m$). Thời gian tìm kiếm, chèn, xóa trung bình là $O(1)$.
  \item \textbf{Hệ số tải (Load Factor):} $\alpha = \frac{n}{m}$ với $n$ là số phần tử và $m$ là kích thước bảng. Khi $\alpha$ vượt ngưỡng (thường là $0.7 - 0.75$), cần mở rộng bảng (Rehash).
  \item \textbf{Phương pháp nối kết (Separate Chaining):} Mỗi ô của bảng là một danh sách liên kết. Hệ số tải $\alpha$ có thể lớn hơn 1 và số phần tử không bị giới hạn bởi kích thước mảng.
  \item \textbf{Địa chỉ mở (Open Addressing):} Tất cả phần tử nằm trực tiếp trong bảng ($\alpha \le 1$). Khi xung đột, tìm ô trống kế tiếp:
    \begin{itemize}[noitemsep]
      \item \textit{Thăm dò tuyến tính (Linear Probing):} $h(k, i) = (h(k) + i) \pmod m$ (dễ bị gom cụm nguyên phát -- Primary Clustering).
      \item \textit{Thăm dò bậc hai (Quadratic Probing):} $h(k, i) = (h(k) + c_1 i + c_2 i^2) \pmod m$ (giảm gom cụm nguyên phát nhưng có thể bị Secondary Clustering).
      \item \textit{Băm kép (Double Hashing):} $h(k, i) = (h_1(k) + i \cdot h_2(k)) \pmod m$ (tạo chuỗi thăm dò tốt nhất, $h_2(k)$ phải nguyên tố cùng nhau với $m$).
    \end{itemize}
\end{itemize}
\end{summarybox}

\begin{question}{5.1}
% [ID gốc: 3.1]
% [Nguồn: Ngân hàng CSD203 (Câu #1)]
A chained hash table has an array size of 255. What is the maximum number of elements that can be placed in the table?
\begin{choices}
\choice 254
\choice There is no limit
\choice 255
\choice 256
\end{choices}
\textbf{Đáp án:} B
\textbf{Giải thích:} Bảng băm nối chuỗi (Separate Chaining) lưu các phần tử trùng vị trí băm vào danh sách liên kết ngoài, nên về mặt lý thuyết không có giới hạn dung lượng phần tử.
\end{question}

\begin{question}{5.2}
% [ID gốc: 3.2]
% [Nguồn: Ngân hàng CSD203 (Câu #11)]
What is a hash table?
\begin{choices}
\choice A structure used for storage
\choice A structure used to implement stack and queue
\choice A structure that maps values to keys
\choice A structure that maps keys to values
\end{choices}
\textbf{Đáp án:} D
\textbf{Giải thích:} Bảng băm là cấu trúc dữ liệu lưu trữ dữ liệu dưới dạng các cặp khóa--giá trị (key-value), ánh xạ khóa thành chỉ số mảng thông qua hàm băm.
\end{question}

\begin{question}{5.3}
% [ID gốc: 3.3]
% [Nguồn: Ngân hàng CSD203 (Câu #18)]
In which of the following hash functions, do consecutive keys map to consecutive hash values?
\begin{choices}
\choice Folding method
\choice Multiplication method
\choice Mid-square method
\choice Division method
\end{choices}
\textbf{Đáp án:} D
\textbf{Giải thích:} Hàm băm chia lấy dư $h(k) = k \bmod m$ ánh xạ các khóa liên tiếp $k, k+1, k+2$ vào các ô liên tiếp nhau trong bảng băm.
\end{question}

\begin{question}{5.4}
% [ID gốc: 3.4]
% [Nguồn: Ngân hàng CSD203 (Câu #34)]
In a hash table, an element with key $k$ is stored at index
\begin{choices}
\choice $\text{Log } k$
\choice $k * k$
\choice $H(k)$
\choice $k$
\end{choices}
\textbf{Đáp án:} C
\textbf{Giải thích:} Duyệt BFS sử dụng hàng đợi (Queue), thăm đỉnh gốc rồi lần lượt thăm tất cả các đỉnh lân cận ở mức hiện tại trước khi xuống mức tiếp theo.
\end{question}

\begin{question}{5.5}
% [ID gốc: 3.5]
% [Nguồn: Ngân hàng CSD203 (Câu #39)]
Given a hash table T with 20 slots that stores 1000 elements, the load factor $\alpha$ for T is
\begin{choices}
\choice 20
\choice 1000
\choice 0.02
\choice 50
\end{choices}
\textbf{Đáp án:} D
\textbf{Giải thích:} Trong danh sách liên kết đơn, thao tác chèn nút vào đầu danh sách (insert at head) chỉ cần cập nhật con trỏ head, thực hiện trong thời gian $O(1)$.
\end{question}

\begin{question}{5.6}
% [ID gốc: 3.6]
% [Nguồn: Ngân hàng CSD203 (Câu #41)]
For the given hash table with size 100 and folding method is used, in what location will the key 2844267 be hashed using probing?
\begin{choices}
\choice 7
\choice 105
\choice 5
\choice 67
\end{choices}
\textbf{Đáp án:} C
\textbf{Giải thích:} Cây tìm kiếm nhị phân có tính chất: khóa của nút con trái nhỏ hơn khóa của nút cha, và khóa của nút con phải lớn hơn khóa của nút cha.
\end{question}

\begin{question}{5.7}
% [ID gốc: 3.7]
% [Nguồn: Ngân hàng CSD203 (Câu #52)]
An advantage of chained hash table (external hashing) over the open addressing scheme is
\begin{choices}
\choice Worst case complexity of search operations is less
\choice Space used is less
\choice Deletion is easier
\choice None of these
\end{choices}
\textbf{Đáp án:} C
\textbf{Giải thích:} Dò bậc hai (Quadratic Probing) sử dụng hàm bước nhảy $i^2$ để giải quyết va chạm trong bảng băm địa chỉ mở, giúp giảm gom cụm sơ cấp.
\end{question}

\begin{question}{5.8}
% [ID gốc: 3.8]
% [Nguồn: Ngân hàng CSD203 (Câu #64)]
In a hash table, what is called if two different keys are hashed to the same location?
\begin{choices}
\choice Collision
\choice Addressing
\choice Hashing
\choice Probing
\end{choices}
\textbf{Đáp án:} A
\textbf{Giải thích:} Hàng đợi hai đầu (Deque) cho phép thêm và xóa phần tử ở cả hai đầu (front và rear) với độ phức tạp $O(1)$.
\end{question}

\begin{question}{5.9}
% [ID gốc: 3.9]
% [Nguồn: Ngân hàng CSD203 (Câu #69)]
A hash function h defined h(key)=key mod 5, linear probing is used to insert sequentially the keys 24,16, 36 into a hash table indexed from 0 to 4. What will the location of the key 36 be?
\begin{choices}
\choice 0
\choice 1
\choice 2
\choice 3
\end{choices}
\textbf{Đáp án:} C
\textbf{Giải thích:} Thuật toán sắp xếp đếm (Counting Sort) hoạt động trong thời gian tuyến tính $O(n + k)$ khi giá trị các phần tử nằm trong khoảng cố định từ 0 đến $k$.
\end{question}

\begin{question}{5.10}
% [ID gốc: 3.10]
% [Nguồn: Ngân hàng CSD203 (Câu #76)]
A map/dictionary is created and the following key-value pairs are added in the following order: \{"a" : 2\}, \{"b": 3\}, \{"a", 4\}, \{"c", 5\}. What is the value associated with the key "a"?
\begin{choices}
\choice 2
\choice 3
\choice 4
\choice There's an error
\end{choices}
\textbf{Đáp án:} C
\textbf{Giải thích:} Dictionary trong Python lưu trữ các cặp key-value; phương thức truy xuất khóa trả về giá trị tương ứng đã được gán.
\end{question}

\begin{question}{5.11}
% [ID gốc: 3.11]
% [Nguồn: Ngân hàng CSD203 (Câu #100)]
What is simple uniform hashing?
\begin{choices}
\choice Every element has equal probability of hashing into any of the slots
\choice A weighted probabilistic method is used to hash elements into the slots
\choice Elements has Random probability of hashing into array slots
\choice Elements are hashed based on priority
\end{choices}
\textbf{Đáp án:} A
\textbf{Giải thích:} Cây nhị phân hoàn chỉnh (Complete Binary Tree) có tất cả các mức đều đầy, ngoại trừ mức cuối cùng được điền từ trái sang phải.
\end{question}

\begin{question}{5.12}
% [ID gốc: 3.12]
% [Nguồn: Ngân hàng CSD203 (Câu #101)]
In hashing, what is direct addressing ?
\begin{choices}
\choice Distinct array position for every possible key
\choice Fewer array positions than keys
\choice Fewer keys than array positions
\choice Same array position for all keys
\end{choices}
\textbf{Đáp án:} A
\textbf{Giải thích:} Thời gian tìm kiếm trung bình trong cây AVL là $O(\log n)$ nhờ tính chất luôn cân bằng nghiêm ngặt.
\end{question}

\begin{question}{5.13}
% [ID gốc: 3.13]
% [Nguồn: Ngân hàng CSD203 (Câu #102)]
What is the hash function used in the division method?
\begin{choices}
\choice $h(k) = k/m$
\choice $h(k) = k \bmod m$
\choice $h(k) = m/k$
\choice $h(k) = m \bmod k$
\end{choices}
\textbf{Đáp án:} B
\textbf{Giải thích:} Hàm băm tốt cần phân bố đều các khóa trên toàn bộ bảng băm (uniform distribution) để giảm thiểu tối đa xác suất xảy ra va chạm.
\end{question}

\begin{question}{5.14}
% [ID gốc: 3.14]
% [Nguồn: Ngân hàng CSD203 (Câu #103)]
Hash function is used
\begin{choices}
\choice to converts each key k into an index in the hash table.
\choice to converts each index in the hash table into a key.
\choice to find a key in the hash table.
\choice to calculate a key k.
\end{choices}
\textbf{Đáp án:} A
\textbf{Giải thích:} Cấu trúc dữ liệu hàng đợi (Queue) mô hình hóa quy trình hàng đợi trong thực tế như xử lý lệnh in ấn (Spooling) và lập lịch tiến trình CPU.
\end{question}

\begin{question}{5.15}
% [ID gốc: 3.15]
% [Nguồn: Ngân hàng CSD203 (Câu #104)]
Given a hash table with size 100 and Division hashing method is used, in what location will the key 124 be hashed?
\begin{choices}
\choice 24
\choice 1
\choice 4
\choice 12
\end{choices}
\textbf{Đáp án:} A
\textbf{Giải thích:} Thuật toán Bellman-Ford có thể tìm đường đi ngắn nhất trên đồ thị có cạnh trọng số âm và phát hiện chu trình âm, với độ phức tạp $O(VE)$.
\end{question}

\begin{question}{5.16}
% [ID gốc: 3.16]
% [Nguồn: Ngân hàng CSD203 (Câu #132)]
A hash function h defined $h(key) = key \pmod{10}$, quadratic probing is used to insert sequentially the keys 2, 4, 7, 17 into a hash table indexed from 0 to 9. What will the location of the key 17 be?
\begin{choices}
\choice 0
\choice 1
\choice 7
\choice 8
\end{choices}
\textbf{Đáp án:} D
\textbf{Giải thích:} Trong cây AVL, phép quay RL (Right-Left rotation) được áp dụng khi cây con phải bị mất cân bằng do chèn phần tử vào nhánh trái của nó.
\end{question}

\begin{question}{5.17}
% [ID gốc: 3.17]
% [Nguồn: Ngân hàng CSD203 (Câu #134)]
Which of the following method is not a collision resolution strategy for open addressing in hash table?
\begin{choices}
\choice Linear probing.
\choice Quadratic probing.
\choice Double hashing.
\choice Chaining.
\end{choices}
\textbf{Đáp án:} D
\textbf{Giải thích:} Thuật toán Floyd-Warshall tìm đường đi ngắn nhất giữa mọi cặp đỉnh trên đồ thị trong thời gian $O(V^3)$.
\end{question}

\begin{question}{5.18}
% [ID gốc: 3.18]
% [Nguồn: Ngân hàng CSD203 (Câu #135)]
What is the value of the Shift Folding Hash Function if K = 723-203-541-213-24 and table size = 1000?
\begin{choices}
\choice 704
\choice 24
\choice 723
\choice 1704
\end{choices}
\textbf{Đáp án:} A
\textbf{Giải thích:} Thao tác tìm kiếm nhị phân yêu cầu dữ liệu đầu vào phải được sắp xếp trước, không thể áp dụng trực tiếp cho mảng chưa sắp xếp.
\end{question}

\begin{question}{5.19}
% [ID gốc: 3.19]
% [Nguồn: Ngân hàng CSD203 (Câu #137)]
A hash function h defined $h(key) = key \pmod{5}$, with Quadratic probing, is used to insert the keys 1, 10, 11, 12 into a table indexed from 0 to 4. What will be the location of key 12?
\begin{choices}
\choice 4
\choice 3
\choice 0
\choice 1
\end{choices}
\textbf{Đáp án:} B
\textbf{Giải thích:} Dò tuyến tính trong bảng băm có công thức hàm thăm dò là $h(k, i) = (h'(k) + i) \bmod m$.
\end{question}

\begin{question}{5.20}
% [ID gốc: 3.20]
% [Nguồn: Ngân hàng CSD203 (Câu #160)]
What will be the output of the following code?
\begin{codebox}
student = {
  1: "Tuan",
  2: "Hoa",
  3: "Nam"
}
x = student.keys()
print(x)
\end{codebox}
\begin{choices}
\choice \texttt{dict\_keys([1, 2, 3])}
\choice \texttt{dict\_values(['Tuan', 'Hoa', 'Nam'])}
\choice \texttt{[1: "Tuan", 2: "Hoa", 3: "Nam"]}
\choice 3
\end{choices}
\textbf{Đáp án:} A
\textbf{Giải thích:} Dictionary \texttt{student} chứa các khóa 1, 2, 3; phương thức \texttt{student.keys()} trả về đối tượng \texttt{dict\_keys([1, 2, 3])}.
\end{question}

\begin{question}{5.21}
% [ID gốc: 3.21]
% [Nguồn: Ngân hàng CSD203 (Câu #163)]
A hash function h defined $h(key) = key \bmod 10$, linear probing is used to insert sequentially the keys 2, 4, 7, 17 into a hash table indexed from 0 to 9. What will the location of the key 8 be?
\begin{choices}
\choice 0
\choice 1
\choice 7
\choice 9
\end{choices}
\textbf{Đáp án:} D
\textbf{Giải thích:} Thuật toán tìm kiếm nhị phân chia đôi không gian tìm kiếm ở mỗi bước, số phần tử cần xem xét giảm theo cấp số nhân cơ số 2.
\end{question}

\begin{question}{5.22}
% [ID gốc: 3.22]
% [Nguồn: Ngân hàng CSD203 (Câu #166)]
Given a 11-entry hash table and a hash function, $h(i) = (3i + 5) \bmod 11$, to hash the keys 2, 12, 15, and 1, assuming no collision is handled. Where is 1 hashed?
\begin{choices}
\choice 1
\choice 2
\choice 8
\choice collision
\end{choices}
\textbf{Đáp án:} D
\textbf{Giải thích:} Trong bảng băm địa chỉ mở, khi xóa một phần tử ta phải đánh dấu ô đó là 'Đã xóa' (Deleted marker) thay vì để trống để không ngắt chuỗi dò.
\end{question}

\begin{question}{5.23}
% [ID gốc: 3.23]
% [Nguồn: Ngân hàng CSD203 (Câu #167)]
A chained hash table has an array size of N (separate chaining for solving collision). What is the maximum number of elements that can be placed in the table?
\begin{choices}
\choice $N$
\choice $N + 1$
\choice $N - 1$
\choice There is no limit
\end{choices}
\textbf{Đáp án:} D
\textbf{Giải thích:} Thuật toán sắp xếp nhanh (QuickSort) thường không ổn định (unstable) do các phép hoán đổi cách xa nhau qua pivot.
\end{question}

\begin{question}{5.24}
% [ID gốc: 3.24]
% [Nguồn: Ngân hàng CSD203 (Câu #168)]
Given a hash table T with 10 slots that stores 5 elements, the load factor $\alpha$ for T is
\begin{choices}
\choice 10
\choice 5
\choice 50
\choice 0.5
\end{choices}
\textbf{Đáp án:} D
\textbf{Giải thích:} Chu trình Hamilton là chu trình đi qua mỗi đỉnh của đồ thị đúng một lần và quay trở về đỉnh xuất phát.
\end{question}

\begin{question}{5.25}
% [ID gốc: 3.25]
% [Nguồn: Ngân hàng CSD203 (Câu #199)]
Consider a hash function that distributes keys uniformly. The hash table size is 20. After hashing of how many keys will the probability that any new key hashed collides with an existing one exceed 0.5
\begin{choices}
\choice 1
\choice 5
\choice 10
\choice 15
\end{choices}
\textbf{Đáp án:} D
\textbf{Giải thích:} Theo nguyên lý Dirichlet / xác suất xung đột bảng băm (hoặc bài toán nghịch lý ngày sinh), với bảng kích thước $m = 20$, sau khi băm một số lượng khóa nhất định thì xác suất va chạm của khóa mới vượt 0.5; theo các phương án của đề thi thì đáp án là 15.
\end{question}

\begin{question}{5.26}
% [ID gốc: 3.26]
% [Nguồn: Ngân hàng CSD203 (Câu #204)]
Given a hash table of size $S$ and a division hash function $h()$, an element with key $k$ is stored at index
\begin{choices}
\choice $k$
\choice $k/S$
\choice $k\,\%\,S$
\choice $h(k/S)$
\end{choices}
\textbf{Đáp án:} C
\textbf{Giải thích:} Trong bảng băm kích thước $S$ sử dụng hàm băm chia lấy dư (Division method), một phần tử có khóa $k$ được lưu trữ tại chỉ số $k \bmod S$ (tức \texttt{k \% S}).
\end{question}

\begin{question}{5.27}
% [ID gốc: 3.27]
% [Nguồn: Ngân hàng CSD203 (Câu #221)]
Question: 17 What is the main goal of collision resolution in hash tables?
\begin{choices}
\choice To decrease the size of the hash table
\choice To handle cases where multiple keys map to the same index
\choice To increase the complexity of the hash function
\choice To sort the keys in the hash table
\end{choices}
\textbf{Đáp án:} B
\textbf{Giải thích:} Mục tiêu cốt lõi của các phương pháp xử lý đụng độ (Collision resolution) trong bảng băm là giải quyết tình huống khi có nhiều khóa khác nhau được hàm băm ánh xạ về cùng một chỉ số.
\end{question}

\begin{question}{5.28}
% [ID gốc: 3.28]
% [Nguồn: Ngân hàng CSD203 (Câu #222)]
What data structure is used to implement a hash table?
\begin{choices}
\choice Arrays
\choice Linked Lists
\choice Trees
\choice Graphs
\end{choices}
\textbf{Đáp án:} A
\textbf{Giải thích:} Cấu trúc dữ liệu nền tảng nhất được sử dụng để hiện thực bảng băm (Hash Table) là mảng (Arrays).
\end{question}

\begin{question}{5.29}
% [ID gốc: 3.29]
% [Nguồn: Ngân hàng CSD203 (Câu #223)]
In the worst-case scenario, what is the time complexity of a search operation in a hash table with separate chaining?
\begin{choices}
\choice $O(n)$
\choice $O(\log n)$
\choice $O(1)$
\choice $O(n^2)$
\end{choices}
\textbf{Đáp án:} A
\textbf{Giải thích:} Trong trường hợp xấu nhất của bảng băm nối chuỗi (Separate Chaining), tất cả $n$ phần tử bị băm vào cùng một ô và nằm trong cùng một danh sách liên kết, thao tác tìm kiếm suy biến thành $O(n)$.
\end{question}

\begin{question}{5.30}
% [ID gốc: 3.30]
% [Nguồn: Ngân hàng CSD203 (Câu #224)]
What effect does a higher load factor have on the efficiency of a hash table?
\begin{choices}
\choice Higher load factors decrease efficiency
\choice Higher load factors have no effect on efficiency
\choice Higher load factors increase efficiency
\choice Efficiency is unrelated to load factor
\end{choices}
\textbf{Đáp án:} A
\textbf{Giải thích:} Hệ số tải (Load factor) càng cao đồng nghĩa bảng băm càng đầy, làm tăng số lượng va chạm và các chuỗi thăm dò/danh sách liên kết dài hơn, dẫn đến giảm hiệu năng thao tác.
\end{question}

\begin{question}{5.31}
% [ID gốc: 3.31]
% [Nguồn: Ngân hàng CSD203 (Câu #226)]
What is the main disadvantage of using chaining to handle collisions?
\begin{choices}
\choice Increased memory usage
\choice Increased complexity of hash function
\choice Decreased average search time
\choice Difficulty in implementing
\end{choices}
\textbf{Đáp án:} A
\textbf{Giải thích:} Nhược điểm chính của phương pháp nối chuỗi (Chaining) là tiêu tốn thêm bộ nhớ cho các con trỏ liên kết trong danh sách và phân bổ bộ nhớ không liên tục.
\end{question}

\begin{question}{5.32}
% [ID gốc: 3.32]
% [Nguồn: Ngân hàng CSD203 (Câu #235)]
What is the load factor in the context of hash tables?
\begin{choices}
\choice The ratio of the number of elements to the total available slots.
\choice The average number of probes required for each search.
\choice The number of collisions that have occurred.
\choice The number of slots available for new elements.
\end{choices}
\textbf{Đáp án:} A
\textbf{Giải thích:} Hệ số tải (Load factor $\alpha$) trong bảng băm là tỷ số giữa số lượng phần tử hiện có trong bảng và tổng số ô chứa có sẵn ($\alpha = n/m$).
\end{question}

\begin{question}{5.33}
% [ID gốc: 3.33]
% [Nguồn: Ngân hàng CSD203 (Câu #252)]
Which collision resolution technique involves storing multiple elements at the same index in a hash table?
\begin{choices}
\choice Linear probing
\choice Chaining
\choice Quadratic probing
\choice Double hashing
\end{choices}
\textbf{Đáp án:} B
\textbf{Giải thích:} Kỹ thuật nối chuỗi (Chaining) giải quyết đụng độ trong bảng băm bằng cách lưu trữ tất cả các phần tử có cùng chỉ số băm vào một danh sách liên kết tại ô đó.
\end{question}

\begin{question}{5.34}
% [ID gốc: 3.34]
% [Nguồn: Ngân hàng CSD203 (Câu #253)]
What is the purpose of a hash function in hash tables?
\begin{choices}
\choice To generate a unique identifier for each element.
\choice To determine the size of the hash table.
\choice To calculate the average time complexity.
\choice To handle collisions.
\choice To create collisions.
\choice To avoid collisions.
\end{choices}
\textbf{Đáp án:} A
\textbf{Giải thích:} Mục đích của hàm băm (Hash function) là ánh xạ khóa dữ liệu đầu vào thành một chỉ số mảng (unique index / identifier) để định vị phần tử trong bảng.
\end{question}

\begin{question}{5.35}
% [ID gốc: 3.35]
% [Nguồn: Ngân hàng CSD203 (Câu #254)]
When rehashing occurs in a hash table due to an increase in load factor, what is the impact on the table's performance?
\begin{choices}
\choice Improved efficiency
\choice Increased collision likelihood
\choice Reduced memory usage
\choice Constant time complexity
\end{choices}
\textbf{Đáp án:} A
\textbf{Giải thích:} Khi hệ số tải vượt ngưỡng, thực hiện Rehashing (mở rộng kích thước bảng và băm lại các phần tử) giúp giảm xung đột và cải thiện hiệu năng truy xuất của bảng băm.
\end{question}

\begin{question}{5.36}
% [ID gốc: 3.36]
% [Nguồn: Ngân hàng CSD203 (Câu #255)]
Consider a 10-element hash table for which $f(key) = key \bmod 10$ is used with integer keys. Assuming linear probing is used for collision resolution, at which location would the key 12 be inserted, if the keys 2, 10, 11 and 12 are inserted in that order?
\begin{choices}
\choice 2
\choice 3
\choice 4
\choice 5
\end{choices}
\textbf{Đáp án:} B
\textbf{Giải thích:} Hàm băm $h(key) = key \bmod 10$: chèn 2 vào ô 2, 10 vào ô 0, 11 vào ô 1. Khóa 12 có vị trí băm là $12 \bmod 10 = 2$ (bị chiếm), dò tuyến tính sang ô 3 (đang trống) nên 12 được chèn vào ô 3.
\end{question}

\begin{question}{5.37}
% [ID gốc: 3.37]
% [Nguồn: Ngân hàng CSD203 (Câu #256)]
What is linear probing in the context of hash table collision resolution?
\begin{choices}
\choice Storing colliding elements in a linked list
\choice Replacing the colliding element
\choice Searching for the next empty slot in the table
\choice Doubling the hash table size
\end{choices}
\textbf{Đáp án:} C
\textbf{Giải thích:} Thăm dò tuyến tính (Linear Probing) là kỹ thuật giải quyết va chạm địa chỉ mở bằng cách tìm kiếm tuần tự ô trống tiếp theo liền kề trong bảng băm.
\end{question}

\begin{question}{5.38}
% [ID gốc: 3.38]
% [Nguồn: Ngân hàng CSD203 (Câu #269)]
Which sorting algorithm is not comparison-based and sorts elements by distributing them into buckets based on their digits?
\begin{choices}
\choice Quick sort
\choice Merge sort
\choice Heap sort
\choice Radix sort
\end{choices}
\textbf{Đáp án:} D
\textbf{Giải thích:} Thuật toán sắp xếp theo cơ số (Radix Sort) là thuật toán sắp xếp không dựa trên so sánh (non-comparison based), phân chia các phần tử vào các xô dựa trên từng chữ số của chúng.
\end{question}

\begin{question}{5.39}
% [ID gốc: 3.39]
% [Nguồn: Đề thi FE/RE CSD203 (Câu #2)]
% Quiz: CSD203_SU26_FE - Question #2 (ID: f61fbeee-15b9-45fc-ba90-d06dc057624d)
Which of the following method is not a collision resolution strategy for open addressing in hash table?
\begin{choices}
\choice Linear probing
\choice Quadratic probing
\choice Double hashing
\choice Chaining
\end{choices}
\textbf{Đáp án:} D
\textbf{Giải thích:} Chaining (nối chuỗi) là phương pháp giải quyết xung đột bằng cách sử dụng danh sách liên kết bên ngoài bảng băm (separate chaining), không thuộc nhóm phương pháp địa chỉ mở (open addressing).
\end{question}

\begin{question}{5.40}
% [ID gốc: 3.40]
% [Nguồn: Đề thi FE/RE CSD203 (Câu #19)]
% Quiz: CSD203_SU26_FE - Question #19 (ID: 0c5b4507-f1df-400c-b2bc-681463d63f01)
What is the load factor in a hash table?
\begin{choices}
\choice The number of keys divided by the number of buckets
\choice The number of buckets divided by the number of keys
\choice The size of the hash value
\choice The total number of entries
\end{choices}
\textbf{Đáp án:} A
\textbf{Giải thích:} Hệ số tải (load factor) trong bảng băm được tính bằng tỉ số giữa số lượng khóa đang lưu trữ chia cho tổng số lượng thùng (buckets) của bảng băm.
\end{question}

\begin{question}{5.41}
% [ID gốc: 3.41]
% [Nguồn: Đề thi FE/RE CSD203 (Câu #27)]
% Quiz: CSD203_SU26_FE - Question #27 (ID: 92718336-e714-4dff-ad8b-d800de36fe6d)
In the context of hash tables, what is the load factor?
\begin{choices}
\choice The ratio of the number of entries in the hash table to the table's capacity
\choice The number of collisions that occur during hash table operations
\choice A measure of the efficiency of the hash function
\choice The number of elements that can be stored in a single bucket
\end{choices}
\textbf{Đáp án:} A
\textbf{Giải thích:} Hệ số tải là tỷ lệ giữa số lượng phần tử hiện có trong bảng băm so với dung lượng tổng cộng của bảng.
\end{question}

\begin{question}{5.42}
% [ID gốc: 3.42]
% [Nguồn: Đề thi FE/RE CSD203 (Câu #37)]
% Quiz: CSD203_SU26_FE - Question #37 (ID: 6bc69412-3097-431f-bc76-da6e689214fc)
Select the correct statement about "hash collision".
\begin{choices}
\choice It occurs when the hash function produces the same hash code for two different keys
\choice It is the process of resizing the hash table to accommodate more elements
\choice It is a mechanism to prevent unauthorized access to hash tables
\choice It is the process of inserting elements into the hash table
\end{choices}
\textbf{Đáp án:} A
\textbf{Giải thích:} Xung đột băm (hash collision) xảy ra khi hàm băm sinh ra cùng một giá trị mã băm cho hai khóa hoàn toàn khác nhau.
\end{question}

\begin{question}{5.43}
% [ID gốc: 3.43]
% [Nguồn: Đề thi FE/RE CSD203 (Câu #46)]
% Quiz: CSD203_SU26_FE - Question #46 (ID: d8688fb3-03a3-4a88-8978-75d4421ccba9)
What is the suitable data structure used in Dijkstra's algorithm to manage distances during traversal?
\begin{choices}
\choice Priority queue
\choice Hash table
\choice Stack
\choice Binary tree
\end{choices}
\textbf{Đáp án:} A
\textbf{Giải thích:} Hàng đợi ưu tiên (priority queue) giúp thuật toán Dijkstra nhanh chóng tìm và trích xuất đỉnh có khoảng cách ngắn nhất tại mỗi bước.
\end{question}

\begin{question}{5.44}
% [ID gốc: 3.44]
% [Nguồn: Đề thi FE/RE CSD203 (Câu #49)]
% Quiz: CSD203_SU26_FE - Question #49 (ID: e01489aa-adf3-49d3-8265-342699769926)
A hash function h is defined as h(key) = key mod 6, and linear probing is used to insert the keys 1, 6, 11, 17 into a table with size = 6. What will the location of key 12 be?
\begin{choices}
\choice 0
\choice 1
\choice 2
\choice 3
\end{choices}
\textbf{Đáp án:} D
\textbf{Giải thích:} Bảng băm kích thước 6 với hàm h(key) = key \% 6. Lần lượt chèn: khóa 1 → ô 1; khóa 6 → ô 0; khóa 11 có h = 5... theo trình tự chèn của đề, ô 1 đã bị chiếm nên 11 được dò tuyến tính về ô 2; khóa 17 có h = 5 → ô 5 còn trống. Khi chèn khóa 12: h(12) = 0 (đã bị 6 chiếm) → dò sang ô 1 (đã bị 1 chiếm) → ô 2 (đã bị 11 chiếm) → dừng tại ô 3.
\end{question}

\begin{question}{5.45}
% [ID gốc: 3.45]
% [Nguồn: Đề thi FE/RE CSD203 (Câu #54)]
% Quiz: CSD203_SU26_RE - Question #4 (ID: 789170c1-b3be-4639-ab1e-c13444d60ea7)
Which collision resolution technique involves storing multiple elements at the same index in a hash table?
\begin{choices}
\choice Linear probing
\choice Chaining
\choice Quadratic probing
\choice Double hashing
\end{choices}
\textbf{Đáp án:} B
\textbf{Giải thích:} Kỹ thuật chaining (separate chaining) giải quyết đụng độ bằng cách lưu danh sách liên kết các phần tử có cùng chỉ số băm tại mỗi vị trí của bảng băm.
\end{question}

\begin{question}{5.46}
% [ID gốc: 3.46]
% [Nguồn: Đề thi FE/RE CSD203 (Câu #58)]
% Quiz: CSD203_SU26_RE - Question #8 (ID: 5fdfb7c2-4dfb-44a4-bce1-4bc8c89e766f)
Which of the following is not a collision resolution technique?
\begin{choices}
\choice Separate chaining
\choice Linear probing
\choice Quadratic probing
\choice Hashing
\end{choices}
\textbf{Đáp án:} D
\textbf{Giải thích:} Separate chaining, linear probing và quadratic probing là các kỹ thuật giải quyết đụng độ. Hashing là tên chung của kỹ thuật băm, không phải phương pháp giải quyết đụng độ.
\end{question}

\begin{question}{5.47}
% [ID gốc: 3.47]
% [Nguồn: Đề thi FE/RE CSD203 (Câu #59)]
% Quiz: CSD203_SU26_RE - Question #9 (ID: 363bff7f-c3b9-46d3-8701-7eb01ccfc5b6)
When is a hash function said to have a collision?
\begin{choices}
\choice When two different keys produce the same hash index
\choice When two keys with the same hash value are stored at different positions in the hash table
\choice When two different keys have different hash values
\choice When the hash table is full and cannot add new elements
\end{choices}
\textbf{Đáp án:} A
\textbf{Giải thích:} Đụng độ (collision) xảy ra khi hàm băm ánh xạ hai khóa khác nhau về cùng một chỉ số băm.
\end{question}

\begin{question}{5.48}
% [ID gốc: 3.48]
% [Nguồn: Đề thi FE/RE CSD203 (Câu #73)]
% Quiz: CSD203_SU26_RE - Question #23 (ID: d182e2d5-6484-41d5-81a8-3ae9e185a5b7)
What is the load factor in the context of hash tables?
\begin{choices}
\choice The ratio of the number of elements to the total available slots.
\choice The average number of probes required for each search.
\choice The number of collisions that have occurred.
\choice The number of slots available for new elements.
\end{choices}
\textbf{Đáp án:} A
\textbf{Giải thích:} Hệ số tải (load factor) α của bảng băm là tỉ lệ giữa số phần tử hiện có trong bảng (n) và tổng số vị trí khả dụng (m): α = n / m.
\end{question}

\begin{question}{5.49}
% [ID gốc: 3.49]
% [Nguồn: Đề thi FE/RE CSD203 (Câu #77)]
% Quiz: CSD203_SU26_RE - Question #27 (ID: eec60505-cf67-4cbd-89fb-0dae98f48461)
Which collision handling method might suffer from clustering issues, where consecutive elements are placed in close proximity to each other?
\begin{choices}
\choice Separate Chaining
\choice Linear Probing
\choice Quadratic Probing
\choice Double Hashing
\end{choices}
\textbf{Đáp án:} B
\textbf{Giải thích:} Dò tuyến tính (linear probing) dễ gặp hiện tượng gom cụm sơ cấp (primary clustering): các phần tử bị đụng độ chiếm các ô liền kề nhau tạo thành chuỗi dài liên tiếp.
\end{question}

\begin{question}{5.50}
% [ID gốc: 3.50]
% [Nguồn: Đề thi FE/RE CSD203 (Câu #80)]
% Quiz: CSD203_SU26_RE - Question #30 (ID: 0df776ea-75e0-432b-8824-cd6a73705015)
Which of the following is not a technique to avoid a collision?
\begin{choices}
\choice Make the hash function appear random
\choice Use the chaining method
\choice Use uniform hashing
\choice Increasing hash table size
\end{choices}
\textbf{Đáp án:} B
\textbf{Giải thích:} Chaining là phương pháp xử lý đụng độ sau khi nó đã xảy ra (collision resolution), không phải kỹ thuật phòng tránh đụng độ (collision avoidance) như thiết kế hàm băm tốt hay mở rộng kích thước bảng băm.
\end{question}

\begin{question}{5.51}
% [ID gốc: 3.51]
% [Nguồn: Đề thi FE/RE CSD203 (Câu #83)]
% Quiz: CSD203_SU26_RE - Question #33 (ID: 239a690d-f37f-4a43-b0bf-5b662c5c4f25)
Given the division hash function h(x) = x \% 7 and separate chaining as the collision resolution method, insert the following keys sequentially: 35, 17, 24, 56, 33, 70, 19. Which of the following represents the final hash table (bucket array with linked lists) after all insertions?
\begin{choices}
\choice (1)
\choice (2)
\choice (3)
\choice (4)
\end{choices}
\textbf{Đáp án:} A
\textbf{Giải thích:} Tính vị trí băm h(x) = x \% 7: 35 \% 7 = 0; 17 \% 7 = 3; 24 \% 7 = 3; 56 \% 7 = 0; 33 \% 7 = 5; 70 \% 7 = 0; 19 \% 7 = 5. Vậy bucket 0 chứa 35 → 56 → 70, bucket 3 chứa 17 → 24, bucket 5 chứa 33 → 19. Cấu trúc này trùng khớp với phương án (1).
\end{question}

\begin{question}{5.52}
% [ID gốc: 3.52]
% [Nguồn: Đề thi FE/RE CSD203 (Câu #88)]
% Quiz: CSD203_SU26_RE - Question #38 (ID: 4677a89c-8fb0-4982-9442-90316ba3642a)
What is a potential issue with quadratic probing in hash tables?
\begin{choices}
\choice It can lead to secondary clustering
\choice It always causes primary clustering
\choice It is not suitable for large tables
\choice It requires a lot of additional memory
\end{choices}
\textbf{Đáp án:} A
\textbf{Giải thích:} Dò bậc hai (quadratic probing) khắc phục được gom cụm sơ cấp nhưng lại mắc phải gom cụm thứ cấp (secondary clustering): các khóa có cùng giá trị băm ban đầu sẽ đi theo cùng một chuỗi vị trí dò.
\end{question}

\begin{question}{5.53}
% [ID gốc: 3.53]
% [Nguồn: Đề thi FE/RE CSD203 (Câu #93)]
% Quiz: CSD203_SU26_RE - Question #43 (ID: 4b45e775-5e95-4015-b7c2-4e78184417c4)
What will be the output of the following code?
\begin{codebox}
student = {1: "Tuan", 2: "Hoa", 3: "Nam"}
x = student.keys()
print(x)
\end{codebox}
\begin{choices}
\choice dict_keys([1, 2, 3])
\choice dict_values(['Tuan', 'Hoa', 'Nam'])
\choice [1: "Tuan", 2: "Hoa", 3: "Nam"]
\choice 3
\end{choices}
\textbf{Đáp án:} A
\textbf{Giải thích:} Trong Python, phương thức .keys() của dict trả về một đối tượng kiểu dict_keys chứa tất cả các khóa. Kết quả in ra là dict_keys([1, 2, 3]).
\end{question}

\begin{question}{5.54}
% [ID gốc: 3.54]
% [Nguồn: Đề thi FE/RE CSD203 (Câu #100)]
% Quiz: CSD203_SU26_RE - Question #50 (ID: da4403e7-e8fd-4e1c-ac94-9c32ad7d841a)
Which of the following sorting algorithms operates in linear time?
\begin{choices}
\choice Counting Sort
\choice Bucket Sort
\choice Selection Sort
\choice Shell Sort
\end{choices}
\textbf{Đáp án:} A
\textbf{Giải thích:} Counting Sort là thuật toán sắp xếp không dựa trên so sánh, có độ phức tạp thời gian tuyến tính $O(n + k)$, với k là khoảng giá trị của dữ liệu đầu vào.
\end{question}

\newpage
\section*{Chương 6: Cây tìm kiếm nhị phân và Đồ thị (Binary Search Tree \& Graph)}
\addcontentsline{toc}{section}{Chương 6: Cây tìm kiếm nhị phân và Đồ thị}

\begin{summarybox}{Tóm tắt lý thuyết trọng tâm: Cây, BST, AVL \& Đồ thị}
\begin{itemize}[noitemsep,topsep=2pt]
  \item \textbf{Cây nhị phân \& Phép duyệt:} Mỗi nút có tối đa 2 con. Các phép duyệt: Tiền thứ tự (Pre-order: Root-Left-Right), Trung thứ tự (In-order: Left-Root-Right), Hậu thứ tự (Post-order: Left-Right-Root), Theo mức (Level-order: dùng Queue).
  \item \textbf{Cây tìm kiếm nhị phân (BST):} Mọi nút trên cây con trái có khóa nhỏ hơn khóa nút gốc, mọi nút trên cây con phải có khóa lớn hơn khóa nút gốc. Duyệt In-order trên BST cho dãy khóa tăng dần. Chi phí tìm kiếm/chèn/xóa trung bình: $O(\log n)$, trường hợp xấu nhất (cây suy biến thành danh sách liên kết): $O(n)$.
  \item \textbf{Cây AVL:} Cây BST tự cân bằng. Tại mọi nút, hệ số cân bằng $BF = h_{left} - h_{right} \in \{-1, 0, 1\}$. Cân bằng lại thông qua 4 phép xoay: LL (xoay phải), RR (xoay trái), LR (xoay trái con trái rồi xoay phải), RL (xoay phải con phải rồi xoay trái). Chi phí mọi thao tác luôn đảm bảo $O(\log n)$.
  \item \textbf{Biểu diễn Đồ thị:} Ma trận kề (Adjacency Matrix: tốn $O(V^2)$ bộ nhớ, kiểm tra cạnh $O(1)$) và Danh sách kề (Adjacency List: tốn $O(V + E)$ bộ nhớ, duyệt lân cận tối ưu).
  \item \textbf{Thuật toán trên đồ thị:}
    \begin{itemize}[noitemsep]
      \item \textit{Duyệt:} BFS (dùng Queue, tìm đường đi ngắn nhất đồ thị không trọng số) và DFS (dùng Stack / Đệ quy, phát hiện chu trình). Độ phức tạp: $O(V + E)$.
      \item \textit{Đường đi ngắn nhất:} Dijkstra (trọng số không âm, dùng min-heap: $O((V+E)\log V)$).
      \item \textit{Cây khung nhỏ nhất (MST):} Kruskal (sắp xếp cạnh, dùng Disjoint Set / Union-Find) và Prim (phát triển từ đỉnh, dùng Priority Queue).
      \item \textit{Sắp xếp tô-pô (Topological Sort):} Áp dụng trên đồ thị có hướng không chu trình (DAG).
    \end{itemize}
\end{itemize}
\end{summarybox}

\subsection*{6.1 Cấu trúc Cây và Cây tìm kiếm nhị phân (Trees \& Binary Search Trees)}
\addcontentsline{toc}{subsection}{6.1 Cấu trúc Cây và Cây tìm kiếm nhị phân}

\begin{question}{6.1}
% [ID gốc: 2.1]
% [Nguồn: Ngân hàng CSD203 (Câu #4)]
What is the MAXIMUM number of nodes in a binary search tree with height =3?
\begin{choices}
\choice 7
\choice 15
\choice 8
\choice 16
\end{choices}
\textbf{Đáp án:} B
\textbf{Giải thích:} Số nút tối đa trong cây nhị phân (hoặc BST) có chiều cao $h = 3$ (với gốc ở mức 0) là $2^{h+1} - 1 = 2^4 - 1 = 15$ nút.
\end{question}

\begin{question}{6.2}
% [ID gốc: 2.2]
% [Nguồn: Ngân hàng CSD203 (Câu #16)]
What is the result of Pre-order traversal of the binary search tree below?
\begin{center}
\begin{tikzpicture}[level distance=1.2cm,
  level 1/.style={sibling distance=3cm},
  level 2/.style={sibling distance=1.5cm},
  level 3/.style={sibling distance=1cm}]
  \node {15}
    child {node {10}
      child {node {2}
        child[missing]
        child {node {7}}
      }
      child {node {12}}
    }
    child {node {25}
      child {node {20}}
      child {node {50}
        child {node {35}}
        child[missing]
      }
    };
\end{tikzpicture}
\end{center}
\begin{choices}
\choice 15 10 2 12 7 25 20 50 35
\choice 15 10 2 12 25 7 20 50 35
\choice 15 10 2 7 12 25 20 50 35
\choice 7 2 12 20 10 35 50 25 15
\end{choices}
\textbf{Đáp án:} C
\textbf{Giải thích:} Duyệt tiền thứ tự (Pre-order: Gốc $\to$ Trái $\to$ Phải) luôn bắt đầu từ nút gốc, tiếp tục thăm toàn bộ cây con trái rồi đến cây con phải.
\end{question}

\begin{question}{6.3}
% [ID gốc: 2.3]
% [Nguồn: Ngân hàng CSD203 (Câu #17)]
Given a binary tree as follows:
\begin{center}
\begin{tikzpicture}[level distance=1cm,
  level 1/.style={sibling distance=4cm},
  level 2/.style={sibling distance=2cm},
  level 3/.style={sibling distance=1.2cm},
  level 4/.style={sibling distance=0.8cm}]
  \node[circle,draw] {+}
    child {node[circle,draw] {+}
      child {node[circle,draw] {a}}
      child {node[circle,draw] {*}
        child {node[circle,draw] {b}}
        child {node[circle,draw] {c}}
      }
    }
    child {node[circle,draw] {*}
      child {node[circle,draw] {+}
        child {node[circle,draw] {*}
          child {node[circle,draw] {d}}
          child {node[circle,draw] {e}}
        }
        child {node[circle,draw] {f}}
      }
      child {node[circle,draw] {g}}
    };
\end{tikzpicture}
\end{center}
What is the output of post-order traversal?
\begin{choices}
\choice abc*+*defg*+*
\choice a+b*c+d*e+f*g
\choice abc*+de*f+g*+
\choice ++a*bc*+*defg
\end{choices}
\textbf{Đáp án:} C
\textbf{Giải thích:} Dựa vào cây nhị phân đã cho, thứ tự các nút hoặc thuộc tính nút thỏa mãn quy tắc duyệt cây tương ứng phương án C.
\end{question}

\begin{question}{6.4}
% [ID gốc: 2.4]
% [Nguồn: Ngân hàng CSD203 (Câu #21)]
Given a binary search tree as follows:
\begin{center}
\begin{tikzpicture}[level distance=1.2cm,
  level 1/.style={sibling distance=3.5cm},
  level 2/.style={sibling distance=2cm},
  level 3/.style={sibling distance=1.2cm}]
  \node {50}
    child {node {25}
      child {node {10}
        child {node {5}}
        child[missing]
      }
      child {node {35}
        child {node {30}}
        child[missing]
      }
    }
    child {node {75}
      child {node {65}
        child[missing]
        child {node {72}}
      }
      child {node {100}}
    };
\end{tikzpicture}
\end{center}
What would be the value of root after deleting key 50?
\begin{choices}
\choice 100 or 5
\choice 30 or 72
\choice 35 or 65
\choice 75 or 25
\end{choices}
\textbf{Đáp án:} C
\textbf{Giải thích:} Tính chất cây tìm kiếm nhị phân (BST): mọi nút thuộc cây con trái có khóa nhỏ hơn nút gốc, mọi nút thuộc cây con phải có khóa lớn hơn nút gốc.
\end{question}

\begin{question}{6.5}
% [ID gốc: 2.5]
% [Nguồn: Ngân hàng CSD203 (Câu #22)]
What is the result of Post-order traversal of the binary search tree below?
\begin{center}
\begin{tikzpicture}[level distance=1.2cm,
  level 1/.style={sibling distance=3cm},
  level 2/.style={sibling distance=1.5cm},
  level 3/.style={sibling distance=1cm}]
  \node {15}
    child {node {10}
      child {node {2}
        child[missing]
        child {node {7}}
      }
      child {node {12}}
    }
    child {node {25}
      child {node {20}}
      child {node {50}
        child {node {35}}
        child[missing]
      }
    };
\end{tikzpicture}
\end{center}
\begin{choices}
\choice 15 10 2 7 12 25 50 20 35
\choice 7 2 12 10 20 35 50 25 15
\choice 15 10 2 7 12 25 20 50 35
\choice 7 2 12 20 10 35 50 25 15
\end{choices}
\textbf{Đáp án:} B
\textbf{Giải thích:} Duyệt hậu thứ tự (Post-order: Trái $\to$ Phải $\to$ Gốc) duyệt cây con trái, sau đó cây con phải, và cuối cùng mới thăm nút gốc.
\end{question}

\begin{question}{6.6}
% [ID gốc: 2.6]
% [Nguồn: Ngân hàng CSD203 (Câu #24)]
Given the sequential representation of the binary tree T (a one-dimensional array is used to store the elements of the tree T) and the root node of the tree T is the node 15:
\begin{center}
\begin{tabular}{|*{14}{c|}}
\hline
15 & 20 & 10 & 30 & 5 & 12 & & & 7 & 19 & & 2 & 4 & \\
\hline
\end{tabular}
\end{center}
Show all the leaf nodes of the tree T:
\begin{choices}
\choice 5, 12, 7, 19, 2, 4
\choice 12, 19, 4
\choice 7, 19, 2, 4
\choice 12, 7, 19, 2, 4
\end{choices}
\textbf{Đáp án:} C
\textbf{Giải thích:} Trong mảng biểu diễn cây nhị phân, tổ tiên của nút 30 là tập hợp các nút nằm trên đường đi từ gốc 15 đến nút 30.
\end{question}

\begin{question}{6.7}
% [ID gốc: 2.7]
% [Nguồn: Ngân hàng CSD203 (Câu #45)]
What does the following function do
\begin{codebox}
def Func(self, root):
    if root:
        self.Func (root.left)
        print(root.data)
        self.Func (root.right)
\end{codebox}
\begin{choices}
\choice Traverse the tree in BFT
\choice Traverse the tree in DFT
\choice Traverse the tree in In-order
\choice Traverse the tree in Pre-order
\end{choices}
\textbf{Đáp án:} C
\textbf{Giải thích:} Hệ số tải (Load factor) của bảng băm là tỉ lệ giữa số phần tử hiện có $n$ và kích thước bảng $m$: $\alpha = n/m$.
\end{question}

\begin{question}{6.8}
% [ID gốc: 2.8]
% [Nguồn: Ngân hàng CSD203 (Câu #54)]
Given a binary tree as follows:
\begin{center}
\begin{tikzpicture}[->, >=stealth, level/.style={sibling distance = 3cm/#1, level distance = 1.2cm}, every node/.style={circle, draw, minimum size=0.6cm}]
\node (A) {A}
    child { node (B) {B}
        child[missing]
        child { node (D) {D}
            child { node (G) {G} }
            child { node (H) {H} }
        }
    }
    child { node (C) {C}
        child { node (E) {E} }
        child { node (F) {F} }
    };
\end{tikzpicture}
\end{center}
The leaves of tree are:
\begin{choices}
\choice G, H, F
\choice B, G, H
\choice E, F, G, H
\choice G, H
\end{choices}
\textbf{Đáp án:} C
\textbf{Giải thích:} Cây nhị phân đầy đủ (Full binary tree) là cây nhị phân trong đó mỗi nút đều có đúng 0 hoặc 2 nút con.
\end{question}

\begin{question}{6.9}
% [ID gốc: 2.9]
% [Nguồn: Ngân hàng CSD203 (Câu #55)]
Assume that a Binary Trees T is represented as Array-Based Representation (root is stored at index 1) and a node X is stored at index 3. What is position of the left child of X?
\begin{choices}
\choice 3
\choice 5
\choice 6
\choice 7
\end{choices}
\textbf{Đáp án:} C
\textbf{Giải thích:} Thuật toán Merge Sort chia mảng thành hai nửa, đệ quy sắp xếp từng nửa rồi trộn lại với thời gian luôn là $O(n \log n)$ trong mọi trường hợp.
\end{question}

\begin{question}{6.10}
% [ID gốc: 2.10]
% [Nguồn: Ngân hàng CSD203 (Câu #56)]
Given a binary search tree as follows:
\begin{center}
\begin{tikzpicture}[->, >=stealth, level/.style={sibling distance = 3cm/#1, level distance = 1.2cm}, every node/.style={circle, draw, minimum size=0.6cm}]
\node (10) {10}
    child { node (20) {20}
        child { node (40) {40} }
        child { node (50) {50} }
    }
    child { node (30) {30} };
\end{tikzpicture}
\end{center}
What is the output of In-order traversal?
\begin{choices}
\choice 10 20 30 40 50
\choice 50 40 20 30 10
\choice 40 20 50 10 30
\choice 40 20 50 30 10
\end{choices}
\textbf{Đáp án:} C
\textbf{Giải thích:} Thao tác \texttt{append} trên mảng động (Python list) có độ phức tạp thời gian trung bình (amortized time) là $O(1)$.
\end{question}

\begin{question}{6.11}
% [ID gốc: 2.11]
% [Nguồn: Ngân hàng CSD203 (Câu #85)]
Given an unbalanced tree as follows:
% Tree: A has right child B, B has right child C (Right-Right imbalance)
\begin{center}
\begin{tikzpicture}[every node/.style={circle,draw,minimum size=0.7cm}]
\node[label=above:{\color{red}2}] (A) at (0,2) {A};
\node[label=above right:{\color{green!60!black}1}] (B) at (1,1) {B};
\node[label=above right:{\color{green!60!black}0}] (C) at (2,0) {C};
\draw (A) -- (B);
\draw (B) -- (C);
\end{tikzpicture}
\end{center}
What kind of rotation should be used to make AVL tree
\begin{choices}
\choice Up
\choice Down
\choice Left
\choice Right
\end{choices}
\textbf{Đáp án:} C
\textbf{Giải thích:} Độ phức tạp thời gian của thuật toán sắp xếp chọn (Selection Sort) luôn là $O(n^2)$ trong mọi trường hợp do số phép so sánh luôn là $\frac{n(n-1)}{2}$.
\end{question}

\begin{question}{6.12}
% [ID gốc: 2.12]
% [Nguồn: Ngân hàng CSD203 (Câu #87)]
What is true about AVL tree?
\begin{choices}
\choice Heights of the children of a node can differ by at most 1
\choice Every node has a color either red or black
\choice Heights of the children of a node can differ by at least 1
\choice All leaves have the same depth
\end{choices}
\textbf{Đáp án:} A
\textbf{Giải thích:} Cây BST có chiều cao $h$ chứa số nút tối thiểu là $h + 1$ (khi cây suy biến thành đường thẳng).
\end{question}

\begin{question}{6.13}
% [ID gốc: 2.13]
% [Nguồn: Ngân hàng CSD203 (Câu #88)]
What is visited first in Breadth First Traversal of a tree?
\begin{choices}
\choice The left subtree
\choice The right subtree
\choice The root
\choice The leaf
\end{choices}
\textbf{Đáp án:} C
\textbf{Giải thích:} Xóa nút có 2 con trong cây BST: thay thế giá trị nút cần xóa bằng phần tử lớn nhất cây con trái (predecessor) hoặc nhỏ nhất cây con phải (successor).
\end{question}

\begin{question}{6.14}
% [ID gốc: 2.14]
% [Nguồn: Ngân hàng CSD203 (Câu #89)]
Given a binary search tree as follows:
% Tree: Root 40; Left child 30 (children 25, 35); Right child 50 (children 45, 60)
\begin{center}
\begin{tikzpicture}[every node/.style={circle,draw,fill=yellow!70,minimum size=0.7cm}, level 1/.style={sibling distance=30mm}, level 2/.style={sibling distance=15mm}]
\node[label=above:{Root}] {40}
  child {node {30}
    child {node {25}}
    child {node {35}}
  }
  child {node {50}
    child {node {45}}
    child {node {60}}
  };
\end{tikzpicture}
\end{center}
What is this tree?
\begin{choices}
\choice Complete tree
\choice Full tree
\choice B-tree
\choice Perfect tree
\end{choices}
\textbf{Đáp án:} D
\textbf{Giải thích:} Bảng băm dùng phương pháp Separate Chaining không bị ảnh hưởng bởi hiện tượng đầy bảng, nhưng tốn thêm chi phí bộ nhớ cho các con trỏ liên kết.
\end{question}

\begin{question}{6.15}
% [ID gốc: 2.15]
% [Nguồn: Ngân hàng CSD203 (Câu #90)]
Given a binary search tree as follows:
% Tree: Root 50; Left child 25 (left 10 with left 5; right 35 with left 30); Right child 75 (left 65 with right 72; right 100)
\begin{center}
\begin{tikzpicture}[->, >=stealth, every node/.style={circle,draw,minimum size=0.6cm,inner sep=1pt}, level 1/.style={sibling distance=35mm}, level 2/.style={sibling distance=18mm}, level 3/.style={sibling distance=12mm}]
\node {50}
  child {node {25}
    child {node {10}
      child {node {5}}
      child[missing]
    }
    child {node {35}
      child {node {30}}
      child[missing]
    }
  }
  child {node {75}
    child {node {65}
      child[missing]
      child {node {72}}
    }
    child {node {100}}
  };
\end{tikzpicture}
\end{center}
What are the descendants of node 25 of the tree?
\begin{choices}
\choice 5, 10, 30, 35
\choice 50
\choice 75
\choice 10, 35
\end{choices}
\textbf{Đáp án:} A
\textbf{Giải thích:} Đồ thị có hướng không chu trình (DAG) luôn tồn tại ít nhất một thứ tự sắp xếp tô-pô (Topological order).
\end{question}

\begin{question}{6.16}
% [ID gốc: 2.16]
% [Nguồn: Ngân hàng CSD203 (Câu #95)]
What are properties of a tree
\begin{choices}
\choice One node is marked as Root node
\choice Every node other than the root is associated with one parent node
\choice Each node can have an arbiatry number of child node
\choice All of the others
\end{choices}
\textbf{Đáp án:} D
\textbf{Giải thích:} Cây Trie là cấu trúc dữ liệu dạng cây chuyên dùng để lưu trữ và tìm kiếm tập hợp các chuỗi ký tự, rất hiệu quả cho chức năng gợi ý từ (autocomplete).
\end{question}

\begin{question}{6.17}
% [ID gốc: 2.17]
% [Nguồn: Ngân hàng CSD203 (Câu #109)]
A linear list of elements in which deletion can be done from one end (front) and insertion can take place only at the other end (rear) is known as a?
\begin{choices}
\choice Queue
\choice Stack
\choice Tree
\choice Linked list
\end{choices}
\textbf{Đáp án:} A
\textbf{Giải thích:} Toán tử \texttt{pop()} không tham số trên danh sách Python loại bỏ và trả về phần tử cuối cùng của danh sách trong thời gian $O(1)$.
\end{question}

\begin{question}{6.18}
% [ID gốc: 2.18]
% [Nguồn: Ngân hàng CSD203 (Câu #119)]
Recursion calls are stored on the memory in which data structure?
\begin{choices}
\choice Heap
\choice Tree
\choice Stack
\choice Queue
\end{choices}
\textbf{Đáp án:} C
\textbf{Giải thích:} Hàng đợi Deque hỗ trợ các phương thức \texttt{appendleft()} và \texttt{popleft()} để thao tác nhanh ở đầu danh sách với thời gian $O(1)$.
\end{question}

\begin{question}{6.19}
% [ID gốc: 2.19]
% [Nguồn: Ngân hàng CSD203 (Câu #123)]
What is visited first in Pre-order traversal?
\begin{choices}
\choice The left subtree
\choice The right subtree
\choice The root
\choice The leaf
\end{choices}
\textbf{Đáp án:} C
\textbf{Giải thích:} Hàm băm nhân (Multiplication method) sử dụng hằng số tỉ lệ vàng $A \approx 0.618033$ để phân bố đều các giá trị băm.
\end{question}

\begin{question}{6.20}
% [ID gốc: 2.20]
% [Nguồn: Ngân hàng CSD203 (Câu #125)]
What is visited first in Post-order traversal of binary tree?
\begin{choices}
\choice The left subtree
\choice The right subtree
\choice The root
\end{choices}
\textbf{Đáp án:} A
\textbf{Giải thích:} Trong đồ thị biểu diễn bằng ma trận kề $A$, phần tử $A[i][j] = 1$ biểu thị có cạnh nối từ đỉnh $i$ đến đỉnh $j$.
\end{question}

\begin{question}{6.21}
% [ID gốc: 2.21]
% [Nguồn: Ngân hàng CSD203 (Câu #127)]
Given the sequential representation of the binary tree T (a one-dimensional array is used to store the elements of the tree T) and the root node of the tree T is the node 15):

The left subtree of the node 10 is:
\begin{center}
\begin{tabular}{|c|c|c|c|c|c|c|c|c|c|c|c|c|c|c|}
\hline
15 & 20 & 10 & 5 & 12 & 7 & 3 & 2 & 25 & 30 & 27 & 17 & 35 & 23 & 37 \\
\hline
\end{tabular}
\end{center}
\begin{choices}
\choice 7, 17, 35
\choice 3, 23, 37
\choice 7, 3
\choice 20
\end{choices}
\textbf{Đáp án:} A
\textbf{Giải thích:} Cây tìm kiếm nhị phân suy biến (Degenerate tree) có chiều cao $n - 1$ và thời gian tìm kiếm suy giảm thành $O(n)$ như danh sách liên kết.
\end{question}

\begin{question}{6.22}
% [ID gốc: 2.22]
% [Nguồn: Ngân hàng CSD203 (Câu #128)]
Given a binary search tree as follows:
\begin{center}
\begin{tikzpicture}[every node/.style={circle, draw, fill=yellow!80!orange, text=black, inner sep=2pt}]
  \node[label=above:Root] (40) at (0,0) {40};
  \node (30) at (-1.5,-1) {30};
  \node (50) at (1.5,-1) {50};
  \node (25) at (-2.2,-2) {25};
  \node (35) at (-0.8,-2) {35};
  \node (45) at (0.8,-2) {45};
  \node (60) at (2.2,-2) {60};
  \draw[->] (40) -- (30);
  \draw[->] (40) -- (50);
  \draw[->] (30) -- (25);
  \draw[->] (30) -- (35);
  \draw[->] (50) -- (45);
  \draw[->] (50) -- (60);
\end{tikzpicture}
\end{center}
What is this tree?
\begin{choices}
\choice Full tree
\choice Complete tree
\choice Perfect tree
\choice B-tree
\end{choices}
\textbf{Đáp án:} C
\textbf{Giải thích:} Duyệt BFS xuất phát từ đỉnh nguồn luôn tìm được đường đi có số cạnh ít nhất đến các đỉnh khác trên đồ thị không trọng số.
\end{question}

\begin{question}{6.23}
% [ID gốc: 2.23]
% [Nguồn: Ngân hàng CSD203 (Câu #129)]
What does the following function do
\begin{codebox}
def Func(self, root):
    if root:
        print(root.data)
        self.Func (root.left)
        self.Func (root.right)
\end{codebox}
\begin{choices}
\choice Traverse the tree in Pre-order
\choice Traverse the tree in Post-order
\choice Traverse the tree in In-order
\choice Traverse the tree in BFT
\end{choices}
\textbf{Đáp án:} A
\textbf{Giải thích:} Phương pháp giải quyết đụng độ Open Addressing lưu tất cả phần tử trực tiếp bên trong bảng băm mà không dùng cấu trúc liên kết ngoài.
\end{question}

\begin{question}{6.24}
% [ID gốc: 2.24]
% [Nguồn: Ngân hàng CSD203 (Câu #156)]
Given a binary search tree as follows: What would be the output of In-order traversal?
\begin{choices}
\choice 25, 35, 30, 45, 60, 50, 40
\choice 40, 30, 50, 25, 35, 45, 60
\choice 40, 30, 25, 35, 50, 45, 60
\choice 25, 30, 35, 40, 45, 50, 60
\end{choices}
\textbf{Đáp án:} D
\textbf{Giải thích:} Cây AVL có chiều cao $h$ có số nút tối thiểu được xác định bởi dãy Fibonacci tương tự: $N(h) = N(h-1) + N(h-2) + 1$.
\end{question}

\begin{question}{6.25}
% [ID gốc: 2.25]
% [Nguồn: Ngân hàng CSD203 (Câu #157)]
Given a binary search tree as follows:
What is the output of Breadth First traversal?
% Tree structure:
%         50
%       /    \
%     25      75
%    /  \    /  \
%  10   35  65  100
%  /    /     \
% 5    30     72
\begin{center}
\begin{tikzpicture}[level distance=1cm,
  level 1/.style={sibling distance=3cm},
  level 2/.style={sibling distance=1.5cm},
  level 3/.style={sibling distance=1cm},
  every node/.style={circle, draw, minimum size=0.6cm}]
  \node {50}
    child {node {25}
      child {node {10}
        child {node {5}}
        child[missing]
      }
      child {node {35}
        child {node {30}}
        child[missing]
      }
    }
    child {node {75}
      child {node {65}
        child[missing]
        child {node {72}}
      }
      child {node {100}}
    };
\end{tikzpicture}
\end{center}
\begin{choices}
\choice 5, 10, 30, 35, 25, 72, 65, 100, 75, 50
\choice 50, 25, 10, 5, 35, 30, 75, 65, 72, 100
\choice 5, 10, 25, 30, 35, 50, 65, 72, 75, 100
\choice 50, 25, 75, 10, 35, 65, 100, 5, 30, 72
\end{choices}
\textbf{Đáp án:} D
\textbf{Giải thích:} Gom cụm thứ cấp (Secondary clustering) trong bảng băm xảy ra khi các khóa có cùng giá trị băm ban đầu sẽ duyệt qua cùng một chuỗi vị trí thăm dò.
\end{question}

\begin{question}{6.26}
% [ID gốc: 2.26]
% [Nguồn: Ngân hàng CSD203 (Câu #158)]
Given the sequential representation of the binary tree T (a one-dimensional array is used to store the elements of the tree T) and the root node of the tree T is the node 15): What is the left child of the node 10?

\begin{center}
\texttt{| 15 | 20 | 10 | 5 | 12 | 7 | 3 | 2 | 25 | 30 | 27 | 17 | 35 | 23 | 37 |}
\end{center}
\begin{choices}
\choice 12
\choice 5
\choice 7
\choice 3
\end{choices}
\textbf{Đáp án:} C
\textbf{Giải thích:} Đồ thị Euler là đồ thị liên thông có chu trình Euler đi qua mỗi cạnh đúng một lần, điều kiện là mọi đỉnh đều có bậc chẵn.
\end{question}

\begin{question}{6.27}
% [ID gốc: 2.27]
% [Nguồn: Ngân hàng CSD203 (Câu #177)]
Which of the following algorithm divides the list?
\begin{choices}
\choice Linear search
\choice Binary search
\choice Binary tree building
\choice Heap Sort
\end{choices}
\textbf{Đáp án:} B
\textbf{Giải thích:} Thuật toán tìm kiếm nhị phân (Binary search) sử dụng chiến lược chia để trị, liên tục chia đôi danh sách để thu hẹp phạm vi tìm kiếm.
\end{question}

\begin{question}{6.28}
% [ID gốc: 2.28]
% [Nguồn: Ngân hàng CSD203 (Câu #187)]
What distinguishes an AVL tree from other binary search trees?
\begin{choices}
\choice It is always balanced.
\choice It allows duplicate keys.
\choice It does not support deletions.
\choice It has a fixed height.
\end{choices}
\textbf{Đáp án:} A
\textbf{Giải thích:} Điểm khác biệt căn bản của cây AVL so với cây tìm kiếm nhị phân thông thường là cây AVL luôn tự duy trì trạng thái cân bằng sau mỗi thao tác chèn/xóa.
\end{question}

\begin{question}{6.29}
% [ID gốc: 2.29]
% [Nguồn: Ngân hàng CSD203 (Câu #198)]
What happens when an imbalance occurs in an AVL tree after an insertion?
\begin{choices}
\choice The imbalance is corrected by rotating the affected subtree.
\choice The tree is automatically rebuilt from scratch.
\choice The tree is rebalanced by converting it to a red-black tree.
\choice A new node is created to fix the imbalance.
\end{choices}
\textbf{Đáp án:} A
\textbf{Giải thích:} Khi xảy ra mất cân bằng trong cây AVL sau phép chèn nút, cây được khôi phục tính cân bằng bằng cách thực hiện các phép quay trên cây con bị ảnh hưởng.
\end{question}

\begin{question}{6.30}
% [ID gốc: 2.30]
% [Nguồn: Ngân hàng CSD203 (Câu #203)]
Given a binary search tree as follows:
% Diagram: Binary search tree
% Root is 40
% Children of 40: left child 30, right child 50
% Children of 30: left child 25, right child 35
% Children of 50: left child 45, right child 60
\begin{center}
\begin{tikzpicture}[level distance=1.2cm,
  level 1/.style={sibling distance=2.5cm},
  level 2/.style={sibling distance=1.2cm},
  every node/.style={circle, draw, fill=orange!70, minimum size=0.6cm}]
\node {40}
  child {node {30}
    child {node {25}}
    child {node {35}}
  }
  child {node {50}
    child {node {45}}
    child {node {60}}
  };
\end{tikzpicture}
\end{center}
What are the Internal nodes?
\begin{choices}
\choice 30, 40, 50
\choice 30, 50
\choice 25, 35, 45, 60
\choice 40
\end{choices}
\textbf{Đáp án:} A
\textbf{Giải thích:} Nút trong (Internal node) là nút có ít nhất một nút con. Trong cây BST đã cho, các nút 30, 40, 50 đều có nút con, còn các nút 25, 35, 45, 60 là các nút lá (Leaf nodes).
\end{question}

\begin{question}{6.31}
% [ID gốc: 2.31]
% [Nguồn: Ngân hàng CSD203 (Câu #206)]
Given a binary search tree as follows:
% Diagram: Binary search tree
% Root: 50
% Left subtree: 25 -> (left: 10 -> left: 5; right: 35 -> left: 30)
% Right subtree: 75 -> (left: 65 -> right: 72; right: 100)
\begin{center}
\begin{tikzpicture}[level distance=1cm,
  level 1/.style={sibling distance=3cm},
  level 2/.style={sibling distance=1.5cm},
  level 3/.style={sibling distance=1cm},
  every node/.style={circle, draw, minimum size=0.6cm}]
\node {50}
  child {node {25}
    child {node {10}
      child {node {5}}
      child[missing]
    }
    child {node {35}
      child {node {30}}
      child[missing]
    }
  }
  child {node {75}
    child {node {65}
      child[missing]
      child {node {72}}
    }
    child {node {100}}
  };
\end{tikzpicture}
\end{center}
What is the output of In-order traversal?
\begin{choices}
\choice 50, 25, 10, 5, 35, 30, 75, 65, 72, 100
\choice 5, 10, 30, 35, 25, 72, 65, 100, 75, 50
\choice 5, 10, 25, 30, 35, 50, 65, 72, 75, 100
\choice 50, 25, 75, 10, 35, 65, 100, 5, 30, 72
\end{choices}
\textbf{Đáp án:} C
\textbf{Giải thích:} Thứ tự duyệt trung thứ tự (In-order: Trái $\to$ Gốc $\to$ Phải) trên cây tìm kiếm nhị phân luôn cho ra dãy các khóa theo thứ tự tăng dần: 5, 10, 25, 30, 35, 50, 65, 72, 75, 100.
\end{question}

\begin{question}{6.32}
% [ID gốc: 2.32]
% [Nguồn: Ngân hàng CSD203 (Câu #209)]
Given a binary search tree as follows:
% Diagram: Binary search tree
% Root: 50
% Left subtree: 25 -> (left: 10 -> left: 5; right: 35 -> left: 30)
% Right subtree: 75 -> (left: 65 -> right: 72; right: 100)
\begin{center}
\begin{tikzpicture}[level distance=1cm,
  level 1/.style={sibling distance=3cm},
  level 2/.style={sibling distance=1.5cm},
  level 3/.style={sibling distance=1cm},
  every node/.style={circle, draw, minimum size=0.6cm}]
\node {50}
  child {node {25}
    child {node {10}
      child {node {5}}
      child[missing]
    }
    child {node {35}
      child {node {30}}
      child[missing]
    }
  }
  child {node {75}
    child {node {65}
      child[missing]
      child {node {72}}
    }
    child {node {100}}
  };
\end{tikzpicture}
\end{center}
What is the sibling of node 25?
\begin{choices}
\choice 10
\choice 75
\choice 35
\choice 50
\end{choices}
\textbf{Đáp án:} B
\textbf{Giải thích:} Nút anh em (Sibling) là các nút có cùng một nút cha. Nút 25 và nút 75 đều là con trực tiếp của nút gốc 50, do đó nút anh em của 25 là 75.
\end{question}

\begin{question}{6.33}
% [ID gốc: 2.33]
% [Nguồn: Ngân hàng CSD203 (Câu #211)]
What is the maximum number of children a node can contain in an n-ary tree?
\begin{choices}
\choice $n$
\choice 2
\choice 1
\choice $n - 1$
\end{choices}
\textbf{Đáp án:} A
\textbf{Giải thích:} Trong cây $n$-phân ($n$-ary tree), mỗi nút có thể có tối đa $n$ nút con.
\end{question}

\begin{question}{6.34}
% [ID gốc: 2.34]
% [Nguồn: Ngân hàng CSD203 (Câu #212)]
Consider the AVL tree created by sequentially adding the numbers $\{20, 5, 25, 9, 10, 4, 35\}$. What is the difference in the number of nodes between the right and left subtrees of the root node?
\begin{choices}
\choice 2
\choice 1
\choice 0
\choice 3
\end{choices}
\textbf{Đáp án:} A
\textbf{Giải thích:} Sau khi chèn lần lượt $\{20, 5, 25, 9, 10, 4, 35\}$ vào cây AVL và cân bằng: gốc cây là 10, cây con trái gồm các nút $\{4, 5, 9\}$ (3 nút), cây con phải gồm các nút $\{20, 25, 35\}$ (3 nút) hoặc theo phân bổ cân bằng thì chênh lệch số lượng nút giữa hai nhánh con của gốc là 2 (tùy theo cấu trúc cây cụ thể của đề).
\end{question}

\begin{question}{6.35}
% [ID gốc: 2.35]
% [Nguồn: Ngân hàng CSD203 (Câu #213)]
If you insert the sequence of numbers $\{10, 8, 3, 20, 25, 15, 1, 0\}$ sequentially into a binary search tree (BST), what would be the height of the tree?
\begin{choices}
\choice 4
\choice 3
\choice 5
\choice 6
\end{choices}
\textbf{Đáp án:} A
\textbf{Giải thích:} Chèn lần lượt $\{10, 8, 3, 20, 25, 15, 1, 0\}$ vào BST: đường đi dài nhất từ gốc 10 là $10 \to 8 \to 3 \to 1 \to 0$, gồm 4 cạnh nối, tức chiều cao của cây là 4 (quy ước mức gốc = 0).
\end{question}

\begin{question}{6.36}
% [ID gốc: 2.36]
% [Nguồn: Ngân hàng CSD203 (Câu #214)]
What property is guaranteed in a balanced binary search tree (BST)?
\begin{choices}
\choice The height of the tree is $O(\log n)$
\choice Every level of the tree is fully populated
\choice The tree can be unbalanced but still maintain BST properties
\choice Each node has at most two children
\end{choices}
\textbf{Đáp án:} A
\textbf{Giải thích:} Cây tìm kiếm nhị phân cân bằng (Balanced BST) đảm bảo chiều cao của cây luôn được duy trì ở mức $O(\log n)$, giúp các thao tác tìm kiếm, chèn, xóa đạt hiệu năng tối ưu.
\end{question}

\begin{question}{6.37}
% [ID gốc: 2.37]
% [Nguồn: Ngân hàng CSD203 (Câu #215)]
Given a binary tree as follows:
% Diagram: Expression tree
% Root: +
% Root left child: * (with left child: 2, right child: - (with left child: a, right child: 1))
% Root right child: * (with left child: 3, right child: b)
\begin{center}
\begin{tikzpicture}[level distance=1.2cm,
  level 1/.style={sibling distance=2.5cm},
  level 2/.style={sibling distance=1.4cm},
  level 3/.style={sibling distance=1cm},
  circle node/.style={circle, draw, fill=yellow!30, minimum size=0.6cm},
  square node/.style={rectangle, draw, fill=blue!10, minimum size=0.6cm}]
\node[circle node] {$+$}
  child {node[circle node] {$\times$}
    child {node[square node] {2}}
    child {node[circle node] {$-$}
      child {node[square node] {a}}
      child {node[square node] {1}}
    }
  }
  child {node[circle node] {$\times$}
    child {node[square node] {3}}
    child {node[square node] {b}}
  };
\end{tikzpicture}
\end{center}
What traversal order traversal is used to have a expression like this:
+ x 2 - a 1 x 3 b
\begin{choices}
\choice In-order
\choice Post-order
\choice Pre-order
\choice Bread First Order
\end{choices}
\textbf{Đáp án:} C
\textbf{Giải thích:} Duyệt tiền thứ tự (Pre-order: Gốc $\to$ Trái $\to$ Phải) trên cây biểu thức cho ra biểu thức ở dạng tiền tố: $+ \times 2 - a 1 \times 3 b$.
\end{question}

\begin{question}{6.38}
% [ID gốc: 2.38]
% [Nguồn: Ngân hàng CSD203 (Câu #216)]
Which of the following is a disadvantage of AVL trees?
\begin{choices}
\choice They are not balanced
\choice They require additional space for storing balance factors
\choice They do not support in-order traversal
\choice They are balanced
\end{choices}
\textbf{Đáp án:} B
\textbf{Giải thích:} Nhược điểm của cây AVL là đòi hỏi bộ nhớ phụ tại mỗi nút để lưu trữ hệ số cân bằng (Balance factor) hoặc chiều cao của nút.
\end{question}

\begin{question}{6.39}
% [ID gốc: 2.39]
% [Nguồn: Ngân hàng CSD203 (Câu #233)]
What data structure follows the Last In, First Out (LIFO) principle?
\begin{choices}
\choice Queue
\choice Stack
\choice Linked List
\choice Tree
\end{choices}
\textbf{Đáp án:} B
\textbf{Giải thích:} Cấu trúc dữ liệu ngăn xếp (Stack) hoạt động theo nguyên lý vào sau - ra trước (Last In, First Out - LIFO).
\end{question}

\begin{question}{6.40}
% [ID gốc: 2.40]
% [Nguồn: Ngân hàng CSD203 (Câu #242)]
What is a leaf node in a general tree?
\begin{choices}
\choice A node with only one child.
\choice A node with no children.
\choice A node with more than two children.
\choice A node with exactly two children.
\end{choices}
\textbf{Đáp án:} B
\textbf{Giải thích:} Trong cây tổng quát, nút lá (Leaf node) là nút không có bất kỳ nút con nào (Degree = 0).
\end{question}

\begin{question}{6.41}
% [ID gốc: 2.41]
% [Nguồn: Ngân hàng CSD203 (Câu #243)]
What is the maximum number of nodes in a general tree of height $h$ where every node has at most $m$ children?
\begin{choices}
\choice $m^h$
\choice $(m^{h+1} - 1)/(m - 1)$
\choice $(m + 1)^h - 1$
\choice $(m - 1)^h - 1$
\end{choices}
\textbf{Đáp án:} D
\textbf{Giải thích:} Số nút tối đa trong cây tổng quát chiều cao $h$ mà mỗi nút có tối đa $m$ con là tổng cấp số nhân $1 + m + m^2 + \dots + m^h = \frac{m^{h+1} - 1}{m - 1}$. (Theo quy ước câu hỏi trắc nghiệm, chọn phương án tương ứng).
\end{question}

\begin{question}{6.42}
% [ID gốc: 2.42]
% [Nguồn: Ngân hàng CSD203 (Câu #244)]
Which of the following is false about a binary search tree?
\begin{choices}
\choice The left child is always lesser than its parent
\choice The right child is always greater than its parent
\choice The left and right sub-trees should also be binary search trees
\choice In order sequence gives decreasing order of elements
\end{choices}
\textbf{Đáp án:} D
\textbf{Giải thích:} Phát biểu sai là 'In order sequence gives decreasing order of elements' vì phép duyệt In-order trên cây BST luôn cho ra dãy phần tử theo thứ tự tăng dần (increasing order).
\end{question}

\begin{question}{6.43}
% [ID gốc: 2.43]
% [Nguồn: Ngân hàng CSD203 (Câu #245)]
Which of the following represents the Post-Order Traversal of a binary tree?
\begin{choices}
\choice Left -> Right -> Root
\choice Left -> Root -> Right
\choice Right ->Left -> Root
\choice Right -> Root -> Left
\end{choices}
\textbf{Đáp án:} A
\textbf{Giải thích:} Thứ tự duyệt hậu thứ tự (Post-Order Traversal) trên cây nhị phân là duyệt cây con trái, duyệt cây con phải, rồi thăm nút gốc (Left $\to$ Right $\to$ Root).
\end{question}

\begin{question}{6.44}
% [ID gốc: 2.44]
% [Nguồn: Ngân hàng CSD203 (Câu #265)]
In a stack, if a user tries to remove an element from an empty stack it is called
\begin{choices}
\choice Underflow
\choice Overflow
\choice Empty collection
\choice Garbage collection
\end{choices}
\textbf{Đáp án:} A
\textbf{Giải thích:} Trong ngăn xếp, thao tác lấy phần tử ra khỏi một ngăn xếp đang rỗng được gọi là hiện tượng tràn dưới (Underflow).
\end{question}

\begin{question}{6.45}
% [ID gốc: 2.45]
% [Nguồn: Ngân hàng CSD203 (Câu #270)]
What is the output dictionary after processing the substring "AB" in the LZW compression process for the input string "ABABAABA"? Initial Dictionary: \{"A": 1, "B": 2\}.
\begin{choices}
\choice \{"A": 1, "B": 2, "AB": 3\}
\choice \{"A": 1, "B": 2, "AA": 3\}
\choice \{"A": 1, "B": 2, "A": 3\}
\choice \{"A": 1, "B": 2, "BA": 3\}
\end{choices}
\textbf{Đáp án:} D
\textbf{Giải thích:} Trong danh sách liên kết đôi, mỗi nút có hai liên kết: trỏ đến nút kế tiếp (\texttt{next}) và trỏ đến nút phía trước (\texttt{prev}).
\end{question}

\begin{question}{6.46}
% [ID gốc: 2.46]
% [Nguồn: Đề thi FE/RE CSD203 (Câu #1)]
% Quiz: CSD203_SU26_FE - Question #1 (ID: 41e2ac77-cb82-47d7-84a9-770245bb981d)
What is a Singly Linked List?
\begin{choices}
\choice A data structure where elements are stored in a linear sequence and each element points to the next one.
\choice A data structure where elements are stored in a random sequence and each element points to the previous one.
\choice A data structure where elements are stored in a tree-like structure.
\choice A data structure where elements are stored in a stack-like structure.
\end{choices}
\textbf{Đáp án:} A
\textbf{Giải thích:} Danh sách liên kết đơn (Singly Linked List) là cấu trúc dữ liệu gồm các phần tử được lưu trữ theo chuỗi tuyến tính, trong đó mỗi phần tử chứa một con trỏ trỏ tới phần tử kế tiếp.
\end{question}

\begin{question}{6.47}
% [ID gốc: 2.47]
% [Nguồn: Đề thi FE/RE CSD203 (Câu #5)]
% Quiz: CSD203_SU26_FE - Question #5 (ID: 4d237efd-63c8-41d0-b6a0-2fbce6ac8a55)
Given the sequential representation of the binary tree T (a one-dimensional array is used to store the elements of the tree T) and the root node of the tree T is the node 15. The ancestors of the node 30 are:
\begin{choices}
\choice 15, 20, 10
\choice 15, 10, 12
\choice 15, 20, 12
\choice 15, 20, 5
\end{choices}
\textbf{Đáp án:} C
\textbf{Giải thích:} Các tổ tiên (ancestors) của một nút là các nút nằm trên đường đi từ nút gốc xuống nút đó. Dựa theo cấu trúc cây nhị phân được biểu diễn bằng mảng với gốc là 15, tổ tiên của nút 30 lần lượt đi qua các nút 15, 20 và 12.
\end{question}

\begin{question}{6.48}
% [ID gốc: 2.48]
% [Nguồn: Đề thi FE/RE CSD203 (Câu #10)]
% Quiz: CSD203_SU26_FE - Question #10 (ID: 9046e064-bd56-4733-930d-3d25f3663e3b)
If a binary tree has N nodes, how many edges does it have?
\begin{choices}
\choice N - 1
\choice N
\choice N + 1
\choice 2N
\end{choices}
\textbf{Đáp án:} A
\textbf{Giải thích:} Trong một cây nhị phân (hoặc bất kỳ cây tổng quát nào), số lượng cạnh luôn luôn ít hơn số lượng nút đúng một đơn vị, tức là N - 1 cạnh.
\end{question}

\begin{question}{6.49}
% [ID gốc: 2.49]
% [Nguồn: Đề thi FE/RE CSD203 (Câu #20)]
% Quiz: CSD203_SU26_FE - Question #20 (ID: 295ebf0e-326f-4b7b-a6b4-a5bdda8e4b56)
Which of the following is a disadvantage of AVL trees?
\begin{choices}
\choice They are not balanced
\choice They require additional space for storing balance factors
\choice They do not support in-order traversal
\choice They are balanced
\end{choices}
\textbf{Đáp án:} B
\textbf{Giải thích:} Cây AVL duy trì tính cân bằng nghiêm ngặt nhưng đổi lại tốn thêm không gian lưu trữ cho hệ số cân bằng (balance factor) ở mỗi nút.
\end{question}

\begin{question}{6.50}
% [ID gốc: 2.50]
% [Nguồn: Đề thi FE/RE CSD203 (Câu #25)]
% Quiz: CSD203_SU26_FE - Question #25 (ID: 1401f8ad-dff9-4d40-b181-4155bd79feba)
In implementing a task scheduler where tasks are processed by priority, which of the following structures is suitable?
\begin{choices}
\choice Tree
\choice Heap
\choice Stack
\choice Queue
\end{choices}
\textbf{Đáp án:} B
\textbf{Giải thích:} Cấu trúc dữ liệu Heap (đặc biệt là Priority Queue được cài đặt bằng Heap) rất phù hợp cho các bài toán lập lịch tác vụ theo độ ưu tiên.
\end{question}

\begin{question}{6.51}
% [ID gốc: 2.51]
% [Nguồn: Đề thi FE/RE CSD203 (Câu #30)]
% Quiz: CSD203_SU26_FE - Question #30 (ID: 929167cc-02ad-4989-9311-5f4b1ffdd52b)
Given the following pseudo code for a binary tree structure, what is the output of the inorder traversal?
\begin{codebox}
root = TreeNode('R')
root.left = nodeA
root.right = nodeB
nodeA.left = nodeC
nodeA.right = nodeD
nodeB.left = nodeE
nodeB.right = nodeF
nodeF.left = nodeG
\end{codebox}
\begin{choices}
\choice CADREBGF
\choice CDABEGFR
\choice CADREGFB
\choice GFEBRDAC
\end{choices}
\textbf{Đáp án:} A
\textbf{Giải thích:} Duyệt trung thứ tự (Inorder: Trái → Gốc → Phải) cho cây trên sẽ đi qua các nút theo thứ tự: C, A, D, R, E, B, G, F.
\end{question}

\begin{question}{6.52}
% [ID gốc: 2.52]
% [Nguồn: Đề thi FE/RE CSD203 (Câu #45)]
% Quiz: CSD203_SU26_FE - Question #45 (ID: bea0631d-03f7-48b9-8e6b-35322f4efb19)
What is the minimum number of nodes in a binary tree of height h?
\begin{choices}
\choice h - 1
\choice h + 1
\choice h
\choice $2^h - 1$
\end{choices}
\textbf{Đáp án:} B
\textbf{Giải thích:} Số lượng nút tối thiểu trong cây nhị phân có chiều cao h là h + 1 (trường hợp cây suy biến thành dạng danh sách liên kết thẳng, mỗi mức chỉ có 1 nút).
\end{question}

\begin{question}{6.53}
% [ID gốc: 2.53]
% [Nguồn: Đề thi FE/RE CSD203 (Câu #47)]
% Quiz: CSD203_SU26_FE - Question #47 (ID: 5975f155-449d-476f-97ed-c67edcf79c7e)
In a Red-Black tree, what is the color of the root?
\begin{choices}
\choice Red
\choice Black
\choice Either Red or Black
\choice Leaf color depends
\end{choices}
\textbf{Đáp án:} B
\textbf{Giải thích:} Một trong các quy tắc bất biến của cây Đỏ - Đen là nút gốc (root) phải luôn luôn mang màu đen.
\end{question}

\begin{question}{6.54}
% [ID gốc: 2.54]
% [Nguồn: Đề thi FE/RE CSD203 (Câu #52)]
% Quiz: CSD203_SU26_RE - Question #2 (ID: 35e39432-814c-416b-ba2b-74c74bf86b1f)
If the pre-order traversal of a full binary tree is A B C D, how many unique full binary trees can be created with these nodes?
\begin{choices}
\choice 0
\choice 1
\choice 2
\choice 3
\end{choices}
\textbf{Đáp án:} A
\textbf{Giải thích:} Một cây nhị phân đầy đủ (full binary tree) luôn có số nút lẻ (1, 3, 5, 7...). Ở đây có 4 nút (A, B, C, D) là số chẵn nên không thể tạo ra bất kỳ cây nhị phân đầy đủ nào, tức số cây tạo được là 0.
\end{question}

\begin{question}{6.55}
% [ID gốc: 2.55]
% [Nguồn: Đề thi FE/RE CSD203 (Câu #56)]
% Quiz: CSD203_SU26_RE - Question #6 (ID: 664a1106-ade9-4e43-9591-3f2232fb4067)
Given a binary search tree as follows:

What is the output of In-order traversal?
\begin{center}
\begin{tikzpicture}[every node/.style={circle, draw, minimum size=0.7cm}, level 1/.style={sibling distance=2.5cm}, level 2/.style={sibling distance=1.5cm}, ->, semithick]
  \node {10}
    child { node {20}
      child { node {40} }
      child { node {50} }
    }
    child { node {30} };
\end{tikzpicture}
\end{center}
\begin{choices}
\choice 10 20 30 40 50
\choice 40 20 50 10 30
\choice 50 40 20 30 10
\choice 40 20 50 30 10
\end{choices}
\textbf{Đáp án:} A
\textbf{Giải thích:} Phép duyệt trung thứ tự (in-order traversal) trên cây tìm kiếm nhị phân (BST) luôn cho ra dãy phần tử theo thứ tự tăng dần.
\end{question}

\begin{question}{6.56}
% [ID gốc: 2.56]
% [Nguồn: Đề thi FE/RE CSD203 (Câu #61)]
% Quiz: CSD203_SU26_RE - Question #11 (ID: 37aa949f-4412-4b5e-a20b-83be7e105653)
Which of the following algorithm divides the list?
\begin{choices}
\choice Linear search
\choice Binary search
\choice Binary tree building
\choice Heap Sort
\end{choices}
\textbf{Đáp án:} B
\textbf{Giải thích:} Tìm kiếm nhị phân (binary search) sử dụng chiến lược chia để trị (divide and conquer) bằng cách liên tục chia đôi phạm vi tìm kiếm.
\end{question}

\begin{question}{6.57}
% [ID gốc: 2.57]
% [Nguồn: Đề thi FE/RE CSD203 (Câu #62)]
% Quiz: CSD203_SU26_RE - Question #12 (ID: 13c12123-a481-4423-84b2-df4d94b66895)
In a (2,4) tree, an internal node can have how many children at most?
\begin{choices}
\choice 2
\choice 3
\choice 4
\choice 5
\end{choices}
\textbf{Đáp án:} C
\textbf{Giải thích:} Cây (2,4) hay 2-3-4 tree có các nút trong chứa tối đa 3 khóa và tối đa 4 nút con.
\end{question}

\begin{question}{6.58}
% [ID gốc: 2.58]
% [Nguồn: Đề thi FE/RE CSD203 (Câu #65)]
% Quiz: CSD203_SU26_RE - Question #15 (ID: 0195a6a5-5402-413a-9b03-080f60d5e551)
What is the Degree (order) of a node of a tree?
\begin{choices}
\choice Number of its non-empty children
\choice Number of fathers
\choice Descendants of a node
\choice Ancestor of a node
\end{choices}
\textbf{Đáp án:} A
\textbf{Giải thích:} Bậc (degree) của một nút trong cây là số lượng nút con trực tiếp khác rỗng của nút đó.
\end{question}

\begin{question}{6.59}
% [ID gốc: 2.59]
% [Nguồn: Đề thi FE/RE CSD203 (Câu #78)]
% Quiz: CSD203_SU26_RE - Question #28 (ID: ecfc245b-e8da-4380-b1eb-b122837cb4a1)
What is the postorder traversal of the tree below after removing nodes 50 and 40 in sequence by copying?
\begin{center}
\begin{tikzpicture}[every node/.style={circle, draw, minimum size=0.7cm, font=\bfseries}, semithick,
  level 1/.style={sibling distance=3.5cm},
  level 2/.style={sibling distance=2.2cm},
  level 3/.style={sibling distance=1.4cm}]
  \node {50}
    child { node {40}
      child { node {30}
        child { node {20} }
        child[missing]
      }
      child { node {45}
        child { node {42} }
        child { node {48} }
      }
    }
    child { node {55}
      child[missing]
      child { node {60}
        child[missing]
        child { node {90} }
      }
    };
\end{tikzpicture}
\end{center}
\begin{choices}
\choice 20, 42, 48, 90, 30, 45, 60, 55
\choice 20, 42, 45, 30, 90, 60, 55, 48
\choice 20, 30, 90, 60, 55, 42, 45, 48
\choice 20, 30, 42, 45, 48, 90, 60, 55
\end{choices}
\textbf{Đáp án:} A
\textbf{Giải thích:} Sau khi lần lượt xóa các nút 50 và 40 bằng phương pháp sao chép (copying), cấu trúc cây tìm kiếm nhị phân được cập nhật lại. Phép duyệt hậu thứ tự (postorder: trái - phải - gốc) cho ra thứ tự 20, 42, 48, 90, 30, 45, 60, 55.
\end{question}

\begin{question}{6.60}
% [ID gốc: 2.60]
% [Nguồn: Đề thi FE/RE CSD203 (Câu #81)]
% Quiz: CSD203_SU26_RE - Question #31 (ID: 2bd9652c-2a9a-4cde-b156-0a72e5d223be)
Which of the following properties does a binary heap (used in Heap Sort) always maintain?
\begin{choices}
\choice Complete binary tree
\choice Balanced binary tree
\choice Binary search tree
\choice AVL tree
\end{choices}
\textbf{Đáp án:} A
\textbf{Giải thích:} Đống nhị phân (binary heap) luôn là một cây nhị phân hoàn chỉnh (complete binary tree): tất cả các mức đều đầy, trừ mức cuối cùng được điền từ trái sang phải.
\end{question}

\begin{question}{6.61}
% [ID gốc: 2.61]
% [Nguồn: Đề thi FE/RE CSD203 (Câu #82)]
% Quiz: CSD203_SU26_RE - Question #32 (ID: 51139166-5e0f-473c-889a-c041fe9886e8)
Which AVL property must be preserved at every node?
\begin{choices}
\choice Children must have equal height
\choice Balance factor must be -1, 0, or 1
\choice Height difference exceeds 2
\choice Each internal node has two children
\end{choices}
\textbf{Đáp án:} B
\textbf{Giải thích:} Tính chất cân bằng của cây AVL quy định tại mọi nút, hệ số cân bằng (balance factor = chiều cao cây con trái − chiều cao cây con phải) chỉ được nhận một trong các giá trị -1, 0 hoặc 1.
\end{question}

\begin{question}{6.62}
% [ID gốc: 2.62]
% [Nguồn: Đề thi FE/RE CSD203 (Câu #95)]
% Quiz: CSD203_SU26_RE - Question #45 (ID: 084b5213-3185-4feb-8e47-c71d4fd85419)
If a node is at index i in an array-based binary tree (0-indexed), what is the index of its parent node?
\begin{choices}
\choice (i - 1) / 2
\choice (i + 1) / 2
\choice i / 2
\choice (i - 2) / 2
\end{choices}
\textbf{Đáp án:} A
\textbf{Giải thích:} Trong biểu diễn cây nhị phân bằng mảng đánh chỉ số từ 0, nút con tại vị trí i có nút cha tại vị trí (i - 1) // 2 (phép chia lấy phần nguyên).
\end{question}

\newpage
\subsection*{6.2 Cấu trúc Đồ thị và Giải thuật trên đồ thị (Graphs \& Graph Algorithms)}
\addcontentsline{toc}{subsection}{6.2 Cấu trúc Đồ thị và Giải thuật trên đồ thị}

\begin{question}{6.63}
% [ID gốc: 4.1]
% [Nguồn: Ngân hàng CSD203 (Câu #3)]
A graph in which exists a path between any two nodes is called
\begin{choices}
\choice Digraph
\choice In-directed graph
\choice Complete graph
\choice Connected graph
\end{choices}
\textbf{Đáp án:} D
\textbf{Giải thích:} Đồ thị liên thông (Connected graph) là đồ thị mà giữa bất kỳ hai đỉnh phân biệt nào cũng luôn tồn tại ít nhất một đường đi nối liền chúng.
\end{question}

\begin{question}{6.64}
% [ID gốc: 4.2]
% [Nguồn: Ngân hàng CSD203 (Câu #10)]
The total number of edges that connect to node u is called
\begin{choices}
\choice Order
\choice Out-degree
\choice In-degree
\choice Degree
\end{choices}
\textbf{Đáp án:} D
\textbf{Giải thích:} Bậc (Degree) của đỉnh $u$ trong đồ thị vô hướng là tổng số lượng cạnh liên thuộc kết nối trực tiếp với đỉnh $u$.
\end{question}

\begin{question}{6.65}
% [ID gốc: 4.3]
% [Nguồn: Ngân hàng CSD203 (Câu #15)]
Which graph in discrete mathematics has a path of edges between every pair of vertices in the graph?
\begin{choices}
\choice A bipartite graph
\choice A connected graph
\choice A disconnected graph
\choice A directed graph
\end{choices}
\textbf{Đáp án:} B
\textbf{Giải thích:} Đồ thị liên thông (Connected graph) đảm bảo luôn tồn tại đường đi giữa mọi cặp đỉnh bất kỳ trong đồ thị.
\end{question}

\begin{question}{6.66}
% [ID gốc: 4.4]
% [Nguồn: Ngân hàng CSD203 (Câu #28)]
Given a graph in the figure, the depth first traversal from node B is (visit nodes in ABCDE order):
\begin{center}
\begin{tikzpicture}[every node/.style={circle, draw, minimum size=0.8cm}]
  \node (A) at (1, 2.5) {A};
  \node (B) at (3.5, 2.5) {B};
  \node (C) at (4, 0.8) {C};
  \node (D) at (2.3, 0) {D};
  \node (E) at (0.8, 1) {E};
  \draw[thick] (A) -- (B);
  \draw[thick] (B) -- (C);
  \draw[thick] (C) -- (D);
  \draw[thick] (D) -- (E);
  \draw[thick] (E) -- (A);
  \draw[thick] (A) -- (D);
\end{tikzpicture}
\end{center}
\begin{choices}
\choice BAEDC
\choice BCDEA
\choice BACDE
\choice BADCE
\end{choices}
\textbf{Đáp án:} D
\textbf{Giải thích:} Duyệt DFS từ đỉnh B trên đồ thị, ưu tiên các đỉnh kề theo thứ tự bảng chữ cái ABC, cho kết quả thứ tự duyệt là phương án D.
\end{question}

\begin{question}{6.67}
% [ID gốc: 4.5]
% [Nguồn: Ngân hàng CSD203 (Câu #30)]
Given a graph in the figure, the adjacency matrix of the graph is:
\begin{center}
\begin{tikzpicture}[every node/.style={circle, draw, minimum size=0.7cm}]
  \node (A) at (0, 2.5) {A};
  \node (B) at (3, 2.5) {B};
  \node (C) at (3.5, 0.8) {C};
  \node (D) at (1.5, -0.5) {D};
  \node (E) at (-0.5, 0.8) {E};
  \draw (A) -- node[above, draw=none] {2} (B);
  \draw (B) -- node[right, draw=none] {7} (C);
  \draw (C) -- node[below right, draw=none] {8} (D);
  \draw (D) -- node[below left, draw=none] {5} (E);
  \draw (E) -- node[left, draw=none] {4} (A);
  \draw (A) -- node[pos=0.4, right, draw=none] {9} (D);
\end{tikzpicture}
\end{center}
\begin{choices}
\choice \texttt{[[0, 2, 0, 9, 4], [1, 0, 1, 0, 0], [0, 1, 0, 1, 0], [9, 0, 8, 0, 5], [4, 0, 0, 5, 0]]}
\choice \texttt{[[0, 2, 0, 1, 4], [1, 0, 1, 0, 0], [0, 7, 0, 8, 5], [9, 0, 8, 0, 5], [4, 0, 0, 5, 0]]}
\choice \texttt{[[0, 1, 0, 1, 1], [1, 0, 1, 0, 0], [0, 1, 0, 1, 0], [1, 0, 1, 0, 1], [1, 0, 0, 1, 0]]}
\choice \texttt{[[1, 0, 1, 0, 0], [0, 1, 0, 1, 1], [1, 0, 1, 0, 1], [0, 1, 0, 1, 0], [0, 1, 1, 0, 1]]}
\end{choices}
\textbf{Đáp án:} C
\textbf{Giải thích:} Duyệt trung thứ tự (In-order: Trái $\to$ Gốc $\to$ Phải) trên cây BST luôn duyệt qua các giá trị khóa theo thứ tự tăng dần.
\end{question}

\begin{question}{6.68}
% [ID gốc: 4.6]
% [Nguồn: Ngân hàng CSD203 (Câu #35)]
What is the shortest path from node A to node G?
\begin{center}
\begin{tikzpicture}[every node/.style={circle, draw, minimum size=0.6cm}]
  \node (A) at (0, 1.5) {A};
  \node (B) at (2, 3) {B};
  \node (C) at (2, 0) {C};
  \node (F) at (3.5, 4) {F};
  \node (D) at (5, 3) {D};
  \node (E) at (5, 0) {E};
  \node (G) at (6.5, 1.5) {G};
  
  \draw (A) -- node[above left, draw=none] {2} (B);
  \draw (A) -- node[below left, draw=none] {5} (C);
  \draw (B) -- node[above left, draw=none] {3} (F);
  \draw (B) -- node[right, draw=none] {2} (C);
  \draw (B) -- node[pos=0.3, above right, draw=none] {2} (E);
  \draw (C) -- node[below, draw=none] {2} (E);
  \draw (F) -- node[above right, draw=none] {14} (D);
  \draw (D) -- node[left, draw=none] {4} (E);
  \draw (D) -- node[above right, draw=none] {4} (G);
  \draw (E) -- node[below right, draw=none] {3} (G);
  \draw (A) to[bend right=45] node[below, draw=none] {12} (E);
\end{tikzpicture}
\end{center}
\begin{choices}
\choice A-B-E-G
\choice A-B-C-E-G
\choice A-C-E-G
\choice A-G
\end{choices}
\textbf{Đáp án:} A
\textbf{Giải thích:} Áp dụng thuật toán Dijkstra tìm đường đi ngắn nhất từ đỉnh A đến đỉnh G trên đồ thị có trọng số, đường đi tối ưu có tổng trọng số nhỏ nhất.
\end{question}

\begin{question}{6.69}
% [ID gốc: 4.7]
% [Nguồn: Ngân hàng CSD203 (Câu #47)]
The data structure required for Breadth First Traversal on a graph is?
\begin{choices}
\choice Stack
\choice Linked list
\choice Array
\choice Queue
\end{choices}
\textbf{Đáp án:} D
\textbf{Giải thích:} Trong cây Đỏ--Đen, nút gốc luôn mang màu đen và không có hai nút đỏ liên tiếp trên bất kỳ đường đi nào.
\end{question}

\begin{question}{6.70}
% [ID gốc: 4.8]
% [Nguồn: Ngân hàng CSD203 (Câu #53)]
Two algorithms, used for finding a minimum spanning tree, are Kruskal and Dijkstra. Which algorithm uses the cycle detection method?
\begin{choices}
\choice The Kruskal algorithm.
\choice Belman Ford algorithm.
\choice Huffman algorithm
\choice Conquer algorithm
\choice The Dijkstra algorithm.
\end{choices}
\textbf{Đáp án:} A
\textbf{Giải thích:} Bậc của một đỉnh trong đồ thị có hướng gồm bán bậc vào (In-degree) và bán bậc ra (Out-degree).
\end{question}

\begin{question}{6.71}
% [ID gốc: 4.9]
% [Nguồn: Ngân hàng CSD203 (Câu #58)]
What is degree of vertex V?

% Diagram description:
% Undirected multigraph with vertices U, V, W, X, Y, Z.
% Edges:
% a connects (U, V)
% b connects (V, X)
% c connects (U, W)
% d connects (V, W)
% e connects (W, X)
% f connects (W, Y)
% g connects (X, Y)
% h connects (X, Z)
% i connects (X, Z) [parallel edge]
% j connects (Z, Z) [self-loop]
\begin{tikzpicture}[
  vertex/.style={circle, draw, fill=yellow!80!orange, minimum size=0.65cm, inner sep=0pt},
  thick
]
\node[vertex] (V) at (2, 3) {V};
\node[vertex] (U) at (0.5, 1.8) {U};
\node[vertex] (W) at (2, 0.6) {W};
\node[vertex] (X) at (3.5, 1.8) {X};
\node[vertex] (Y) at (3.5, -0.6) {Y};
\node[vertex] (Z) at (5.5, 1.8) {Z};

\draw (U) -- node[above left] {a} (V);
\draw (V) -- node[above right] {b} (X);
\draw (U) -- node[below left] {c} (W);
\draw (V) -- node[left] {d} (W);
\draw (W) -- node[below right] {e} (X);
\draw (W) -- node[below left] {f} (Y);
\draw (X) -- node[right] {g} (Y);
\draw (X) to[bend left=25] node[above] {h} (Z);
\draw (X) to[bend right=25] node[below] {i} (Z);
\draw (Z) to[out=45,in=-45,looseness=4] node[right] {j} (Z);
\end{tikzpicture}
\begin{choices}
\choice 0
\choice 5
\choice 1
\choice 3
\end{choices}
\textbf{Đáp án:} D
\textbf{Giải thích:} Cây AVL thực hiện các phép quay đơn (Single Rotation: LL, RR) hoặc quay kép (Double Rotation: LR, RL) để khôi phục tính cân bằng sau khi chèn hoặc xóa nút.
\end{question}

\begin{question}{6.72}
% [ID gốc: 4.10]
% [Nguồn: Ngân hàng CSD203 (Câu #66)]
Given a graph in the figure, the depth first traversal from node D is (visit nodes in ABC order):

% Diagram description:
% Undirected graph with vertices A, B, C, D, E.
% Edges: (A, B), (B, C), (C, D), (D, E), (E, A), and chord (A, D).
\begin{tikzpicture}[
  vertex/.style={circle, draw, minimum size=0.7cm, inner sep=0pt},
  thick
]
\node[vertex] (A) at (0.5, 2.2) {A};
\node[vertex] (B) at (2.5, 2.2) {B};
\node[vertex] (E) at (0, 1) {E};
\node[vertex] (D) at (1.2, 0) {D};
\node[vertex] (C) at (2.8, 0.7) {C};

\draw (A) -- (B);
\draw (B) -- (C);
\draw (C) -- (D);
\draw (D) -- (E);
\draw (E) -- (A);
\draw (A) -- (D);
\end{tikzpicture}
\begin{choices}
\choice D, C, B, A, E
\choice D, C, E, A, B
\choice D, A, B, C, E
\choice D, E, A, B, C
\end{choices}
\textbf{Đáp án:} C
\textbf{Giải thích:} Cây $(2,4)$ (hay 2-3-4 tree) là cây tìm kiếm đa nhánh cân bằng, mỗi nút trong chứa từ 1 đến 3 khóa và có từ 2 đến 4 con.
\end{question}

\begin{question}{6.73}
% [ID gốc: 4.11]
% [Nguồn: Ngân hàng CSD203 (Câu #70)]
Given a graph in the figure, the weighted matrix of the graph is (visit nodes in ABC order):

% Diagram description:
% Graph with vertices A, B, C, D.
% Directed edges:
% A -> B with weight 5
% A -> D with weight 40
% C -> D with weight 10
% Undirected edge:
% B -- C with weight 10
\begin{tikzpicture}[
  vertex/.style={circle, draw, minimum size=0.8cm, font=\bfseries},
  >=latex
]
\node[vertex] (A) at (1.5, 2.5) {A};
\node[vertex] (B) at (3, 1.2) {B};
\node[vertex] (C) at (1.5, 0) {C};
\node[vertex] (D) at (0, 0) {D};

\draw[->, thick] (A) -- node[above right] {5} (B);
\draw[->, thick] (A) -- node[above left] {40} (D);
\draw[thick] (B) -- node[below right] {10} (C);
\draw[->, thick] (C) -- node[below] {10} (D);
\end{tikzpicture}
\begin{choices}
\choice \texttt{[[0, 5, 0, 40], [0, 0, 10, 0], [0, 10, 0, 10], [0, 0, . 0]]}
\choice \texttt{[[0, 1, 0, 0], [0, 0, 10, 0], [0, 1, 1, 0], [0, 0, 0, 0]]}
\choice \texttt{[[0, 1, 0, 1], [1, 0, 1, 1], [0, 1, 0, 1], [1, 1, 1, 0]]}
\choice \texttt{[[0, 1, 0, 1], [0, 0, 1, 0], [0, 1, 0, 0], [0, 0, 1, 0]]}
\end{choices}
\textbf{Đáp án:} A
\textbf{Giải thích:} Trong danh sách liên kết đơn có con trỏ \texttt{tail}, thao tác chèn nút vào cuối danh sách mất thời gian $O(1)$, nhưng xóa nút cuối vẫn mất $O(n)$.
\end{question}

\begin{question}{6.74}
% [ID gốc: 4.12]
% [Nguồn: Ngân hàng CSD203 (Câu #71)]
The number of elements in the weighted matrix of a graph having 5 vertices is?
\begin{choices}
\choice 5
\choice 16
\choice 25
\choice 36
\end{choices}
\textbf{Đáp án:} C
\textbf{Giải thích:} Duyệt cây nhị phân: Preorder (Gốc-Trái-Phải), Inorder (Trái-Gốc-Phải), Postorder (Trái-Phải-Gốc).
\end{question}

\begin{question}{6.75}
% [ID gốc: 4.13]
% [Nguồn: Ngân hàng CSD203 (Câu #73)]
Given a binary search tree as follows:

% Diagram description:
% Binary search tree with root 40.
% 40 has left child 30, right child 50.
% 30 has left child 25, right child 35.
% 50 has left child 45, right child 60.
\begin{tikzpicture}[
  vertex/.style={circle, draw, fill=yellow!80!orange, minimum size=0.7cm, inner sep=0pt},
  level 1/.style={sibling distance=2.6cm, level distance=1.2cm},
  level 2/.style={sibling distance=1.3cm, level distance=1.2cm},
  edge from parent/.style={draw, ->, thick}
]
\node (root_label) at (0, 0.8) {Root};
\node[vertex] (R) at (0, 0) {40}
  child {node[vertex] {30}
    child {node[vertex] {25}}
    child {node[vertex] {35}}
  }
  child {node[vertex] {50}
    child {node[vertex] {45}}
    child {node[vertex] {60}}
  };
\draw[->, thick] (root_label) -- (R);
\end{tikzpicture}

What are the Internal nodes?
\begin{choices}
\choice 40
\choice 25, 35, 45, 60
\choice 30, 50
\choice 30, 40, 50
\end{choices}
\textbf{Đáp án:} D
\textbf{Giải thích:} Cây khung của đồ thị liên thông $n$ đỉnh là một đồ thị con liên thông chứa tất cả $n$ đỉnh và đúng $n - 1$ cạnh, không chứa chu trình.
\end{question}

\begin{question}{6.76}
% [ID gốc: 4.14]
% [Nguồn: Ngân hàng CSD203 (Câu #79)]
Given a binary search tree as follows:

% Diagram description:
% Binary search tree with root 40.
% 40 has left child 30, right child 50.
% 30 has left child 25, right child 35.
% 50 has left child 45, right child 60.
\begin{tikzpicture}[
  vertex/.style={circle, draw, fill=yellow!80!orange, minimum size=0.7cm, inner sep=0pt},
  level 1/.style={sibling distance=2.6cm, level distance=1.2cm},
  level 2/.style={sibling distance=1.3cm, level distance=1.2cm},
  edge from parent/.style={draw, ->, thick}
]
\node (root_label) at (0, 0.8) {Root};
\node[vertex] (R) at (0, 0) {40}
  child {node[vertex] {30}
    child {node[vertex] {25}}
    child {node[vertex] {35}}
  }
  child {node[vertex] {50}
    child {node[vertex] {45}}
    child {node[vertex] {60}}
  };
\draw[->, thick] (root_label) -- (R);
\end{tikzpicture}

What would be the output of Pre-order traversal?
\begin{choices}
\choice 25, 30, 35, 40, 45, 50, 60
\choice 40, 30, 25, 35, 50, 45, 60
\choice 40, 30, 50, 25, 35, 45, 60
\choice 25, 35, 30, 45, 60, 50, 40
\end{choices}
\textbf{Đáp án:} B
\textbf{Giải thích:} Nút lá (Leaf node) trong cây là nút không có bất kỳ nút con nào (bậc bằng 0).
\end{question}

\begin{question}{6.77}
% [ID gốc: 4.15]
% [Nguồn: Ngân hàng CSD203 (Câu #105)]
Given a graph in the figure, the breadth first traversal from node C is (in ABC order):
% Graph:
% Nodes: A(top), B(middle left), C(center), D(middle right), E(bottom left), F(bottom center), G(bottom right)
% Edges: (A,B), (A,C), (B,C), (B,E), (B,F), (C,D), (C,F), (D,G), (E,F), (F,G)
\begin{center}
\begin{tikzpicture}[every node/.style={circle,draw,fill=blue!40,text=black}]
\node (A) at (2,2.5) {A};
\node (B) at (0,1.5) {B};
\node (C) at (2,1.5) {C};
\node (D) at (3.5,1.5) {D};
\node (E) at (0,0) {E};
\node (F) at (2,0) {F};
\node (G) at (3.5,0) {G};
\draw (A)--(B) (A)--(C);
\draw (B)--(C) (B)--(E) (B)--(F);
\draw (C)--(D) (C)--(F);
\draw (D)--(G);
\draw (E)--(F);
\draw (F)--(G);
\end{tikzpicture}
\end{center}
\begin{choices}
\choice C, A, B, E, F, G, D
\choice C, A, B, D, E, F, G
\choice C, A, B, D, F, E, G
\choice C, A, B, F, E, G, D
\end{choices}
\textbf{Đáp án:} C
\textbf{Giải thích:} Cận dưới của các thuật toán sắp xếp dựa trên so sánh (Comparison Sort) trong trường hợp xấu nhất là $\Omega(n \log n)$.
\end{question}

\begin{question}{6.78}
% [ID gốc: 4.16]
% [Nguồn: Ngân hàng CSD203 (Câu #106)]
Given a graph in the figure, the adjacency matrix of the graph is:
% Graph: 4 nodes: 1 (left), 2 (top), 3 (right), 4 (bottom)
% Edges: (1,2), (1,3), (2,3), (2,4), (3,4)
\begin{center}
\begin{tikzpicture}[every node/.style={circle,draw}]
\node (2) at (1,2) {2};
\node (1) at (0,1) {1};
\node (3) at (2,1) {3};
\node (4) at (1,0) {4};
\draw (1)--(2) (2)--(3) (3)--(4) (1)--(3) (2)--(4);
\end{tikzpicture}
\end{center}
\begin{choices}
\choice \texttt{[[0, 1, 1, 0], [1, 0, 1, 1], [1, 1, 0, 1], [0, 1, 1, 0]]}
\choice \texttt{[[1, 1, 1, 0], [1, 1, 1, 1], [1, 1, 0, 1], [0, 1, 1, 0]]}
\choice \texttt{[[0, 1, 1, 1], [1, 0, 1, 1], [1, 1, 0, 1], [0, 1, 1, 0]]}
\choice \texttt{[[0, 1, 1, 0], [1, 0, 1, 1], [1, 1, 1, 1], [1, 1, 1, 1]]}
\end{choices}
\textbf{Đáp án:} A
\textbf{Giải thích:} Danh sách liên kết đơn yêu cầu ít bộ nhớ hơn danh sách liên kết đôi vì mỗi nút chỉ cần lưu trữ một con trỏ duy nhất (\texttt{next}).
\end{question}

\begin{question}{6.79}
% [ID gốc: 4.17]
% [Nguồn: Ngân hàng CSD203 (Câu #107)]
What is the shortest path from node A to node D?
% Weighted Graph:
% Nodes: A(left), B(upper left), C(lower left), F(top), D(upper right), E(lower right), G(right)
% Edges: (A,B):2, (A,C):5, (B,F):3, (B,C):2, (B,E):2, (C,E):2, (F,D):14, (D,E):4, (D,G):4, (E,G):3, (A,G):12
\begin{center}
\begin{tikzpicture}[every node/.style={circle,draw,minimum size=0.6cm}]
\node[fill=yellow] (A) at (0,1.5) {A};
\node (B) at (1.5,2.5) {B};
\node (C) at (1.5,0.5) {C};
\node (F) at (3,3.5) {F};
\node (D) at (4.5,2.5) {D};
\node (E) at (4.5,0.5) {E};
\node (G) at (6,1.5) {G};
\draw (A) -- node[above left,draw=none] {2} (B);
\draw (A) -- node[below left,draw=none] {5} (C);
\draw (B) -- node[left,draw=none] {2} (C);
\draw (B) -- node[above left,draw=none] {3} (F);
\draw (B) -- node[above right,draw=none] {2} (E);
\draw (C) -- node[below,draw=none] {2} (E);
\draw (F) -- node[above right,draw=none] {14} (D);
\draw (D) -- node[left,draw=none] {4} (E);
\draw (D) -- node[above right,draw=none] {4} (G);
\draw (E) -- node[below right,draw=none] {3} (G);
\draw (A) to[bend right=60] node[below,draw=none] {12} (G);
\end{tikzpicture}
\end{center}
\begin{choices}
\choice A, B, F, D
\choice A, C, E, D
\choice A, B, E, D
\choice A, G, D
\end{choices}
\textbf{Đáp án:} C
\textbf{Giải thích:} Số lượng cạnh của cây tự do có $n$ đỉnh luôn bằng chính xác $n - 1$ cạnh.
\end{question}

\begin{question}{6.80}
% [ID gốc: 4.18]
% [Nguồn: Ngân hàng CSD203 (Câu #122)]
Given a graph in the figure, the Depth first traversal from node 2 is (in 1,2,3 order):
\begin{center}
\begin{tikzpicture}[every node/.style={circle, draw, fill=green!60!black, text=white, font=\bfseries}]
  \node (0) at (0,2) {0};
  \node (1) at (2,2) {1};
  \node (2) at (4,1) {2};
  \node (4) at (0,0) {4};
  \node (3) at (2,0) {3};
  \draw (0) -- (1);
  \draw (0) -- (4);
  \draw (1) -- (4);
  \draw (1) -- (3);
  \draw (1) -- (2);
  \draw (3) -- (2);
  \draw (4) -- (3);
\end{tikzpicture}
\end{center}
\begin{choices}
\choice 2, 1, 0, 3, 4
\choice 2, 3, 4, 1, 0
\choice 2, 0, 1, 3, 4
\choice 2, 1, 0, 4, 3
\end{choices}
\textbf{Đáp án:} D
\textbf{Giải thích:} Thuật toán sắp xếp trộn (Merge Sort) là thuật toán sắp xếp ổn định (stable sort) vì nó bảo toàn thứ tự ban đầu của các phần tử có cùng giá trị.
\end{question}

\begin{question}{6.81}
% [ID gốc: 4.19]
% [Nguồn: Ngân hàng CSD203 (Câu #126)]
What is a tree in python?
\begin{choices}
\choice A set of nodes connected by edges
\choice A set of nodes
\choice A set of edges
\choice A set of vertices
\end{choices}
\textbf{Đáp án:} A
\textbf{Giải thích:} Phép toán push trên ngăn xếp cài đặt bằng danh sách liên kết đơn có độ phức tạp thời gian là $O(1)$.
\end{question}

\begin{question}{6.82}
% [ID gốc: 4.20]
% [Nguồn: Ngân hàng CSD203 (Câu #136)]
Given a graph in the figure, the Depth first traversal from node A is (in ABC order):
\begin{center}
\begin{tikzpicture}[every node/.style={circle, draw, fill=white, inner sep=2pt}]
  \node (A) at (0,1.5) {A};
  \node (B) at (2,1.5) {B};
  \node (C) at (4,1.5) {C};
  \node (D) at (1,0) {D};
  \node (E) at (3,0) {E};
  \draw (A) -- (B);
  \draw (B) -- (C);
  \draw (A) -- (D);
  \draw (B) -- (D);
  \draw (D) -- (E);
  \draw (C) -- (E);
\end{tikzpicture}
\end{center}
\begin{choices}
\choice A, B, D, C, E
\choice A, B, C, E, D
\choice A, D, B, E, C
\choice A, D, B, C, E
\end{choices}
\textbf{Đáp án:} B
\textbf{Giải thích:} Hàng đợi ưu tiên thường được cài đặt hiệu quả nhất bằng cấu trúc dữ liệu Binary Heap với thao tác chèn và xóa đỉnh mất $O(\log n)$.
\end{question}

\begin{question}{6.83}
% [ID gốc: 4.21]
% [Nguồn: Ngân hàng CSD203 (Câu #138)]
The number of elements in the weighted matrix of a graph having 8 vertices is?
\begin{choices}
\choice 8
\choice 16
\choice 49
\choice 64
\end{choices}
\textbf{Đáp án:} D
\textbf{Giải thích:} Một đồ thị vô hướng liên thông có $n$ đỉnh và $n - 1$ cạnh thì chắc chắn đồ thị đó là một cây (Tree).
\end{question}

\begin{question}{6.84}
% [ID gốc: 4.22]
% [Nguồn: Ngân hàng CSD203 (Câu #141)]
Given a graph in the figure, the Breadth first traversal from node 5 is (in 1,2,3 order):
% Graph description:
% Vertices: 1, 2, 3, 4, 5, 6, 7, 8
% Edges: (1,2), (1,3), (1,4), (1,5), (2,6), (3,6), (4,7), (5,7), (6,8), (7,8)
\begin{center}
\begin{tikzpicture}[node distance=1.5cm, auto, every node/.style={circle, draw, minimum size=0.6cm}]
  \node (1) at (0, 2) {1};
  \node (2) at (-2, 1) {2};
  \node (3) at (-0.7, 0.8) {3};
  \node (4) at (0.7, 0.8) {4};
  \node (5) at (2, 1) {5};
  \node (6) at (-1.35, -0.2) {6};
  \node (7) at (1.35, -0.2) {7};
  \node (8) at (0, -1) {8};
  \draw (1) -- (2);
  \draw (1) -- (3);
  \draw (1) -- (4);
  \draw (1) -- (5);
  \draw (2) -- (6);
  \draw (3) -- (6);
  \draw (4) -- (7);
  \draw (5) -- (7);
  \draw (6) -- (8);
  \draw (7) -- (8);
\end{tikzpicture}
\end{center}
\begin{choices}
\choice 5, 1, 7, 2, 3, 4, 8, 6
\choice 5, 1, 2, 6, 7, 8, 3, 4
\choice 5, 7, 8, 6, 2, 1, 3, 4
\choice 5, 4, 1, 6, 7, 8, 3, 2
\end{choices}
\textbf{Đáp án:} A
\textbf{Giải thích:} Cây nhị phân tìm kiếm tối ưu (Optimal BST) sử dụng quy hoạch động để giảm thiểu thời gian tìm kiếm trung bình dựa trên tần suất truy cập khóa.
\end{question}

\begin{question}{6.85}
% [ID gốc: 4.23]
% [Nguồn: Ngân hàng CSD203 (Câu #155)]
What is the shortest path from node A to node E?
% Graph description:
% Vertices: A, B, C, D, E, F, G
% Weighted edges:
% (A, B) = 2, (A, C) = 5, (A, E) = 12
% (B, F) = 3, (B, C) = 2, (B, E) = 2
% (C, E) = 2
% (F, D) = 14
% (D, E) = 4, (D, G) = 4
% (E, G) = 3
\begin{center}
\begin{tikzpicture}[auto, node distance=2cm, every node/.style={circle, draw, minimum size=0.6cm}]
  \node (A) at (0, 1.5) {A};
  \node (B) at (2, 2.5) {B};
  \node (C) at (2, 0.5) {C};
  \node (F) at (4, 3.5) {F};
  \node (E) at (6, 0.5) {E};
  \node (D) at (6, 2.5) {D};
  \node (G) at (8, 1.5) {G};
  \draw (A) -- node[above] {2} (B);
  \draw (A) -- node[below] {5} (C);
  \draw (B) -- node[above] {3} (F);
  \draw (B) -- node[left] {2} (C);
  \draw (B) -- node[above] {2} (E);
  \draw (C) -- node[below] {2} (E);
  \draw (F) -- node[above] {14} (D);
  \draw (D) -- node[left] {4} (E);
  \draw (D) -- node[above] {4} (G);
  \draw (E) -- node[below] {3} (G);
  \draw (A) to[bend right=60] node[below] {12} (E);
\end{tikzpicture}
\end{center}
\begin{choices}
\choice A, B, C, E
\choice A, G, E
\choice A, B, E
\choice A, C, E
\end{choices}
\textbf{Đáp án:} C
\textbf{Giải thích:} Thuật toán sắp xếp chọn (Selection Sort) thực hiện ít phép hoán đổi nhất trong số các thuật toán cơ bản, tối đa $n - 1$ lần hoán đổi.
\end{question}

\begin{question}{6.86}
% [ID gốc: 4.24]
% [Nguồn: Ngân hàng CSD203 (Câu #172)]
Given a graph in the figure, the breadth first traversal from node A is (in ABC order):
\begin{center}
\begin{tikzpicture}[every node/.style={circle, draw, minimum size=0.7cm}]
  \node (A) at (0, 1.5) {A};
  \node (B) at (2, 1.5) {B};
  \node (C) at (4, 1.5) {C};
  \node (D) at (1, 0) {D};
  \node (E) at (3, 0) {E};
  \draw (A) -- (B);
  \draw (B) -- (C);
  \draw (A) -- (D);
  \draw (B) -- (D);
  \draw (D) -- (E);
  \draw (C) -- (E);
\end{tikzpicture}
\end{center}
\begin{choices}
\choice A, D, E, C, B
\choice A, B, C, D, E
\choice A, B, D, C, E
\choice A, B, C, D, E
\end{choices}
\textbf{Đáp án:} C
\textbf{Giải thích:} Duyệt BFS từ A theo thứ tự bảng chữ cái: thăm A $\to$ thăm các đỉnh kề B, D $\to$ thăm các đỉnh kề tiếp theo C, E. Thứ tự duyệt là A, B, D, C, E.
\end{question}

\begin{question}{6.87}
% [ID gốc: 4.25]
% [Nguồn: Ngân hàng CSD203 (Câu #186)]
Given a graph in the figure, the adjacency matrix of the graph is (visit nodes in ABCD order):
\begin{center}
\begin{tikzpicture}[every node/.style={circle, draw, minimum size=0.7cm}, >=stealth]
  \node (A) at (0, 2) {A};
  \node (B) at (2.5, 2) {B};
  \node (D) at (1, 0) {D};
  \node (C) at (3, 0.5) {C};
  \draw[->, thick] (A) -- (B);
  \draw[->, thick] (A) -- (D);
  \draw[->, thick] (D) -- (B);
  \draw[->, thick] (B) -- (C);
  \draw[->, thick] (D) -- (C);
\end{tikzpicture}
\end{center}
\begin{choices}
\choice [[0, 1, 0, 1], [0, 0, 1, 0], [0, 0, 0, 0], [0, 1, 1, 0]]
\choice [[0, 2, 0, 0], [0, 0, 7, 0], [0, 0, 0, 0], [0, 4, 8, 0]]
\choice [[0, 1, 0, 1], [1, 0, 1, 1], [0, 1, 0, 1], [1, 1, 1, 0]]
\choice [[0, 2, 0, 9], [0, 0, 7, 4], [0, 7, 0, 8], [9, 4, 8, 0]]
\end{choices}
\textbf{Đáp án:} A
\textbf{Giải thích:} Ma trận kề của đồ thị có hướng theo thứ tự A, B, C, D: hàng A có cạnh tới B, D; hàng B tới C; hàng C không có cạnh ra; hàng D tới B, C $\to$ \texttt{[[0,1,0,1],[0,0,1,0],[0,0,0,0],[0,1,1,0]]}.
\end{question}

\begin{question}{6.88}
% [ID gốc: 4.26]
% [Nguồn: Ngân hàng CSD203 (Câu #189)]
What are incident edges on vertex V?
\begin{center}
\begin{tikzpicture}[every node/.style={circle, draw, minimum size=0.6cm}]
  \node (U) at (0, 1.5) {U};
  \node (V) at (1.5, 3) {V};
  \node (W) at (1.5, 0) {W};
  \node (X) at (3, 1.5) {X};
  \node (Y) at (3, -1.5) {Y};
  \node (Z) at (5, 1.5) {Z};
  \draw (U) -- node[above left, draw=none] {a} (V);
  \draw (V) -- node[above right, draw=none] {b} (X);
  \draw (U) -- node[below left, draw=none] {c} (W);
  \draw (V) -- node[right, draw=none] {d} (W);
  \draw (W) -- node[below right, draw=none] {e} (X);
  \draw (W) -- node[below left, draw=none] {f} (Y);
  \draw (X) -- node[right, draw=none] {g} (Y);
  \draw (X) to[bend left=30] node[above, draw=none] {h} (Z);
  \draw (X) to[bend right=30] node[below, draw=none] {i} (Z);
  \draw (Z) to[out=45, in=-45, looseness=4] node[above right, draw=none] {j} (Z);
\end{tikzpicture}
\end{center}
\begin{choices}
\choice a, b, d
\choice d
\choice a
\choice b
\end{choices}
\textbf{Đáp án:} A
\textbf{Giải thích:} Các cạnh liên thuộc với đỉnh V trong đồ thị đã cho là các cạnh nối trực tiếp với V: a (nối U), b (nối X), d (nối W).
\end{question}

\begin{question}{6.89}
% [ID gốc: 4.27]
% [Nguồn: Ngân hàng CSD203 (Câu #202)]
Given a graph in the figure, the weighted matrix of the graph is:
% Graph vertices: A, B, C, D, E
% Edges: (A,B)=2, (B,C)=7, (C,D)=8, (D,E)=5, (E,A)=4
\begin{center}
\begin{tikzpicture}[every node/.style={circle, draw, minimum size=0.7cm}]
\node (A) at (0, 2) {A};
\node (B) at (2, 2) {B};
\node (C) at (2.6, 0.5) {C};
\node (D) at (1, -0.5) {D};
\node (E) at (-0.6, 0.5) {E};
\draw (A) -- node[above, draw=none] {2} (B);
\draw (B) -- node[right, draw=none] {7} (C);
\draw (C) -- node[below right, draw=none] {8} (D);
\draw (D) -- node[below left, draw=none] {5} (E);
\draw (E) -- node[left, draw=none] {4} (A);
\end{tikzpicture}
\end{center}
\begin{choices}
\choice [[1, 0, 1, 0, 0], [0, 1, 0, 1, 1], [1, 0, 1, 0, 1], [0, 1, 0, 1, 0],[0, 1, 1, 0, 1]]
\choice [[0, 2, 0, 1, 4], [1, 0, 1, 0, 0], [0, 7, 0, 8, 5], [9, 0, 8, 0, 5],[4, 0, 0, 5, 0]]
\choice [[0, 1, 0, 1, 1], [1, 0, 1, 0, 0], [0, 1, 0, 1, 0], [1, 0, 1, 0, 1],[1, 0, 0, 1, 0]]
\choice [[0, 2, 0, 0, 4], [2, 0, 7, 0, 0], [0, 7, 0, 8, 0], [0, 0, 8, 0, 5],[4, 0, 0, 5, 0]]
\end{choices}
\textbf{Đáp án:} D
\textbf{Giải thích:} Ma trận trọng số ghi nhận trọng số cạnh giữa các cặp đỉnh: $A-B=2$, $B-C=7$, $C-D=8$, $D-E=5$, $E-A=4$, các vị trí không có cạnh nối trực tiếp mang giá trị 0, khớp chính xác với ma trận ở phương án D.
\end{question}

\begin{question}{6.90}
% [ID gốc: 4.28]
% [Nguồn: Ngân hàng CSD203 (Câu #225)]
What does an adjacency list represent in a graph data structure?
\begin{choices}
\choice It represents the vertices of the graph.
\choice It represents the edges of the graph.
\choice It stores the adjacent vertices for each vertex.
\choice It stores the weights of the edges in the graph.
\end{choices}
\textbf{Đáp án:} C
\textbf{Giải thích:} Danh sách kề (Adjacency List) biểu diễn đồ thị bằng cách lưu trữ danh sách các đỉnh kề (adjacent vertices) cho mỗi đỉnh trong đồ thị.
\end{question}

\begin{question}{6.91}
% [ID gốc: 4.29]
% [Nguồn: Ngân hàng CSD203 (Câu #227)]
What is the minimum spanning tree (MST)?
\begin{choices}
\choice Spanning subgraph
\choice Spanning tree
\choice Spanning tree of a weighted graph with minimum total edge weight
\choice Shortest path
\end{choices}
\textbf{Đáp án:} C
\textbf{Giải thích:} Cây khung nhỏ nhất (Minimum Spanning Tree - MST) của đồ thị vô hướng có trọng số là một cây khung nối tất cả các đỉnh sao cho tổng trọng số của các cạnh là nhỏ nhất có thể.
\end{question}

\begin{question}{6.92}
% [ID gốc: 4.30]
% [Nguồn: Ngân hàng CSD203 (Câu #228)]
What are incident edges on vertex V?
% Graph description:
% Vertices: U, V, W, X, Y, Z
% Edges:
%   a connects U and V
%   b connects V and X
%   c connects U and W
%   d connects V and W
%   e connects W and X
%   f connects W and Y
%   g connects X and Y
%   h, i connect X and Z (multiple edges)
%   j is a self-loop on Z
\begin{center}
\begin{tikzpicture}[
    vertex/.style={circle, draw=blue!70, fill=yellow!40, thick, inner sep=0pt, minimum size=7mm, font=\small\bfseries},
    edge label/.style={font=\small, fill=white, inner sep=1pt}
]
    \node[vertex] (U) at (0, 2) {U};
    \node[vertex] (V) at (2, 4) {V};
    \node[vertex] (W) at (2, 0) {W};
    \node[vertex] (X) at (4, 2) {X};
    \node[vertex] (Y) at (4, -2) {Y};
    \node[vertex] (Z) at (6.5, 2) {Z};

    \draw (U) -- node[edge label, above left] {a} (V);
    \draw (V) -- node[edge label, above right] {b} (X);
    \draw (U) -- node[edge label, below left] {c} (W);
    \draw (V) -- node[edge label, right] {d} (W);
    \draw (W) -- node[edge label, above left] {e} (X);
    \draw (W) -- node[edge label, below left] {f} (Y);
    \draw (X) -- node[edge label, right] {g} (Y);
    \draw (X) to[bend left=30] node[edge label, above] {h} (Z);
    \draw (X) to[bend right=30] node[edge label, below] {i} (Z);
    \draw (Z) to[out=40, in=-40, looseness=5] node[edge label, right] {j} (Z);
\end{tikzpicture}
\end{center}
\begin{choices}
\choice a
\choice b
\choice d
\choice a, b, d
\end{choices}
\textbf{Đáp án:} D
\textbf{Giải thích:} Cạnh liên thuộc (Incident edge) với một đỉnh là cạnh có một đầu mút gắn trực tiếp vào đỉnh đó. Trên đồ thị, các cạnh nối trực tiếp với đỉnh $V$ gồm có $a$ (nối với $U$), $b$ (nối với $X$), và $d$ (nối với $W$), tức $\{a, b, d\}$.
\end{question}

\begin{question}{6.93}
% [ID gốc: 4.31]
% [Nguồn: Ngân hàng CSD203 (Câu #229)]
In a directed graph represented using an adjacency matrix, what does a value of 1 at matrix entry M[i][j] mean?
\begin{choices}
\choice There is an edge from vertex i to vertex j.
\choice There is an edge from vertex j to vertex i.
\choice There is no edge between vertex i and vertex j.
\choice The graph is undirected.
\choice There is a self-loop at vertex i.
\choice There is a weighted edge from vertex j to vertex i.
\end{choices}
\textbf{Đáp án:} A
\textbf{Giải thích:} Trong đồ thị có hướng được biểu diễn bằng ma trận kề $M$, giá trị $M[i][j] = 1$ biểu thị rằng có một cung (cạnh có hướng) đi từ đỉnh $i$ sang đỉnh $j$.
\end{question}

\begin{question}{6.94}
% [ID gốc: 4.32]
% [Nguồn: Ngân hàng CSD203 (Câu #257)]
How does the Prim algorithm select the next edge to add to the minimum spanning tree?
\begin{choices}
\choice It chooses the edge with the minimum weight that connects a vertex in the tree to a vertex not yet in the tree.
\choice It chooses the edge with the minimum weight that connects two vertices not yet in the tree.
\choice It chooses an edge randomly from the set of edges that connect a vertex in the tree to a vertex not yet in the tree.
\choice It chooses the edge with the maximum weight that connects a vertex in the tree to a vertex not yet in the tree.
\end{choices}
\textbf{Đáp án:} A
\textbf{Giải thích:} Thuật toán Prim xây dựng cây khung nhỏ nhất bằng cách tại mỗi bước luôn chọn cạnh có trọng số nhỏ nhất nối một đỉnh đã thuộc cây với một đỉnh chưa thuộc cây.
\end{question}

\begin{question}{6.95}
% [ID gốc: 4.33]
% [Nguồn: Ngân hàng CSD203 (Câu #258)]
Which of the following algorithms guarantees the shortest path tree from a single source vertex to all other vertices in a weighted graph, assuming non-negative edge weights?
\begin{choices}
\choice Depth-First Search (DFS)
\choice Breadth-First Search (BFS)
\choice Prim's algorithm
\choice Dijkstra's algorithm
\end{choices}
\textbf{Đáp án:} D
\textbf{Giải thích:} Thuật toán Dijkstra đảm bảo tìm được cây đường đi ngắn nhất từ một đỉnh nguồn đến tất cả các đỉnh còn lại trong đồ thị có trọng số các cạnh không âm.
\end{question}

\begin{question}{6.96}
% [ID gốc: 4.34]
% [Nguồn: Ngân hàng CSD203 (Câu #259)]
Given a graph in the figure, the Breadth first traversal from node 1 is (in 1,2,3 order):
\begin{center}
\begin{tikzpicture}[node distance=2cm, every node/.style={circle, draw, fill=green!60!black, text=white, font=\bfseries}]
  \node (0) at (0, 2) {0};
  \node (1) at (2, 2) {1};
  \node (2) at (3.5, 1) {2};
  \node (3) at (2, 0) {3};
  \node (4) at (0, 0) {4};
  \draw (0) -- (1);
  \draw (0) -- (4);
  \draw (1) -- (4);
  \draw (1) -- (3);
  \draw (1) -- (2);
  \draw (2) -- (3);
  \draw (4) -- (3);
\end{tikzpicture}
\end{center}
\begin{choices}
\choice 1, 0, 2, 3, 4
\choice 0, 1, 2, 3, 4
\choice 1, 2, 3, 4, 0
\choice 1, 0, 4, 3, 2
\end{choices}
\textbf{Đáp án:} A
\textbf{Giải thích:} Duyệt BFS từ đỉnh 1: thăm đỉnh 1, đưa các đỉnh kề chưa duyệt của 1 vào hàng đợi theo thứ tự nhãn tăng dần (0, 2, 3, 4), kết quả thứ tự duyệt là 1, 0, 2, 3, 4.
\end{question}

\begin{question}{6.97}
% [ID gốc: 4.35]
% [Nguồn: Ngân hàng CSD203 (Câu #260)]
How can the shortest-path tree be reconstructed after executing Dijkstra's algorithm?
\begin{choices}
\choice By maintaining a 'parent' pointer for each vertex, indicating the predecessor vertex on the shortest path from the source vertex.
\choice By tracing back from the destination vertex to the source vertex using a depth-first search.
\choice By iterating through the vertices in the order they were added to the "cloud".
\choice The shortest-path tree is implicitly constructed during the algorithm's execution and doesn't require a separate reconstruction step.
\end{choices}
\textbf{Đáp án:} A
\textbf{Giải thích:} Sau khi chạy thuật toán Dijkstra, cây đường đi ngắn nhất được tái dựng lại nhờ lưu mảng/con trỏ đỉnh cha (\texttt{parent} pointer) cho mỗi đỉnh ghi nhận đỉnh đi trước trên đường đi ngắn nhất.
\end{question}

\begin{question}{6.98}
% [ID gốc: 4.36]
% [Nguồn: Ngân hàng CSD203 (Câu #262)]
Given a set $V = \{V1, V2, ...Vn\}$ of $n$ vertices, how many undirected graphs (not necessarily connected) can be constructed out of $V$?
\begin{choices}
\choice $n(n - 1)/2$
\choice $2^n$
\choice $n!$
\choice $2^{n(n-1)/2}$
\end{choices}
\textbf{Đáp án:} D
\textbf{Giải thích:} Số cạnh tối đa có thể có giữa $n$ đỉnh vô hướng là $C_n^2 = \frac{n(n-1)}{2}$. Với mỗi cặp đỉnh ta có 2 lựa chọn (có cạnh hoặc không), nên tổng số đồ thị vô hướng có thể tạo là $2^{n(n-1)/2}$.
\end{question}

\begin{question}{6.99}
% [ID gốc: 4.37]
% [Nguồn: Đề thi FE/RE CSD203 (Câu #8)]
% Quiz: CSD203_SU26_FE - Question #8 (ID: 87d9d576-6d70-4352-b987-04348d415e73)
A graph in which exists a path between any two nodes is called:
\begin{choices}
\choice Complete graph
\choice Connected graph
\choice Digraph
\choice In-directed graph
\end{choices}
\textbf{Đáp án:} B
\textbf{Giải thích:} Đồ thị liên thông (connected graph) là đồ thị mà giữa bất kỳ hai đỉnh nào cũng tồn tại ít nhất một đường đi nối liền chúng.
\end{question}

\begin{question}{6.100}
% [ID gốc: 4.38]
% [Nguồn: Đề thi FE/RE CSD203 (Câu #12)]
% Quiz: CSD203_SU26_FE - Question #12 (ID: 6540fd1b-da22-4ae1-a171-e7a1f064d522)
The data structure required for Breadth First Traversal on a graph is:
\begin{choices}
\choice Stack
\choice Queue
\choice Array
\choice Linked list
\end{choices}
\textbf{Đáp án:} B
\textbf{Giải thích:} Thuật toán duyệt theo chiều rộng (Breadth First Search - BFS) trên đồ thị sử dụng cấu trúc dữ liệu hàng đợi (queue) để quản lý thứ tự thăm các đỉnh.
\end{question}

\begin{question}{6.101}
% [ID gốc: 4.39]
% [Nguồn: Đề thi FE/RE CSD203 (Câu #16)]
% Quiz: CSD203_SU26_FE - Question #16 (ID: ffad1dc9-2434-4701-8405-4363e1e4a14a)
Which algorithm is used in route finding on maps?
\begin{choices}
\choice DFS
\choice Prim's
\choice Dijkstra's
\choice BFS
\end{choices}
\textbf{Đáp án:} C
\textbf{Giải thích:} Thuật toán Dijkstra được sử dụng rộng rãi trong việc tìm đường đi ngắn nhất trên bản đồ (route finding).
\end{question}

\begin{question}{6.102}
% [ID gốc: 4.40]
% [Nguồn: Đề thi FE/RE CSD203 (Câu #32)]
% Quiz: CSD203_SU26_FE - Question #32 (ID: 965c3e20-6f0c-4519-ab38-6dc53a45a47a)
Which problem is solved using the Longest Common Subsequence (LCS) dynamic programming technique?
\begin{choices}
\choice Data compression
\choice Graph traversal
\choice DNA sequence comparison
\choice Matrix multiplication
\end{choices}
\textbf{Đáp án:} C
\textbf{Giải thích:} Kỹ thuật quy hoạch động tìm chuỗi con chung dài nhất (LCS) được áp dụng phổ biến trong việc so sánh chuỗi gene DNA.
\end{question}

\begin{question}{6.103}
% [ID gốc: 4.41]
% [Nguồn: Đề thi FE/RE CSD203 (Câu #41)]
% Quiz: CSD203_SU26_FE - Question #41 (ID: e287fff1-efe3-4f02-84bd-7c5982c47690)
What is the function of the following pseudocode?
\begin{codebox}
Algorithm Traversal(G, u):
    Input: A graph G and a vertex u of G
    Output: A collection of vertices reachable from u, with their discovery edges
    for each outgoing edge e = (u, v) of u do
        if vertex v has not been visited then
            Mark vertex v as visited (via edge e)
            Recursively call Traversal(G, v)
\end{codebox}
\begin{choices}
\choice It performs a breadth-first search (BFS) starting from vertex u
\choice It performs a depth-first search (DFS) starting from vertex u
\choice It finds the shortest path from u to all other vertices
\choice It detects cycles in the graph starting from vertex u
\end{choices}
\textbf{Đáp án:} B
\textbf{Giải thích:} Đoạn mã giả thể hiện chiến lược thăm đỉnh lân cận rồi đệ quy đi sâu xuống các đỉnh chưa thăm, đây chính là đặc trưng của thuật toán tìm kiếm theo chiều sâu (DFS).
\end{question}

\begin{question}{6.104}
% [ID gốc: 4.42]
% [Nguồn: Đề thi FE/RE CSD203 (Câu #42)]
% Quiz: CSD203_SU26_FE - Question #42 (ID: fcdb73cb-b157-4a9d-9e1a-3a9a9573f396)
Given a set V = {V1, V2, ..., Vn} of n vertices, how many undirected graphs (not necessarily connected) can be constructed out of V?
\begin{choices}
\choice $\frac{n(n-1)}{2}$
\choice $2^n$
\choice n!
\choice $2^{$\frac{n(n-1)}{2}$}$
\end{choices}
\textbf{Đáp án:} D
\textbf{Giải thích:} Một đồ thị vô hướng với n đỉnh có tối đa $\frac{n(n-1)}{2}$ cạnh có thể xảy ra. Mỗi cạnh có 2 trạng thái (tồn tại hoặc không), nên tổng số đồ thị có thể tạo ra là $2^{$\frac{n(n-1)}{2}$}$.
\end{question}

\begin{question}{6.105}
% [ID gốc: 4.43]
% [Nguồn: Đề thi FE/RE CSD203 (Câu #50)]
% Quiz: CSD203_SU26_FE - Question #50 (ID: 90986b76-0dca-4f2b-9eff-afc971504741)
What is the shortest path from node A to node G?
\begin{center}
\begin{tikzpicture}[every node/.style={circle, draw, minimum size=0.8cm}, auto, semithick]
  \node[fill=yellow!50] (A) at (0, 2) {A};
  \node (B) at (2, 3.5) {B};
  \node (C) at (2, 0.5) {C};
  \node (F) at (4, 4.8) {F};
  \node (E) at (6, 0.5) {E};
  \node (D) at (6, 3.5) {D};
  \node (G) at (8, 2) {G};
  
  \draw (A) -- node {2} (B);
  \draw (A) -- node {5} (C);
  \draw (B) -- node {2} (C);
  \draw (B) -- node {3} (F);
  \draw (B) -- node[near start] {2} (E);
  \draw (C) -- node {2} (E);
  \draw (F) -- node {14} (D);
  \draw (D) -- node {4} (E);
  \draw (D) -- node {4} (G);
  \draw (E) -- node {3} (G);
  \draw (A) to[bend right=45] node[below] {12} (E);
\end{tikzpicture}
\end{center}
\begin{choices}
\choice A-B-C-E-G
\choice A-G
\choice A-C-E-G
\choice A-B-E-G
\end{choices}
\textbf{Đáp án:} D
\textbf{Giải thích:} Danh sách cạnh đọc được từ hình: A–B = 2, A–C = 5, A–E = 12 (đường cong dưới cùng), B–C = 2, B–F = 3, B–E = 2 (đường chéo từ B xuống E), C–E = 2, F–D = 14, D–E = 4, D–G = 4, E–G = 3.
\end{question}

\begin{question}{6.106}
% [ID gốc: 4.44]
% [Nguồn: Đề thi FE/RE CSD203 (Câu #60)]
% Quiz: CSD203_SU26_RE - Question #10 (ID: 1216fe2a-8e21-4be7-bd90-e6eb175f0694)
Given a graph in the figure, the depth first traversal from node D is (visit nodes in ABC order):
\begin{center}
\begin{tikzpicture}[every node/.style={circle, draw, minimum size=0.7cm}, semithick]
  \node (A) at (1, 3) {A};
  \node (B) at (3.5, 3) {B};
  \node (C) at (4, 1.2) {C};
  \node (D) at (2.25, 0) {D};
  \node (E) at (0.5, 1.5) {E};
  \draw (A) -- (B);
  \draw (B) -- (C);
  \draw (C) -- (D);
  \draw (D) -- (E);
  \draw (E) -- (A);
  \draw (A) -- (D);
\end{tikzpicture}
\end{center}
\begin{choices}
\choice D, E, A, B, C
\choice D, A, B, C, E
\choice D, C, B, A, E
\choice D, C, E, A, B
\end{choices}
\textbf{Đáp án:} D
\textbf{Giải thích:} Duyệt theo chiều sâu (DFS) bắt đầu từ D, ưu tiên đỉnh kề theo thứ tự bảng chữ cái: từ D chọn C trước, từ C đi tiếp theo nhánh sâu nhất rồi mới quay lui. Thứ tự thăm đúng là D, C, E, A, B.
\end{question}

\begin{question}{6.107}
% [ID gốc: 4.45]
% [Nguồn: Đề thi FE/RE CSD203 (Câu #64)]
% Quiz: CSD203_SU26_RE - Question #14 (ID: 7a4f97ad-270d-480b-8e25-41e08f34c2f0)
Which algorithm is used in route finding on maps?
\begin{choices}
\choice DFS
\choice Prim's
\choice Dijkstra's
\choice BFS
\end{choices}
\textbf{Đáp án:} C
\textbf{Giải thích:} Thuật toán Dijkstra được dùng để tìm đường đi ngắn nhất giữa các đỉnh trong đồ thị có trọng số không âm, ứng dụng thực tế trong hệ thống định vị bản đồ.
\end{question}

\begin{question}{6.108}
% [ID gốc: 4.46]
% [Nguồn: Đề thi FE/RE CSD203 (Câu #67)]
% Quiz: CSD203_SU26_RE - Question #17 (ID: 47549ab9-00aa-4ec3-8027-d0928d80a3c4)
What is the BFS traversal order of the graph starting from vertex 0?
\begin{center}
\begin{tikzpicture}[every node/.style={circle, draw=red, thick, minimum size=0.7cm, text=red, font=\bfseries}, >=stealth, thick]
  \node (0) at (1, 0) {0};
  \node (1) at (3, -0.2) {1};
  \node (2) at (0.3, 2.2) {2};
  \node (3) at (2.2, 2) {3};
  \node (4) at (4, 1.8) {4};
  \node (5) at (1.3, 4.3) {5};
  \node (6) at (3.3, 4.2) {6};
  
  \draw[->] (0) -- (1);
  \draw[->] (0) -- (2);
  \draw[->] (1) -- (3);
  \draw[->] (2) -- (3);
  \draw[->] (2) -- (5);
  \draw[->] (3) -- (4);
  \draw[<->] (5) -- (6);
  \draw[<->] (6) -- (4);
\end{tikzpicture}
\end{center}
\begin{choices}
\choice 0, 2, 1, 3, 5, 4, 6
\choice 0, 2, 1, 3, 5, 6, 4
\choice 0, 1, 2, 5, 3, 6, 4
\choice 0, 1, 2, 3, 5, 4, 6
\end{choices}
\textbf{Đáp án:} D
\textbf{Giải thích:} Duyệt theo chiều rộng (BFS) từ đỉnh 0 thăm tất cả đỉnh kề trực tiếp của 0 trước (1, 2), sau đó lần lượt thăm các đỉnh kề của chúng theo thứ tự hàng đợi (3, 5, 4, 6).
\end{question}

\begin{question}{6.109}
% [ID gốc: 4.47]
% [Nguồn: Đề thi FE/RE CSD203 (Câu #75)]
% Quiz: CSD203_SU26_RE - Question #25 (ID: 58f99c17-3378-43eb-b7cf-31a96d7fae06)
The number of elements in the weighted matrix of a graph having 5 vertices is _____
\begin{choices}
\choice 5
\choice 16
\choice 25
\choice 36
\end{choices}
\textbf{Đáp án:} C
\textbf{Giải thích:} Ma trận trọng số (weighted matrix / adjacency matrix) của đồ thị n đỉnh có kích thước n × n. Với n = 5, số phần tử là 5 × 5 = 25.
\end{question}

\begin{question}{6.110}
% [ID gốc: 4.48]
% [Nguồn: Đề thi FE/RE CSD203 (Câu #79)]
% Quiz: CSD203_SU26_RE - Question #29 (ID: 5d209585-8d58-4c7a-a62d-464494e927ed)
Consider the following undirected graph with edge weights as the following image.

The number of minimum-weight spanning trees of the graph is ________.
\begin{center}
\begin{tikzpicture}[every node/.style={circle, fill=black, minimum size=0.45cm}, auto, semithick]
  \node (v11) at (0, 3) {};
  \node (v12) at (3, 3) {};
  \node (v13) at (6, 3) {};
  
  \node (v21) at (0, 1.5) {};
  \node (v22) at (3, 1.5) {};
  \node (v23) at (6, 1.5) {};
  
  \node (v31) at (0, 0) {};
  \node (v32) at (3, 0) {};
  \node (v33) at (6, 0) {};
  
  \draw (v11) -- node[draw=none, fill=none, above] {\large 0.1} (v12);
  \draw (v12) -- node[draw=none, fill=none, above] {\large 0.1} (v13);
  \draw (v21) -- node[draw=none, fill=none, above] {\large 0.1} (v22);
  \draw (v22) -- node[draw=none, fill=none, above] {\large 0.9} (v23);
  \draw (v31) -- node[draw=none, fill=none, above] {\large 0.9} (v32);
  \draw (v32) -- node[draw=none, fill=none, above] {\large 0.1} (v33);
  
  \draw (v11) -- node[draw=none, fill=none, left] {\large 0.9} (v21);
  \draw (v21) -- node[draw=none, fill=none, left] {\large 0.1} (v31);
  \draw (v12) -- node[draw=none, fill=none, left] {\large 0.9} (v22);
  \draw (v22) -- node[draw=none, fill=none, left] {\large 0.1} (v32);
  \draw (v13) -- node[draw=none, fill=none, left] {\large 0.9} (v23);
  \draw (v23) -- node[draw=none, fill=none, left] {\large 0.1} (v33);
\end{tikzpicture}
\end{center}
\begin{choices}
\choice 6
\choice 5
\choice 4
\choice 3
\end{choices}
\textbf{Đáp án:} D
\textbf{Giải thích:} Áp dụng thuật toán Kruskal hoặc Prim để chọn các cạnh trọng số nhỏ nhất sao cho không tạo thành chu trình, số lượng cây khung nhỏ nhất (MST) khác nhau tạo được cho đồ thị này là 3.
\end{question}

\begin{question}{6.111}
% [ID gốc: 4.49]
% [Nguồn: Đề thi FE/RE CSD203 (Câu #94)]
% Quiz: CSD203_SU26_RE - Question #44 (ID: 2d2d01c6-cbf9-4f00-87a5-f10b4fde08f1)
What two techniques are combined in an optimized Union-Find structure? (Choose 2 correct answers)
\begin{choices}
\choice Path compression
\choice Topological sorting
\choice Union by size
\choice Greedy choice
\end{choices}
\textbf{Đáp án:} AC
\textbf{Giải thích:} Cấu trúc dữ liệu Union-Find đạt hiệu năng gần như hằng số bằng cách kết hợp hai kỹ thuật: nén đường đi (path compression) và hợp nhất theo kích thước/hạng (union by size hoặc union by rank).
\end{question}

\begin{question}{6.112}
% [ID gốc: 4.50]
% [Nguồn: Đề thi FE/RE CSD203 (Câu #98)]
% Quiz: CSD203_SU26_RE - Question #48 (ID: 882e9e87-45e3-4c32-8941-feed22c0b2ac)
What algorithm is commonly used to find the shortest path in a weighted graph?
\begin{choices}
\choice Depth-First Search (DFS)
\choice Breadth-First Search (BFS)
\choice Dijkstra's Algorithm
\choice Kruskal's Algorithm
\end{choices}
\textbf{Đáp án:} C
\textbf{Giải thích:} Thuật toán Dijkstra được dùng phổ biến nhất để tìm đường đi ngắn nhất từ một đỉnh nguồn đến các đỉnh còn lại trong đồ thị có trọng số không âm.
\end{question}

\newpage
\section*{Chương 7: Thuật toán sắp xếp (Sorting Algorithms)}
\addcontentsline{toc}{section}{Chương 7: Thuật toán sắp xếp (Sorting Algorithms)}

\begin{summarybox}{Tóm tắt lý thuyết trọng tâm: Các thuật toán sắp xếp (Sorting Algorithms)}
\begin{itemize}[noitemsep,topsep=2pt]
  \item \textbf{Sắp xếp đổi chỗ (Bubble Sort):} So sánh các cặp phần tử liền kề và hoán đổi. Tốt nhất: $O(n)$ (khi mảng đã sắp xếp và có cờ kiểm tra), Xấu nhất/Trung bình: $O(n^2)$. Ổn định (Stable).
  \item \textbf{Sắp xếp chọn (Selection Sort):} Mỗi bước tìm phần tử nhỏ nhất trong mảng con chưa sắp xếp và đưa về vị trí đầu. Luôn tốn $\frac{n(n-1)}{2}$ phép so sánh, độ phức tạp mọi trường hợp là $O(n^2)$. Không ổn định (Unstable).
  \item \textbf{Sắp xếp chèn (Insertion Sort):} Chèn từng phần tử vào đúng vị trí trong mảng con đã sắp xếp. Tốt nhất: $O(n)$ (rất hiệu quả khi mảng gần như đã sắp xếp), Xấu nhất: $O(n^2)$. Ổn định.
  \item \textbf{Sắp xếp nhanh (Quick Sort):} Chia để trị (Divide and Conquer), chọn phần tử chốt (pivot), phân hoạch mảng thành hai nửa. Trung bình/Tốt nhất: $O(n \log n)$, Xấu nhất: $O(n^2)$ (khi pivot luôn là phần tử nhỏ nhất/lớn nhất). Thường không ổn định.
  \item \textbf{Sắp xếp trộn (Merge Sort):} Chia mảng thành 2 nửa đệ quy rồi trộn (merge) lại. Mọi trường hợp luôn đạt $O(n \log n)$, tốn thêm $O(n)$ bộ nhớ phụ. Rất thích hợp sắp xếp danh sách liên kết và dữ liệu lớn. Ổn định.
  \item \textbf{Sắp xếp cơ số (Radix Sort) \& Đếm (Counting Sort):} Sắp xếp phi so sánh, đạt độ phức tạp tuyến tính $O(n \cdot k)$ hoặc $O(n + k)$.
  \item \textbf{Cận dưới lý thuyết:} Bất kỳ thuật toán sắp xếp dựa trên so sánh nào cũng cần ít nhất $\Omega(n \log n)$ phép so sánh trong trường hợp xấu nhất.
\end{itemize}
\end{summarybox}

\begin{question}{7.1}
% [ID gốc: 5.4]
% [Nguồn: Ngân hàng CSD203 (Câu #7)]
Given an array A = \{6, 7, 4, 1, 2, 9\} and Quick Sort is used to sort the array A in increasing order. What is the sequence after pass 2?
\begin{choices}
\choice 4 6 7 1 2 9
\choice 1 4 6 7 2 9
\choice 6 7 4 1 2 9
\choice 1 2 4 7 6 9
\end{choices}
\textbf{Đáp án:} D
\textbf{Giải thích:} Phân hoạch QuickSort với pivot là 4 đưa các phần tử nhỏ hơn 4 sang trái và lớn hơn 4 sang phải, kết quả dãy sau pha đầu là phương án D.
\end{question}

\begin{question}{7.2}
% [ID gốc: 5.8]
% [Nguồn: Ngân hàng CSD203 (Câu #20)]
Which sorting algorithm is the best if the list is already in order?
\begin{choices}
\choice Insertion Sort
\choice Selection Sort
\choice Merge Sort
\choice Quick Sort
\end{choices}
\textbf{Đáp án:} A
\textbf{Giải thích:} Khi danh sách đã được sắp xếp sẵn, Insertion Sort chỉ cần 1 phép so sánh ở mỗi vị trí và không cần hoán đổi, đạt độ phức tạp tuyến tính $O(n)$.
\end{question}

\begin{question}{7.3}
% [ID gốc: 5.9]
% [Nguồn: Ngân hàng CSD203 (Câu #23)]
The following sorting algorithms maintain two sub-lists, one sorted and one to be sorted?
\begin{choices}
\choice Quick Sort
\choice Selection Sort
\choice None of these
\choice Merge Sort
\end{choices}
\textbf{Đáp án:} B
\textbf{Giải thích:} Thuật toán sắp xếp chọn (Selection Sort) hoặc sắp xếp chèn (Insertion Sort) duy trì hai danh sách con: phần đã sắp xếp và phần chưa sắp xếp.
\end{question}

\begin{question}{7.4}
% [ID gốc: 5.22]
% [Nguồn: Ngân hàng CSD203 (Câu #62)]
Given an array A = \{6, 7, 3, 4, 1, 2, 9\} and Quick Sort is used to sort the aray A in increasing order with pivot value is 4. How is the array like after the first round?
\begin{choices}
\choice 6, 7, 3, 4, 1, 2, 9
\choice 2, 1, 3, 4, 7, 6, 9
\choice 1, 2, 3, 4, 7, 6, 9
\choice 1, 2, 3, 4, 6, 7, 9
\end{choices}
\textbf{Đáp án:} B
\textbf{Giải thích:} Độ phức tạp thời gian của phép tìm kiếm trong bảng băm trong trường hợp tốt nhất và trung bình là $O(1)$.
\end{question}

\begin{question}{7.5}
% [ID gốc: 5.23]
% [Nguồn: Ngân hàng CSD203 (Câu #63)]
Which sorting technique will be most appropriate to sort 1 GB of data with only 100 MB of available main memory?
\begin{choices}
\choice Insertion sort
\choice Heap sort
\choice Merge sort
\choice Quick sort
\end{choices}
\textbf{Đáp án:} C
\textbf{Giải thích:} Biểu thức hậu tố (Postfix / Reverse Polish Notation) không cần dấu ngoặc vì thứ tự ưu tiên của các phép toán được xác định hoàn toàn qua vị trí xuất hiện của toán tử.
\end{question}

\begin{question}{7.6}
% [ID gốc: 5.26]
% [Nguồn: Ngân hàng CSD203 (Câu #75)]
Given an array as follows: [4, 1, 2, 9, 8]. How is the array like after the first round of Bubble sort?
\begin{choices}
\choice 4, 1, 2, 9, 8
\choice 1, 2, 4, 8, 9
\choice 1, 4, 2, 8, 9
\choice 1, 4, 2, 9, 8
\end{choices}
\textbf{Đáp án:} B
\textbf{Giải thích:} Cấu trúc dữ liệu Stack phù hợp để thực hiện chức năng Undo/Redo trong các phần mềm soạn thảo văn bản.
\end{question}

\begin{question}{7.7}
% [ID gốc: 5.28]
% [Nguồn: Ngân hàng CSD203 (Câu #78)]
The number of swappings needed to sort the numbers 8, 22, 7, 9, 31 in ascending order, using bubble sort is .......
\begin{choices}
\choice 2
\choice 3
\choice 4
\choice 5
\end{choices}
\textbf{Đáp án:} B
\textbf{Giải thích:} Thuật toán sắp xếp vun đống (Heap Sort) có độ phức tạp thời gian $O(n \log n)$ trong mọi trường hợp (tốt nhất, trung bình và xấu nhất).
\end{question}

\begin{question}{7.8}
% [ID gốc: 5.29]
% [Nguồn: Ngân hàng CSD203 (Câu #80)]
Given an array $A = \{3, 7, 6, 9, 1\}$ how is A like after the first phase of selection sort in ascending order?
\begin{choices}
\choice 3, 7, 6, 9, 1
\choice 1, 3, 7, 6, 9
\choice 1, 7, 6, 9, 3
\choice 9, 3, 7, 6, 9
\end{choices}
\textbf{Đáp án:} C
\textbf{Giải thích:} Trong bảng băm, gom cụm sơ cấp (Primary clustering) xảy ra khi các ô bị chiếm đóng tạo thành một khối liên tục dài, làm tăng số lần thăm dò.
\end{question}

\begin{question}{7.9}
% [ID gốc: 5.36]
% [Nguồn: Ngân hàng CSD203 (Câu #115)]
What is the best case for sorting in Merge Sort?
\begin{choices}
\choice $O(n * n)$
\choice $O(n)$
\choice $O(nlogn)$
\choice $O(logn)$
\end{choices}
\textbf{Đáp án:} C
\textbf{Giải thích:} Trong ngăn xếp, nếu thực hiện thao tác \texttt{pop()} khi ngăn xếp đang rỗng sẽ xảy ra lỗi tràn dưới (Stack Underflow).
\end{question}

\begin{question}{7.10}
% [ID gốc: 5.38]
% [Nguồn: Ngân hàng CSD203 (Câu #117)]
Which of the following sorting algorithms should be preferred so that the number of swap operations are minimized in general?
\begin{choices}
\choice Heap Sort
\choice Selection Sort
\choice Insertion Sort
\choice Merge Sort
\end{choices}
\textbf{Đáp án:} B
\textbf{Giải thích:} Thuật toán sắp xếp nhanh (QuickSort) chọn một phần tử làm chốt (pivot) và phân hoạch mảng thành hai phần nhỏ hơn và lớn hơn pivot.
\end{question}

\begin{question}{7.11}
% [ID gốc: 5.39]
% [Nguồn: Ngân hàng CSD203 (Câu #118)]
Given an array X = [ 3, 1, 2, 7, 9, 6], Selection sort is used to sort X in an ascending order. After 2 steps , how does the array look like?
\begin{choices}
\choice [1, 2, 3, 7, 9, 6]
\choice [3, 1, 2, 7, 9, 6]
\choice [1, 3, 2, 7, 9, 6]
\choice [1, 2, 3, 6, 7, 9]
\end{choices}
\textbf{Đáp án:} A
\textbf{Giải thích:} Đồ thị liên thông mạnh (Strongly connected graph) là đồ thị có hướng mà giữa hai đỉnh bất kỳ luôn có đường đi theo chiều mũi tên.
\end{question}

\begin{question}{7.12}
% [ID gốc: 5.58]
% [Nguồn: Ngân hàng CSD203 (Câu #176)]
The number of swappings needed to sort the numbers 2, 7, 8, 9, 1 in ascending order, using bubble sort is ........
\begin{choices}
\choice 2
\choice 3
\choice 4
\choice 5
\end{choices}
\textbf{Đáp án:} C
\textbf{Giải thích:} Sắp xếp 2, 7, 8, 9, 1 bằng Bubble Sort: số lần đổi chỗ đưa số 1 về đầu là 4 lần đổi chỗ liên tiếp với 9, 8, 7, 2.
\end{question}

\begin{question}{7.13}
% [ID gốc: 5.60]
% [Nguồn: Ngân hàng CSD203 (Câu #180)]
Given an array $A = \{6, 7, 4, 1, 2, 9\}$ and Quick Sort is used to sort the aray A in increasing order. What is the sequence after the first phase, the pivot is 4?
\begin{choices}
\choice 6 7 4 1 2 9
\choice 2 1 4 7 6 9
\choice 1 2 4 7 6 9
\choice 1 2 4 6 7 9
\end{choices}
\textbf{Đáp án:} B
\textbf{Giải thích:} Phân hoạch QuickSort cho mảng {6, 7, 4, 1, 2, 9} với pivot = 4 đưa các phần tử nhỏ hơn 4 ({1, 2}) lên trước, kết quả dãy sau pha đầu là 2 1 4 7 6 9.
\end{question}

\begin{question}{7.14}
% [ID gốc: 5.68]
% [Nguồn: Ngân hàng CSD203 (Câu #191)]
What is the worst-case time complexity of QuickSort?
\begin{choices}
\choice $O(n^2)$
\choice $O(n \log n)$
\choice $O(n)$
\choice $O(2^n)$
\end{choices}
\textbf{Đáp án:} A
\textbf{Giải thích:} Trường hợp xấu nhất của QuickSort xảy ra khi pivot chọn luôn là phần tử nhỏ nhất hoặc lớn nhất, dẫn đến độ phức tạp thời gian là $O(n^2)$.
\end{question}

\begin{question}{7.15}
% [ID gốc: 5.69]
% [Nguồn: Ngân hàng CSD203 (Câu #193)]
Which is sorting algorithms that maintain two sub-lists, one sorted and one to be sorted?
\begin{choices}
\choice Selection Sort
\choice Merge Sort
\choice Quick Sort
\choice None of these
\end{choices}
\textbf{Đáp án:} A
\textbf{Giải thích:} Thuật toán sắp xếp chọn (Selection Sort) chia danh sách thành hai phần: danh sách con đã sắp xếp và danh sách con chưa sắp xếp.
\end{question}

\begin{question}{7.16}
% [ID gốc: 5.72]
% [Nguồn: Ngân hàng CSD203 (Câu #197)]
A sorting algorithm that uses the divide and conquer technique is?
\begin{choices}
\choice Quick sort
\choice Selection sort
\choice Bubble sort
\choice Insertion sort
\end{choices}
\textbf{Đáp án:} A
\textbf{Giải thích:} Thuật toán sắp xếp áp dụng kỹ thuật chia để trị (Divide and Conquer) tiêu biểu là Quick Sort và Merge Sort.
\end{question}

\begin{question}{7.17}
% [ID gốc: 5.77]
% [Nguồn: Ngân hàng CSD203 (Câu #236)]
In which sorting, elements in the array are compared with a pivot value?
\begin{choices}
\choice Bubble Sort
\choice Selection Sort
\choice Merge Sort
\choice Quick Sort
\end{choices}
\textbf{Đáp án:} D
\textbf{Giải thích:} Trong thuật toán sắp xếp nhanh (Quick Sort), các phần tử trong mảng được so sánh và phân hoạch dựa trên một giá trị phần tử chốt (pivot value).
\end{question}

\begin{question}{7.18}
% [ID gốc: 5.78]
% [Nguồn: Ngân hàng CSD203 (Câu #237)]
Which sorting algorithm use the divide-and-conquer approach?
\begin{choices}
\choice Quick Sort
\choice Insertion Sort
\choice Bubble Sort
\choice Selection Sort
\end{choices}
\textbf{Đáp án:} A
\textbf{Giải thích:} Thuật toán Quick Sort là thuật toán sắp xếp kinh điển áp dụng phương pháp chia để trị (Divide-and-Conquer).
\end{question}

\begin{question}{7.19}
% [ID gốc: 5.80]
% [Nguồn: Ngân hàng CSD203 (Câu #239)]
In which sorting, consecutive adjacent pairs of elements in the array are compared with each other?
\begin{choices}
\choice Bubble sort
\choice Selection sort
\choice Merge sort
\choice Quick sort
\end{choices}
\textbf{Đáp án:} A
\textbf{Giải thích:} Trong thuật toán sắp xếp nổi bọt (Bubble Sort), các cặp phần tử liền kề (adjacent pairs) liên tục được so sánh và hoán đổi nếu chúng đứng sai thứ tự.
\end{question}

\begin{question}{7.20}
% [ID gốc: 5.85]
% [Nguồn: Ngân hàng CSD203 (Câu #271)]
Which sorting algorithm has a time complexity of $O(n^2)$ in the worst-case scenario, but can perform better than both Bubble Sort and Insertion Sort in practice due to its fewer number of swaps?
\begin{choices}
\choice Insertion Sort
\choice Bubble Sort
\choice Selection Sort
\choice Quick Sort
\end{choices}
\textbf{Đáp án:} C
\textbf{Giải thích:} Thuật toán sắp xếp chọn (Selection Sort) có độ phức tạp xấu nhất $O(n^2)$ nhưng thực hiện tối đa chỉ $n - 1$ lần hoán đổi (swaps), ít hơn đáng kể so với Bubble Sort và Insertion Sort.
\end{question}

\begin{question}{7.21}
% [ID gốc: 5.86]
% [Nguồn: Ngân hàng CSD203 (Câu #272)]
Which sorting algorithm divides the array into halves recursively and then merges them by sorting each half in the middle?
\begin{choices}
\choice Quick Sort
\choice Merge Sort
\choice Heap Sort
\choice Bubble Sort
\end{choices}
\textbf{Đáp án:} B
\textbf{Giải thích:} Thuật toán sắp xếp trộn (Merge Sort) chia đôi mảng một cách đệ quy cho đến khi còn từng phần tử đơn lẻ rồi trộn các mảng con đã sắp xếp lại với nhau.
\end{question}

\begin{question}{7.22}
% [ID gốc: 5.89]
% [Nguồn: Ngân hàng CSD203 (Câu #275)]
Given an array as follows: [6, 1, 5, 3, 7]. How is the array like after the first round of Selection sort?
\begin{choices}
\choice 1, 6, 5, 3, 7
\choice 7, 1, 5, 3, 6
\choice 1, 3, 5, 6, 7
\choice 7, 6, 5, 3, 1
\end{choices}
\textbf{Đáp án:} A
\textbf{Giải thích:} Vòng đầu tiên của Selection Sort tìm phần tử nhỏ nhất trong \texttt{[6, 1, 5, 3, 7]} là 1 và hoán đổi với phần tử đầu tiên là 6, dãy sau vòng 1 là \texttt{1, 6, 5, 3, 7}.
\end{question}

\begin{question}{7.23}
% [ID gốc: 5.91]
% [Nguồn: Đề thi FE/RE CSD203 (Câu #3)]
% Quiz: CSD203_SU26_FE - Question #3 (ID: 4852ccb6-f4d2-4495-9667-8eb50f8aff24)
Which sorting algorithm is most suitable for sorting linked lists due to its efficient use of memory and ability to operate without random access to elements?
\begin{choices}
\choice Merge Sort
\choice Quick Sort
\choice Bubble Sort
\choice Selection Sort
\end{choices}
\textbf{Đáp án:} A
\textbf{Giải thích:} Thuật toán trộn (Merge Sort) hoạt động tuần tự, không yêu cầu truy cập ngẫu nhiên đến các phần tử như Quick Sort, nên rất phù hợp để sắp xếp danh sách liên kết.
\end{question}

\begin{question}{7.24}
% [ID gốc: 5.97]
% [Nguồn: Đề thi FE/RE CSD203 (Câu #26)]
% Quiz: CSD203_SU26_FE - Question #26 (ID: f3a76c5b-289c-44d4-800d-fdfa12d17aa0)
Given the array [47, 12, 85, 33, 9, 71, 56, 24, 63, 18]. What will the array look like after the 3rd iteration of Selection Sort?
\begin{choices}
\choice [9, 12, 18, 33, 47, 71, 56, 24, 63, 85]
\choice [9, 12, 18, 47, 33, 71, 56, 24, 63, 85]
\choice [9, 12, 33, 18, 47, 71, 56, 24, 63, 85]
\choice [9, 12, 18, 33, 85, 71, 56, 24, 63, 47]
\end{choices}
\textbf{Đáp án:} A
\textbf{Giải thích:} Sau 3 bước lặp của thuật toán Selection Sort, 3 phần tử nhỏ nhất được tìm thấy và đưa về đúng vị trí đầu mảng theo thứ tự tăng dần là [9, 12, 18], các phần tử còn lại giữ nguyên trật tự tương đối.
\end{question}

\begin{question}{7.25}
% [ID gốc: 5.99]
% [Nguồn: Đề thi FE/RE CSD203 (Câu #29)]
% Quiz: CSD203_SU26_FE - Question #29 (ID: d00ba0a6-57dc-4b41-a3b6-5188d0e30758)
Which sorting algorithm is stable, meaning it preserves the relative order of equal elements?
\begin{choices}
\choice Merge Sort
\choice Heap Sort
\choice Quick Sort
\choice Selection Sort
\end{choices}
\textbf{Đáp án:} A
\textbf{Giải thích:} Merge Sort là thuật toán sắp xếp ổn định (stable sort), đảm bảo các phần tử có cùng giá trị giữ nguyên thứ tự ban đầu đối với nhau.
\end{question}

\begin{question}{7.26}
% [ID gốc: 5.100]
% [Nguồn: Đề thi FE/RE CSD203 (Câu #31)]
% Quiz: CSD203_SU26_FE - Question #31 (ID: 3008d373-1661-47db-98c0-5f708795e09e)
What is the lower bound for any comparison-based sorting algorithm?
\begin{choices}
\choice $O(n)$
\choice $O(n \log n)$
\choice $O(n^2)$
\choice $O(\log n)$
\end{choices}
\textbf{Đáp án:} B
\textbf{Giải thích:} Giới hạn dưới của các thuật toán sắp xếp dựa trên so sánh trong trường hợp xấu nhất là $O(n \log n)$.
\end{question}

\begin{question}{7.27}
% [ID gốc: 5.107]
% [Nguồn: Đề thi FE/RE CSD203 (Câu #48)]
% Quiz: CSD203_SU26_FE - Question #48 (ID: 22a47d9b-a3cd-4413-b71f-8463b986be72)
Radix Sort is most efficient when sorting:
\begin{choices}
\choice Integers or strings with a fixed length
\choice Data with a wide range of possible values
\choice Small datasets
\choice Unordered linked lists
\end{choices}
\textbf{Đáp án:} A
\textbf{Giải thích:} Thuật toán sắp xếp Radix Sort đạt hiệu suất tối ưu khi xử lý tập dữ liệu số nguyên hoặc chuỗi có độ dài cố định.
\end{question}

\begin{question}{7.28}
% [ID gốc: 5.111]
% [Nguồn: Đề thi FE/RE CSD203 (Câu #57)]
% Quiz: CSD203_SU26_RE - Question #7 (ID: 34a4eafd-a93f-4bfd-99b4-b397f3f62a73)
Which Python keyword/argument allows you to define a custom sorting rule in the sorted() function?
\begin{choices}
\choice key=
\choice cmp=
\choice lambda
\choice func=
\end{choices}
\textbf{Đáp án:} A
\textbf{Giải thích:} Trong Python 3, hàm sorted() dùng tham số key= để nhận một hàm định nghĩa quy tắc sắp xếp tùy chỉnh. Tham số cmp= đã bị loại bỏ từ Python 3.
\end{question}

\begin{question}{7.29}
% [ID gốc: 5.120]
% [Nguồn: Đề thi FE/RE CSD203 (Câu #91)]
% Quiz: CSD203_SU26_RE - Question #41 (ID: aa2f6662-5af5-4c08-aa7d-100658034782)
The number of swappings needed to sort the numbers 8, 22, 7, 9, 31 in ascending order, using bubble sort is _____
\begin{choices}
\choice 2
\choice 3
\choice 4
\choice 5
\end{choices}
\textbf{Đáp án:} B
\textbf{Giải thích:} Mảng ban đầu [8, 22, 7, 9, 31]. Lượt 1: đổi 22 và 7 → [8, 7, 22, 9, 31]; đổi 22 và 9 → [8, 7, 9, 22, 31] (2 lần đổi). Lượt 2: đổi 8 và 7 → [7, 8, 9, 22, 31] (lần thứ 3). Mảng đã sắp xếp xong, tổng số lần hoán đổi là 3.
\end{question}

\newpage
\section*{Chương 8: Các giải thuật bổ trợ (Độ phức tạp, Đệ quy, Tìm kiếm, Huffman \& Xử lý chuỗi)}
\addcontentsline{toc}{section}{Chương 8: Các giải thuật bổ trợ}

\begin{summarybox}{Tóm tắt lý thuyết trọng tâm: Phân tích độ phức tạp, Đệ quy, Tìm kiếm, Huffman \& KMP}
\begin{itemize}[noitemsep,topsep=2pt]
  \item \textbf{Độ phức tạp thuật toán (Big-O):} Đánh giá tốc độ tăng tiệm cận của thời gian chạy hoặc không gian bộ nhớ. Thứ tự tăng dần: $O(1) < O(\log n) < O(n) < O(n \log n) < O(n^2) < O(2^n) < O(n!)$.
  \item \textbf{Thuật toán tìm kiếm:}
    \begin{itemize}[noitemsep]
      \item \textit{Tìm kiếm tuyến tính (Linear Search):} Áp dụng cho mảng bất kỳ. Độ phức tạp: $O(n)$.
      \item \textit{Tìm kiếm nhị phân (Binary Search):} Bắt buộc mảng phải được sắp xếp trước. Mỗi bước so sánh chia đôi không gian tìm kiếm. Độ phức tạp: $O(\log n)$. Số phép so sánh tối đa cho mảng $n$ phần tử là $\lfloor \log_2 n \rfloor + 1$.
    \end{itemize}
  \item \textbf{Giải thuật đệ quy (Recursion):} Hàm tự gọi chính nó. Bắt buộc phải có điều kiện dừng (Base Case). Đệ quy đuôi (Tail Recursion) xảy ra khi lệnh gọi đệ quy là thao tác cuối cùng của hàm, có thể dễ dàng tối ưu hóa khử đệ quy bằng vòng lặp.
  \item \textbf{Nén dữ liệu \& Mã hóa Huffman:} Thuật toán tham lam (Greedy) xây dựng cây nhị phân từ dưới lên bằng cách liên tục ghép 2 nút có tần suất nhỏ nhất (dùng Min-Heap). Ký tự có tần suất xuất hiện càng cao thì mã nhị phân càng ngắn (Prefix Code -- không có mã nào là tiền tố của mã khác).
  \item \textbf{So khớp chuỗi (String Matching \& KMP):} Thuật toán Knuth-Morris-Pratt xây dựng mảng tiền tố (lps array) cho mẫu cần tìm để không bao giờ phải quay lui con trỏ trên chuỗi văn bản chính. Độ phức tạp: $O(n + m)$.
\end{itemize}
\end{summarybox}

\begin{question}{8.1}
% [ID gốc: 5.1]
% [Nguồn: Ngân hàng CSD203 (Câu #2)]
Given a string S = "ABCDEFABCD" and a pattern P = "BC". What is the return value when applying KMP algorithm? (the string is indexed from 0)
\begin{choices}
\choice 2, 8
\choice 0
\choice 1
\choice 1, 7
\end{choices}
\textbf{Đáp án:} D
\textbf{Giải thích:} Mẫu 'BC' xuất hiện trong 'ABCDEFABCD' tại hai vị trí chỉ mục là 1 (trong 'ABCDEF') và 7 (trong 'ABCD'), với chuỗi đánh chỉ số từ 0.
\end{question}

\begin{question}{8.2}
% [ID gốc: 5.2]
% [Nguồn: Ngân hàng CSD203 (Câu #5)]
What will be the output of the following code?
\begin{codebox}
def fun(i,j):
  if(i==0):
    return j
  else:
    return fun(i-1,j+1)

print(fun(4,8))
\end{codebox}
\begin{choices}
\choice 16
\choice 8
\choice 12
\choice 4
\end{choices}
\textbf{Đáp án:} C
\textbf{Giải thích:} Hàm đệ quy \texttt{fun(i, j)} giảm $i$ đi 1 và tăng $j$ lên 1 ở mỗi bước: \texttt{fun(4, 8)} $\to$ \texttt{fun(3, 9)} $\to$ \texttt{fun(2, 10)} $\to$ \texttt{fun(1, 11)} $\to$ \texttt{fun(0, 12)} trả về 12.
\end{question}

\begin{question}{8.3}
% [ID gốc: 5.3]
% [Nguồn: Ngân hàng CSD203 (Câu #6)]
Using Huffman coding to fill missed data into the following table:
\begin{center}
\begin{tabular}{|c|c|c|c|}
\hline
\textbf{ID} & \textbf{Symbol} & \textbf{Frequency} & \textbf{Huffman code} \\
\hline
1 & A & 1 & \\
\hline
2 & B & 2 & \\
\hline
3 & C & 4 & \\
\hline
4 & D & 5 & \\
\hline
5 & E & 6 & \\
\hline
6 & F & 8 & \\
\hline
7 & G & 12 & \\
\hline
8 & H & 32 & \\
\hline
\end{tabular}
\end{center}
\begin{choices}
\choice A: 01111, B: 000, C: 0110, D: 0011, E: 0010, F: 010, G: 01110, H: 1
\choice A: 1, B: 01110, C: 0110, D: 0011, E: 0010, F: 010, G: 000, H: 1
\choice A: 1, B: 01110, C: 0110, D: 0011, E: 0010, F: 010, G: 000, H: 01111
\choice A: 01111, B: 01110, C: 0110, D: 0011, E: 0010, F: 010, G: 000, H: 1
\end{choices}
\textbf{Đáp án:} D
\textbf{Giải thích:} Mã hóa Huffman gán từ mã ngắn hơn cho các ký tự có tần suất xuất hiện cao hơn nhằm tối ưu dung lượng lưu trữ.
\end{question}

\begin{question}{8.4}
% [ID gốc: 5.5]
% [Nguồn: Ngân hàng CSD203 (Câu #9)]
Given data compression rate is 0.25 and length of input data is 8. What is the length of output data?
\begin{choices}
\choice 8
\choice 0
\choice 6
\choice 4
\end{choices}
\textbf{Đáp án:} C
\textbf{Giải thích:} Tỷ lệ nén = (Kích thước sau nén) / (Kích thước gốc) = 0.25. Kích thước dữ liệu nén là $8 \times 0.25 = 2$.
\end{question}

\begin{question}{8.5}
% [ID gốc: 5.6]
% [Nguồn: Ngân hàng CSD203 (Câu #12)]
What will be the output of the following code?
\begin{codebox}
def f(n):
    if (n < 10):
        return n - 5
    return f(n+1);

print(f(4))
\end{codebox}
\begin{choices}
\choice 5
\choice -1
\choice 1
\choice 0
\end{choices}
\textbf{Đáp án:} B
\textbf{Giải thích:} Với $n = 4 < 10$, hàm $f(n)$ trả về giá trị $n - 5 = 4 - 5 = -1$.
\end{question}

\begin{question}{8.6}
% [ID gốc: 5.7]
% [Nguồn: Ngân hàng CSD203 (Câu #14)]
A sorted array contains 16 items. Using binary search, the maximum number of comparisons to search for an item in this array is ............
\begin{choices}
\choice 5
\choice 4
\choice 3
\choice None of above
\end{choices}
\textbf{Đáp án:} A
\textbf{Giải thích:} Tìm kiếm nhị phân trên mảng 16 phần tử cần tối đa $\log_2(16) = 4$ bước so sánh để tìm thấy hoặc kết luận không có phần tử.
\end{question}

\begin{question}{8.7}
% [ID gốc: 5.10]
% [Nguồn: Ngân hàng CSD203 (Câu #25)]
Which is the most appropriate definition for recursion?
\begin{choices}
\choice An in-built method that is automatically called.
\choice A class method that calls another class method.
\choice A function that calls another execution instance of the same function.
\choice A function that calls another function.
\end{choices}
\textbf{Đáp án:} C
\textbf{Giải thích:} Đệ quy là phương pháp lập trình trong đó một hàm tự gọi lại chính nó để giải quyết bài toán con có kích thước nhỏ hơn.
\end{question}

\begin{question}{8.8}
% [ID gốc: 5.11]
% [Nguồn: Ngân hàng CSD203 (Câu #27)]
Which of the following statement is false?
\begin{choices}
\choice Infinite recursion can occur if the base case is not properly defined.
\choice Every recursive function must have a return value.
\choice A recursive function makes the code easier to understand.
\choice Every recursive function must have a base case.
\end{choices}
\textbf{Đáp án:} B
\textbf{Giải thích:} Khẳng định sai là phương án B vì mô tả không đúng nguyên lý hoạt động của cấu trúc dữ liệu hoặc thuật toán tương ứng.
\end{question}

\begin{question}{8.9}
% [ID gốc: 5.12]
% [Nguồn: Ngân hàng CSD203 (Câu #31)]
Which of the following algorithm does not divide the list?
\begin{choices}
\choice Merge Sort
\choice Binary search
\choice Quick Sort
\choice Linear search
\end{choices}
\textbf{Đáp án:} D
\textbf{Giải thích:} Độ phức tạp thời gian của thuật toán tìm kiếm nhị phân trên mảng đã sắp xếp là $O(\log n)$.
\end{question}

\begin{question}{8.10}
% [ID gốc: 5.13]
% [Nguồn: Ngân hàng CSD203 (Câu #33)]
Which search algorithm would be best to use with ordered data?
\begin{choices}
\choice A binary search
\choice Neither binary search nor a linear search
\choice Either binary search or a linear search
\choice A linear search
\end{choices}
\textbf{Đáp án:} A
\textbf{Giải thích:} Trong cấu trúc dữ liệu Heap nhị phân, mỗi nút cha luôn thỏa mãn điều kiện lớn hơn hoặc bằng các con (Max-Heap) hoặc nhỏ hơn hoặc bằng các con (Min-Heap).
\end{question}

\begin{question}{8.11}
% [ID gốc: 5.14]
% [Nguồn: Ngân hàng CSD203 (Câu #36)]
Given a list of 64 elements, how many elements will be checked to look for a value that is larger than the largest value in the list using binary search?
\begin{choices}
\choice 8
\choice 6
\choice 7
\choice 9
\end{choices}
\textbf{Đáp án:} C
\textbf{Giải thích:} Hàm băm hoàn hảo (Perfect hash function) là hàm băm ánh xạ mỗi khóa vào một vị trí duy nhất trong bảng băm mà không xảy ra bất kỳ xung đột nào.
\end{question}

\begin{question}{8.12}
% [ID gốc: 5.15]
% [Nguồn: Ngân hàng CSD203 (Câu #37)]
What will be the output of the following code?
\begin{codebox}
def rec(n):
    if n > 5:
        return n
    else:
        return n + rec(n+2)

print(rec(1))
\end{codebox}
\begin{choices}
\choice 5
\choice 16
\choice 1
\choice 15
\end{choices}
\textbf{Đáp án:} B
\textbf{Giải thích:} Cấu trúc dữ liệu ngăn xếp (Stack) hoạt động theo nguyên lý LIFO (Last In First Out), phần tử đưa vào sau cùng sẽ được lấy ra đầu tiên.
\end{question}

\begin{question}{8.13}
% [ID gốc: 5.16]
% [Nguồn: Ngân hàng CSD203 (Câu #38)]
Which of the following is longest proper prefix which is also suffix of "ABCA"?
\begin{choices}
\choice "ABC"
\choice "ABCA"
\choice "A"
\choice "AB"
\end{choices}
\textbf{Đáp án:} C
\textbf{Giải thích:} Tiền tố đúng dài nhất vừa là hậu tố của 'ABCA' là chuỗi 'A' (độ dài 1).
\end{question}

\begin{question}{8.14}
% [ID gốc: 5.17]
% [Nguồn: Ngân hàng CSD203 (Câu #44)]
Which of the following cases occurs when searching an array using linear search the value to be searched is equal to the first element of the array?
\begin{choices}
\choice The best case
\choice The amortized case
\choice The average case
\choice The worst case
\end{choices}
\textbf{Đáp án:} A
\textbf{Giải thích:} Thuật toán sắp xếp nổi bọt (Bubble Sort) so sánh và hoán đổi các cặp phần tử liền kề nếu chúng sai thứ tự, lặp lại cho đến khi danh sách được sắp xếp.
\end{question}

\begin{question}{8.15}
% [ID gốc: 5.18]
% [Nguồn: Ngân hàng CSD203 (Câu #49)]
The complexity of Linear search algorithm is
\begin{choices}
\choice $O(n)$
\choice $O(\log n)$
\choice $O(n2)$
\choice $O(n \log n)$
\end{choices}
\textbf{Đáp án:} A
\textbf{Giải thích:} Độ phức tạp thời gian trường hợp trung bình của QuickSort là $O(n \log n)$, trong khi trường hợp xấu nhất là $O(n^2)$.
\end{question}

\begin{question}{8.16}
% [ID gốc: 5.19]
% [Nguồn: Ngân hàng CSD203 (Câu #57)]
The concept of prefix and suffix is used in which of the following algorithms?
\begin{choices}
\choice KMP
\choice Boyer-Moore
\choice Brute Force
\choice Advanced Brute Force
\end{choices}
\textbf{Đáp án:} A
\textbf{Giải thích:} Trong danh sách liên kết vòng, nút cuối trỏ về nút đầu (head) giúp việc duyệt có thể tiếp tục liên tục mà không gặp giá trị \texttt{None}.
\end{question}

\begin{question}{8.17}
% [ID gốc: 5.20]
% [Nguồn: Ngân hàng CSD203 (Câu #60)]
What will be the output of the following code?
\begin{codebox}
def rec(n):
    if (n <= 2):
        return n
    else:
        return rec(n - 2) + rec(n-1)

print(rec(4))
\end{codebox}
\begin{choices}
\choice 5
\choice 4
\choice 0
\choice Error
\end{choices}
\textbf{Đáp án:} A
\textbf{Giải thích:} Ma trận kề của đồ thị vô hướng $n$ đỉnh luôn là ma trận đối xứng qua đường chéo chính với kích thước $n \times n$.
\end{question}

\begin{question}{8.18}
% [ID gốc: 5.21]
% [Nguồn: Ngân hàng CSD203 (Câu #61)]
What is disadvantage of Recursion?
\begin{choices}
\choice Recursive calls are expensive (inefficient) as they take up a lot of memory and time
\choice Sequence generation is more difficult with recursion than using some nested iteration.
\choice A complex task can not be broken down into simpler sub-problems using recursion.
\choice None of these
\end{choices}
\textbf{Đáp án:} A
\textbf{Giải thích:} Thuật toán Prim bắt đầu từ một đỉnh tùy ý và mở rộng dần cây khung bằng cách chọn cạnh có trọng số nhỏ nhất nối cây với một đỉnh chưa thuộc cây.
\end{question}

\begin{question}{8.19}
% [ID gốc: 5.24]
% [Nguồn: Ngân hàng CSD203 (Câu #67)]
Given a string S = ``1234512345'' and a pattern P = ``123''. What is brute-force algorithm returns? (the string is indexed from 0)
\begin{choices}
\choice 0
\choice 1
\choice 5
\choice 0, 5
\end{choices}
\textbf{Đáp án:} D
\textbf{Giải thích:} Thuật toán Brute-Force tìm kiếm chuỗi con '123' trong chuỗi '1234512345' tại các vị trí chỉ số 0 và 5.
\end{question}

\begin{question}{8.20}
% [ID gốc: 5.25]
% [Nguồn: Ngân hàng CSD203 (Câu #68)]
Given a definition as follows: ``a recursive function in which the first statement is a recursive call and then the other operations are performed'' What is this?
\begin{choices}
\choice Tail-Recursion
\choice NonTail-Recursion
\choice Recursion
\choice None of these
\end{choices}
\textbf{Đáp án:} B
\textbf{Giải thích:} Độ phức tạp không gian của thuật toán đệ quy phụ thuộc vào độ sâu tối đa của ngăn xếp lời gọi hàm (call stack depth).
\end{question}

\begin{question}{8.21}
% [ID gốc: 5.27]
% [Nguồn: Ngân hàng CSD203 (Câu #77)]
Given an array as follows: [6, 1, 5, 3, 7]. What is the number of key comparisons done to check if 3 is in the array or not by using linear search?
\begin{choices}
\choice 1
\choice 2
\choice 3
\choice 4
\end{choices}
\textbf{Đáp án:} D
\textbf{Giải thích:} Đồ thị đầy đủ (Complete graph) $K_n$ với $n$ đỉnh có tổng cộng $\frac{n(n-1)}{2}$ cạnh vô hướng nối từng cặp đỉnh.
\end{question}

\begin{question}{8.22}
% [ID gốc: 5.30]
% [Nguồn: Ngân hàng CSD203 (Câu #83)]
Which is the most appropriate definition for recursion?
\begin{choices}
\choice A function that calls another execution instance of the same function.
\choice A function that calls another function.
\choice An in-built method that is automatically called.
\choice A class method that calls another class method.
\end{choices}
\textbf{Đáp án:} A
\textbf{Giải thích:} Thuật toán Dijkstra không hoạt động chính xác trên đồ thị có cạnh mang trọng số âm vì nó giả định khoảng cách chỉ có thể tăng dần.
\end{question}

\begin{question}{8.23}
% [ID gốc: 5.31]
% [Nguồn: Ngân hàng CSD203 (Câu #86)]
Given an array as follows: [4, 2, 5, 8, 6] and a search key x using linear search? The Best case of searching occurs when value of x is
\begin{choices}
\choice 2
\choice 4
\choice 6
\choice None of these
\end{choices}
\textbf{Đáp án:} B
\textbf{Giải thích:} Chuyển biểu thức trung tố sang hậu tố sử dụng ngăn xếp Stack để lưu trữ tạm thời các toán tử theo thứ tự ưu tiên.
\end{question}

\begin{question}{8.24}
% [ID gốc: 5.32]
% [Nguồn: Ngân hàng CSD203 (Câu #110)]
Given a string S = ``ABCDEFABCD'' and a pattern P = ``BC''. What is brute-force algorithm returns? (the string is indexed from 0)
\begin{choices}
\choice 0
\choice 1
\choice 7
\choice 1, 7
\end{choices}
\textbf{Đáp án:} D
\textbf{Giải thích:} Mẫu 'BC' xuất hiện trong chuỗi 'ABCDEFABCD' tại chỉ số 1 và 7.
\end{question}

\begin{question}{8.25}
% [ID gốc: 5.33]
% [Nguồn: Ngân hàng CSD203 (Câu #112)]
What does the function \texttt{s.find(pattern)} do?
\begin{choices}
\choice Return the index starting the leftmost occurrence of pattern; else -1
\choice Return the index starting the rightmost occurrence of pattern; else -1
\choice Return 1 as occurrence of pattern; else -1
\choice Return the length of pattern
\end{choices}
\textbf{Đáp án:} A
\textbf{Giải thích:} Phương thức \texttt{s.find(pattern)} trả về chỉ mục bắt đầu của lần xuất hiện đầu tiên của chuỗi mẫu trong chuỗi \texttt{s}, hoặc trả về -1 nếu không tìm thấy.
\end{question}

\begin{question}{8.26}
% [ID gốc: 5.34]
% [Nguồn: Ngân hàng CSD203 (Câu #113)]
What is the condition for applying binary search in an array?
\begin{choices}
\choice The array should not be too long
\choice The array should not be short long
\choice The array should be sorted
\choice The array should not be sorted
\end{choices}
\textbf{Đáp án:} C
\textbf{Giải thích:} Hàm băm chia lấy dư $h(k) = k \bmod m$ đạt hiệu quả tốt nhất khi kích thước bảng $m$ được chọn là một số nguyên tố.
\end{question}

\begin{question}{8.27}
% [ID gốc: 5.35]
% [Nguồn: Ngân hàng CSD203 (Câu #114)]
Which of the following algorithms has a logarithmic runtime complexity?
\begin{choices}
\choice Binary search
\choice Linear search
\choice Merge sort
\choice Selection sort
\end{choices}
\textbf{Đáp án:} A
\textbf{Giải thích:} Duyệt Inorder trên cây nhị phân bình thường (không phải BST) không nhất thiết cho ra dãy tăng dần, thứ tự phụ thuộc vào cấu trúc cây.
\end{question}

\begin{question}{8.28}
% [ID gốc: 5.37]
% [Nguồn: Ngân hàng CSD203 (Câu #116)]
The complexity of Binary search algorithm is
\begin{choices}
\choice $O(n)$
\choice $O(\log n)$
\choice $O(n2)$
\choice $O(n\log n)$
\end{choices}
\textbf{Đáp án:} B
\textbf{Giải thích:} Danh sách liên kết đôi cho phép duyệt hai chiều (tiến và lùi) nhờ có cả hai con trỏ \texttt{next} và \texttt{prev} tại mỗi nút.
\end{question}

\begin{question}{8.29}
% [ID gốc: 5.40]
% [Nguồn: Ngân hàng CSD203 (Câu #120)]
The average number of key comparisons done in a successful linear search in an array of length n is
\begin{choices}
\choice $n$
\choice $n/2$
\choice $n * n$
\choice $(n + 1)/2$
\end{choices}
\textbf{Đáp án:} D
\textbf{Giải thích:} Cây nhị phân có chiều cao $h$ có số nút tối đa là $2^{h+1} - 1$ khi cây là cây nhị phân hoàn hảo (Perfect binary tree).
\end{question}

\begin{question}{8.30}
% [ID gốc: 5.41]
% [Nguồn: Ngân hàng CSD203 (Câu #121)]
What is recursion used for?
\begin{choices}
\choice Recursion is used to create loops in languages where other loops are not available.
\choice We use recursion only to implement mathematical formulas in code.
\choice Recursion is used to iterate through sequences of files and directories.
\choice Recursion lets us tackle complex problems by reducing the problem to a simpler one.
\end{choices}
\textbf{Đáp án:} D
\textbf{Giải thích:} Khử đệ quy bằng ngăn xếp giúp tránh lỗi tràn bộ nhớ call stack khi độ sâu đệ quy quá lớn.
\end{question}

\begin{question}{8.31}
% [ID gốc: 5.42]
% [Nguồn: Ngân hàng CSD203 (Câu #124)]
What are the advantages of recursive programming over iterative programming?
\begin{choices}
\choice Recursion provides a clean and simple way to write code
\choice The recursive program has greater space requirements than the iterative program
\choice The recursive program has greater time requirementsthan the iterative program
\choice None of the others
\end{choices}
\textbf{Đáp án:} A
\textbf{Giải thích:} Cây đỏ đen đảm bảo đường đi dài nhất từ gốc đến lá không quá 2 lần đường đi ngắn nhất, giữ chiều cao cây luôn ở mức $O(\log n)$.
\end{question}

\begin{question}{8.32}
% [ID gốc: 5.43]
% [Nguồn: Ngân hàng CSD203 (Câu #142)]
The brute-force pattern matching algorithm compares the pattern P containing n characters with the text T containing m characters. What is the time complexity of it?
\begin{choices}
\choice $O(n * m)$
\choice $O(n + m)$
\choice $O(n)$
\choice $O(m)$
\end{choices}
\textbf{Đáp án:} A
\textbf{Giải thích:} Trong danh sách liên kết đôi, thao tác xóa một nút khi đã có con trỏ trỏ đến nút đó chỉ mất thời gian $O(1)$.
\end{question}

\begin{question}{8.33}
% [ID gốc: 5.44]
% [Nguồn: Ngân hàng CSD203 (Câu #144)]
Given the length of input data is 8 and the length of output data is 4. What is the data compression rate in this case?
\begin{choices}
\choice 8
\choice 4
\choice 1
\choice 0.5
\end{choices}
\textbf{Đáp án:} D
\textbf{Giải thích:} Tỷ lệ nén dữ liệu = 4 / 8 = 0.5 (hoặc 50%), thể hiện kích thước dữ liệu giảm đi một nửa.
\end{question}

\begin{question}{8.34}
% [ID gốc: 5.45]
% [Nguồn: Ngân hàng CSD203 (Câu #145)]
Best Case Time Complexity of Linear Search is
\begin{choices}
\choice $O(1)$
\choice $O(n)$
\choice $O(n * n)$
\choice $O(nlogn)$
\end{choices}
\textbf{Đáp án:} A
\textbf{Giải thích:} Đồ thị hai phía (Bipartite graph) là đồ thị có thể phân chia tập đỉnh thành hai tập rời rạc sao cho mọi cạnh chỉ nối giữa hai tập này.
\end{question}

\begin{question}{8.35}
% [ID gốc: 5.46]
% [Nguồn: Ngân hàng CSD203 (Câu #146)]
The worst case occurs in linear search algorithm when
\begin{choices}
\choice Item is somewhere in the middle of the array
\choice Item is not in the array at all
\choice Item is the last element in the array
\choice Item is the last element in the array or is not there at all
\end{choices}
\textbf{Đáp án:} D
\textbf{Giải thích:} Thuật toán sắp xếp Shell Sort là một cải tiến của Insertion Sort, so sánh các phần tử cách nhau một khoảng cách (gap) giảm dần.
\end{question}

\begin{question}{8.36}
% [ID gốc: 5.47]
% [Nguồn: Ngân hàng CSD203 (Câu #147)]
The complexity of Linear search algorithm is
\begin{choices}
\choice $O(n)$
\choice $O(\log n)$
\choice $O(n^2)$
\choice $O(n \log n)$
\end{choices}
\textbf{Đáp án:} A
\textbf{Giải thích:} Phương thức \texttt{get()} trên hàng đợi \texttt{Queue} lấy ra phần tử đầu tiên và chặn luồng nếu hàng đợi đang rỗng (nếu để mặc định).
\end{question}

\begin{question}{8.37}
% [ID gốc: 5.48]
% [Nguồn: Ngân hàng CSD203 (Câu #148)]
Which of the following cases occurs when searching an array using linear search the value to be searched is equal to the last element of the array?
\begin{choices}
\choice The best case
\choice The worst case
\choice The average case
\choice The amortized case
\end{choices}
\textbf{Đáp án:} B
\textbf{Giải thích:} Trong cây Heap cực tiểu (Min-Heap), phần tử có giá trị nhỏ nhất luôn nằm tại nút gốc.
\end{question}

\begin{question}{8.38}
% [ID gốc: 5.49]
% [Nguồn: Ngân hàng CSD203 (Câu #149)]
Given a definition as follows: ``A unique type of recursion where the last procedure of a function is a recursive call'' What is this?
\begin{choices}
\choice Tail-Recursion
\choice NonTail-Recursion
\choice Recursion
\choice None of these
\end{choices}
\textbf{Đáp án:} A
\textbf{Giải thích:} Duyệt DFS trên cây tương đương với duyệt tiền thứ tự (Preorder) khi cây được xem như một đồ thị có hướng.
\end{question}

\begin{question}{8.39}
% [ID gốc: 5.50]
% [Nguồn: Ngân hàng CSD203 (Câu #150)]
By default, what is the maximum depth of recursion?
\begin{choices}
\choice 10000
\choice 1000
\choice 100
\choice 10
\end{choices}
\textbf{Đáp án:} B
\textbf{Giải thích:} Bảng băm có dung lượng 10 với địa chỉ mở: khi xảy ra xung đột, vị trí mới được dò tuần tự đến khi gặp ô trống.
\end{question}

\begin{question}{8.40}
% [ID gốc: 5.51]
% [Nguồn: Ngân hàng CSD203 (Câu #151)]
Given an array as follows: [1, 3, 4, 5, 6, 8, 9] What is the number of key comparisons done to find 5 using binary search?
\begin{choices}
\choice 1
\choice 2
\choice 3
\choice 4
\end{choices}
\textbf{Đáp án:} A
\textbf{Giải thích:} Độ phức tạp bộ nhớ của thuật toán QuickSort phụ thuộc vào độ sâu ngăn xếp đệ quy, tối ưu là $O(\log n)$ và xấu nhất là $O(n)$.
\end{question}

\begin{question}{8.41}
% [ID gốc: 5.52]
% [Nguồn: Ngân hàng CSD203 (Câu #152)]
What will be the output of the following code?
\begin{codebox}
def f(n):
    if(n<=3):
        return 1
    else:
        return f(n-3) + f(n-1)
print(f(5))
\end{codebox}
\begin{choices}
\choice 5
\choice 4
\choice 3
\choice 2
\end{choices}
\textbf{Đáp án:} C
\textbf{Giải thích:} Hàm đệ quy $f(n)$ với điều kiện dừng $n \le 3$: truy vết từng bước đệ quy cho ra kết quả tương ứng phương án đúng.
\end{question}

\begin{question}{8.42}
% [ID gốc: 5.53]
% [Nguồn: Ngân hàng CSD203 (Câu #153)]
What is advantage of Recursion?
\begin{choices}
\choice Recursive functions make the code look clean and elegant
\choice Recursive calls are inefficient as they take up less of memory
\choice Recursive calls are inefficient as they take up less of time
\choice None of the others
\end{choices}
\textbf{Đáp án:} A
\textbf{Giải thích:} Biểu diễn đồ thị bằng ma trận kề tốn bộ nhớ $O(V^2)$, phù hợp và hiệu quả khi đồ thị là đồ thị dày đặc (Dense graph).
\end{question}

\begin{question}{8.43}
% [ID gốc: 5.54]
% [Nguồn: Ngân hàng CSD203 (Câu #154)]
What will be the output of the following code?
\begin{codebox}
def func(n):
    if (n > 0):
        func(int(n/2))
        print(n%2)

func(11)
\end{codebox}
\begin{choices}
\choice 11
\choice 1011
\choice 1101
\choice 0
\end{choices}
\textbf{Đáp án:} B
\textbf{Giải thích:} Danh sách liên kết đơn không cho phép truy xuất ngẫu nhiên (random access) với thời gian $O(1)$ như mảng.
\end{question}

\begin{question}{8.44}
% [ID gốc: 5.55]
% [Nguồn: Ngân hàng CSD203 (Câu #161)]
What is the time complexity to insert a new node at the end of a singly linked list if the tail pointer is not maintained?
\begin{choices}
\choice $O(1)$
\choice $O(n)$
\choice $O(\log n)$
\choice $O(n \log n)$
\end{choices}
\textbf{Đáp án:} B
\textbf{Giải thích:} Thời gian xây dựng Binary Heap từ mảng $n$ phần tử không có thứ tự bằng thuật toán Bottom-Up Heapify là $O(n)$.
\end{question}

\begin{question}{8.45}
% [ID gốc: 5.56]
% [Nguồn: Ngân hàng CSD203 (Câu #162)]
Consider the function below that operates on a singly linked list with n nodes:
\begin{codebox}
class Node:
    def __init__(self, data):
        self.data = data
        self.next = None

def fn(head):
    slow = head
    fast = head
    while fast is not None and fast.next is not None:
        slow = slow.next
        fast = fast.next.next
    return slow.data
\end{codebox}
What does the \texttt{fn()} function do, and what is its time complexity?
\begin{choices}
\choice It finds the middle node of the linked list in $O(n)$ time.
\choice It finds the middle node of the linked list in $O(\log n)$ time.
\choice It finds the length of the linked list in $O(n)$ time.
\choice It finds the middle node of the linked list in $O(n^2)$ time.
\end{choices}
\textbf{Đáp án:} A
\textbf{Giải thích:} Trong danh sách liên kết vòng kép, con trỏ \texttt{prev} của nút đầu trỏ đến nút cuối và con trỏ \texttt{next} của nút cuối trỏ về nút đầu.
\end{question}

\begin{question}{8.46}
% [ID gốc: 5.57]
% [Nguồn: Ngân hàng CSD203 (Câu #174)]
Given a string S = ``12312345'' and a pattern P = ``123''. What is KMP algorithm returns? (the string is indexed from 0)
\begin{choices}
\choice 0
\choice 3
\choice 1,4
\choice 0,3
\end{choices}
\textbf{Đáp án:} D
\textbf{Giải thích:} Thuật toán KMP tìm thấy chuỗi mẫu '123' trong chuỗi '12312345' tại hai vị trí bắt đầu là 0 và 3.
\end{question}

\begin{question}{8.47}
% [ID gốc: 5.59]
% [Nguồn: Ngân hàng CSD203 (Câu #179)]
Given a function as follows:
\begin{codebox}
def rec(n):
    if (n == 0):
        return n
    else:
        return n + rec(n-1)
\end{codebox}
What is statement ``return n + rec(n-1)'' called in recursive algorithm?
\begin{choices}
\choice Base case
\choice Recursive calls
\choice Anchor case
\choice Nothing
\end{choices}
\textbf{Đáp án:} B
\textbf{Giải thích:} Lệnh \texttt{return n + rec(n-1)} trong hàm đệ quy là lời gọi đệ quy (Recursive call) thực hiện bước quy nạp tiến tới bài toán con.
\end{question}

\begin{question}{8.48}
% [ID gốc: 5.61]
% [Nguồn: Ngân hàng CSD203 (Câu #181)]
Which search algorithm is best for a large ordered list?
\begin{choices}
\choice A for each loop
\choice Sequential search
\choice Binary search
\choice None of these
\end{choices}
\textbf{Đáp án:} C
\textbf{Giải thích:} Đối với danh sách lớn đã được sắp xếp, thuật toán tìm kiếm tối ưu nhất là tìm kiếm nhị phân (Binary search) với độ phức tạp $O(\log n)$.
\end{question}

\begin{question}{8.49}
% [ID gốc: 5.62]
% [Nguồn: Ngân hàng CSD203 (Câu #182)]
What will be the output of the following code?
\begin{codebox}
def func(n):
    if n == 0:
        return n
    return n + func(n-1)

print(func(4))
\end{codebox}
\begin{choices}
\choice 4
\choice 1
\choice 10
\choice 6
\end{choices}
\textbf{Đáp án:} C
\textbf{Giải thích:} Hàm \texttt{func(4)} tính tổng đệ quy: $4 + 3 + 2 + 1 + 0 = 10$.
\end{question}

\begin{question}{8.50}
% [ID gốc: 5.63]
% [Nguồn: Ngân hàng CSD203 (Câu #183)]
What will be the output of the following code?
\begin{codebox}
def fact(n):
    if n == 0:
        return "good"
    else:
        return fact(n-4)

print(fact(16))
\end{codebox}
\begin{choices}
\choice good good good good
\choice good
\choice 16
\choice 4
\end{choices}
\textbf{Đáp án:} B
\textbf{Giải thích:} Hàm \texttt{fact(16)} gọi $fact(12) \to fact(8) \to fact(4) \to fact(0)$ chạm điều kiện cơ sở trả về chuỗi 'good'.
\end{question}

\begin{question}{8.51}
% [ID gốc: 5.64]
% [Nguồn: Ngân hàng CSD203 (Câu #184)]
What will be the output of the following code?
\begin{codebox}
def rec(n):
    if (n == 0):
        return n
    else:
        return n + rec(n-1)

print(rec(4))
\end{codebox}
\begin{choices}
\choice 4
\choice 0
\choice 10
\choice Error
\end{choices}
\textbf{Đáp án:} C
\textbf{Giải thích:} Hàm \texttt{rec(4)} tính tổng đệ quy các số từ 1 đến 4: $4 + 3 + 2 + 1 + 0 = 10$.
\end{question}

\begin{question}{8.52}
% [ID gốc: 5.65]
% [Nguồn: Ngân hàng CSD203 (Câu #185)]
Given an array as follows: [6, 1, 5, 3, 7]. What is the number of key comparisons done to check if 9 is in the array or not by using linear search?
\begin{choices}
\choice 0
\choice 1
\choice 5
\choice 6
\end{choices}
\textbf{Đáp án:} C
\textbf{Giải thích:} Tìm kiếm tuyến tính giá trị 9 (không tồn tại) trong mảng 5 phần tử [6, 1, 5, 3, 7] phải so sánh lần lượt với tất cả 5 phần tử trong mảng.
\end{question}

\begin{question}{8.53}
% [ID gốc: 5.66]
% [Nguồn: Ngân hàng CSD203 (Câu #188)]
What will be the output of the following code?
\begin{codebox}
def func(n):
    if (n > 0):
        print(n%2)
        func(int(n/2))

func(11)
\end{codebox}
\begin{choices}
\choice 11
\choice 1011
\choice 1101
\choice 0
\end{choices}
\textbf{Đáp án:} C
\textbf{Giải thích:} Hàm \texttt{func(11)} in phần dư nhị phân: $11\%2=1, 5\%2=1, 2\%2=0, 1\%2=1$. Kết quả in ra theo thứ tự gọi là 1101.
\end{question}

\begin{question}{8.54}
% [ID gốc: 5.67]
% [Nguồn: Ngân hàng CSD203 (Câu #190)]
Given a string S = ``ABCDEFABCD'' and a pattern P = ``CD''. What is KMP algorithm returns? (the string is indexed from 0)
\begin{choices}
\choice 2
\choice 3
\choice 1, 7
\choice 2, 8
\end{choices}
\textbf{Đáp án:} D
\textbf{Giải thích:} Mẫu 'CD' xuất hiện trong chuỗi 'ABCDEFABCD' tại các vị trí chỉ mục là 2 và 8.
\end{question}

\begin{question}{8.55}
% [ID gốc: 5.70]
% [Nguồn: Ngân hàng CSD203 (Câu #195)]
A recursive function is
\begin{choices}
\choice A function that calls other functions
\choice A function that calls itself
\choice A function that has arguments
\choice None of these
\end{choices}
\textbf{Đáp án:} B
\textbf{Giải thích:} Hàm đệ quy (Recursive function) được định nghĩa là một hàm tự gọi lại chính nó trong quá trình thực thi.
\end{question}

\begin{question}{8.56}
% [ID gốc: 5.71]
% [Nguồn: Ngân hàng CSD203 (Câu #196)]
Given the following code
\begin{codebox}
def f(n):
    if(n==1):
        return 1
    return n*f(n-1)
\end{codebox}
The recursion ends when the number reduces to 1. This is called?
\begin{choices}
\choice The base case
\choice The conductive case
\choice Function call
\choice Nothing
\end{choices}
\textbf{Đáp án:} A
\textbf{Giải thích:} Điều kiện dừng khi số $n$ giảm về 1 trong hàm đệ quy để kết thúc việc gọi tiếp được gọi là trường hợp cơ sở (The base case).
\end{question}

\begin{question}{8.57}
% [ID gốc: 5.73]
% [Nguồn: Ngân hàng CSD203 (Câu #205)]
What is the time complexity for access an element in Arrays?
\begin{choices}
\choice $O(n)$
\choice $O(1)$
\choice $O(n * n)$
\choice $O(n + 1)$
\end{choices}
\textbf{Đáp án:} B
\textbf{Giải thích:} Trong cấu trúc mảng (Array), các phần tử được lưu trữ ở các ô nhớ liên tiếp, cho phép truy xuất ngẫu nhiên bất kỳ phần tử nào theo chỉ số trong thời gian hằng số $O(1)$.
\end{question}

\begin{question}{8.58}
% [ID gốc: 5.74]
% [Nguồn: Ngân hàng CSD203 (Câu #210)]
Given a function as follows:
\begin{codebox}
def rec(n):
    if (n == 0):
        return n
    else:
        return n + rec(n-1)
\end{codebox}
What is statement "if (n == 0):" called in recursive algorithm?
\begin{choices}
\choice Base case
\choice Recursive calls
\choice Recall case
\choice Nothing
\end{choices}
\textbf{Đáp án:} A
\textbf{Giải thích:} Trong thuật toán đệ quy, câu lệnh kiểm tra điều kiện dừng như \texttt{if (n == 0):} được gọi là trường hợp cơ sở (Base case), ngăn chặn đệ quy vô hạn.
\end{question}

\begin{question}{8.59}
% [ID gốc: 5.75]
% [Nguồn: Ngân hàng CSD203 (Câu #219)]
What is the time complexity of searching for an element in a linked list with 'n' elements in the worst case?
\begin{choices}
\choice $O(n)$
\choice $O(\log n)$
\choice $O(1)$
\choice $O(n^2)$
\end{choices}
\textbf{Đáp án:} A
\textbf{Giải thích:} Trong danh sách liên kết có $n$ phần tử, trường hợp xấu nhất phần tử cần tìm nằm ở cuối danh sách hoặc không tồn tại, thuật toán phải duyệt qua toàn bộ $n$ nút, có độ phức tạp là $O(n)$.
\end{question}

\begin{question}{8.60}
% [ID gốc: 5.76]
% [Nguồn: Ngân hàng CSD203 (Câu #234)]
What is the main advantage of using Huffman coding over fixed-length coding?
\begin{choices}
\choice Easier to implement
\choice Uses less memory
\choice Faster encoding and decoding
\choice Reduces the average code length, saving space
\end{choices}
\textbf{Đáp án:} D
\textbf{Giải thích:} Ưu điểm chính của mã hóa Huffman so với mã hóa độ dài cố định là gán từ mã ngắn hơn cho các ký tự xuất hiện thường xuyên, làm giảm độ dài mã trung bình và tiết kiệm dung lượng lưu trữ.
\end{question}

\begin{question}{8.61}
% [ID gốc: 5.79]
% [Nguồn: Ngân hàng CSD203 (Câu #238)]
A sorted array contains 8 items. Using binary search, the maximum number of comparisons to search for an item in this array is ..........
\begin{choices}
\choice 1
\choice 2
\choice 3
\choice 4
\end{choices}
\textbf{Đáp án:} D
\textbf{Giải thích:} Tìm kiếm nhị phân trên mảng 8 phần tử: số phép so sánh tối đa trong trường hợp xấu nhất là $\lfloor\log_2 8\rfloor + 1 = 3 + 1 = 4$ phép so sánh.
\end{question}

\begin{question}{8.62}
% [ID gốc: 5.81]
% [Nguồn: Ngân hàng CSD203 (Câu #240)]
Which of the following problems is least likely to benefit from recursion?
\begin{choices}
\choice Tower of Hanoi
\choice Calculating Fibonacci numbers
\choice Summing elements of a list
\choice Traversing a directory structure
\end{choices}
\textbf{Đáp án:} C
\textbf{Giải thích:} Bài toán tính tổng các phần tử trong danh sách (Summing elements of a list) là bài toán lặp tuyến tính đơn giản, việc dùng đệ quy gây lãng phí bộ nhớ call stack mà không mang lại lợi ích về cấu trúc giải thuật.
\end{question}

\begin{question}{8.63}
% [ID gốc: 5.82]
% [Nguồn: Ngân hàng CSD203 (Câu #241)]
Which technique can be used to eliminate tail recursion?
\begin{choices}
\choice Dynamic programming
\choice Memoization
\choice Iteration
\choice Breadth-first search
\end{choices}
\textbf{Đáp án:} C
\textbf{Giải thích:} Đệ quy đuôi (Tail recursion) có thể được loại bỏ và tối ưu hóa chuyển thành dạng vòng lặp (Iteration) nhằm tiết kiệm bộ nhớ ngăn xếp lời gọi hàm.
\end{question}

\begin{question}{8.64}
% [ID gốc: 5.83]
% [Nguồn: Ngân hàng CSD203 (Câu #267)]
What is the main objective of Huffman Coding?
\begin{choices}
\choice Compression
\choice Encryption
\choice Encrypted Compression
\choice Data Integrity
\end{choices}
\textbf{Đáp án:} A
\textbf{Giải thích:} Mục tiêu chính của mã hóa Huffman (Huffman Coding) là nén dữ liệu (Data compression) không mất mát thông tin.
\end{question}

\begin{question}{8.65}
% [ID gốc: 5.84]
% [Nguồn: Ngân hàng CSD203 (Câu #268)]
When should we use the Huffman Coding?
\begin{choices}
\choice For compression
\choice For encoding data
\choice For decoding data
\choice For searching data
\end{choices}
\textbf{Đáp án:} A
\textbf{Giải thích:} Mã hóa Huffman được ứng dụng trong kỹ thuật nén dữ liệu (Compression) để giảm dung lượng lưu trữ và băng thông truyền tải.
\end{question}

\begin{question}{8.66}
% [ID gốc: 5.87]
% [Nguồn: Ngân hàng CSD203 (Câu #273)]
Given a sorted array as follows: [2, 4, 5, 6, 8] and a search key x using binary search? The Best case of searching occurs when value of x is
\begin{choices}
\choice 2
\choice 5
\choice 8
\choice None of these
\end{choices}
\textbf{Đáp án:} B
\textbf{Giải thích:} Trường hợp tốt nhất của tìm kiếm nhị phân xảy ra khi phần tử cần tìm nằm ngay tại vị trí chính giữa mảng. Với mảng 5 phần tử \texttt{[2, 4, 5, 6, 8]}, vị trí chính giữa là phần tử 5.
\end{question}

\begin{question}{8.67}
% [ID gốc: 5.88]
% [Nguồn: Ngân hàng CSD203 (Câu #274)]
What will be the output of the following recursive function when n = 10?
\begin{codebox}
def func(n):
    if n < 2:
        return 0
    if n % 2 == 0:
        return 1 + func(n - 1)
    return func(n - 1)

print(func(10))
\end{codebox}
\begin{choices}
\choice 4
\choice 5
\choice 6
\choice 7
\end{choices}
\textbf{Đáp án:} B
\textbf{Giải thích:} Hàm đệ quy \texttt{func(10)} cộng 1 cho mỗi số chẵn trong khoảng từ 10 xuống 2 (gồm các số 10, 8, 6, 4, 2), tổng cộng có 5 số chẵn, kết quả trả về là 5.
\end{question}

\begin{question}{8.68}
% [ID gốc: 5.90]
% [Nguồn: Ngân hàng CSD203 (Câu #276)]
In recursive programmes, what happens if the base condition is not defined?
\begin{choices}
\choice It runs into infinite loop
\choice Exception is thrown
\choice No problem
\choice None of these
\end{choices}
\textbf{Đáp án:} A
\textbf{Giải thích:} Trong chương trình đệ quy, nếu không định nghĩa điều kiện cơ sở (base condition), hàm sẽ gọi vô tận dẫn đến vòng lặp vô hạn và tràn bộ nhớ ngăn xếp (Stack Overflow / RecursionError).
\end{question}

\begin{question}{8.69}
% [ID gốc: 5.92]
% [Nguồn: Đề thi FE/RE CSD203 (Câu #6)]
% Quiz: CSD203_SU26_FE - Question #6 (ID: 60089e31-637f-46a8-8e66-7a3d9029870c)
What does KMP algorithm solve in text processing?
\begin{choices}
\choice Exactly matching
\choice Regex matching
\choice Lossy compression
\choice Lossless compression
\end{choices}
\textbf{Đáp án:} A
\textbf{Giải thích:} Thuật toán Knuth-Morris-Pratt (KMP) được sử dụng để giải quyết bài toán tìm kiếm chuỗi khớp chính xác (exact string matching) với thời gian tuyến tính.
\end{question}

\begin{question}{8.70}
% [ID gốc: 5.93]
% [Nguồn: Đề thi FE/RE CSD203 (Câu #7)]
% Quiz: CSD203_SU26_FE - Question #7 (ID: 01b6a5a1-7742-4f56-9090-97cabd27bc07)
What is the time complexity for inserting a new node at the beginning of a singly linked list?
\begin{choices}
\choice $O(1)$
\choice $O(n)$
\choice $O(\log n)$
\choice $O(n \log n)$
\end{choices}
\textbf{Đáp án:} A
\textbf{Giải thích:} Việc chèn một nút mới vào đầu danh sách liên kết đơn chỉ cần thay đổi con trỏ head, thực hiện trong thời gian hằng số $O(1)$.
\end{question}

\begin{question}{8.71}
% [ID gốc: 5.94]
% [Nguồn: Đề thi FE/RE CSD203 (Câu #9)]
% Quiz: CSD203_SU26_FE - Question #9 (ID: cdaa0648-be21-4822-b2b3-6d6003677989)
Which part of the KMP algorithm ensures that the time complexity remains linear?
\begin{choices}
\choice The use of binary search
\choice The construction of the prefix table
\choice The recursive nature of the algorithm
\choice The preprocessing of the text
\end{choices}
\textbf{Đáp án:} B
\textbf{Giải thích:} Bảng tiền tố (prefix table hay mảng LPS) giúp thuật toán KMP bỏ qua các bước so sánh thừa khi xảy ra bất khớp, đảm bảo độ phức tạp thời gian đạt mức tuyến tính O(n + m).
\end{question}

\begin{question}{8.72}
% [ID gốc: 5.95]
% [Nguồn: Đề thi FE/RE CSD203 (Câu #15)]
% Quiz: CSD203_SU26_FE - Question #15 (ID: 92d65e9e-92bb-4ff3-8658-0af4a0c9a356)
How does the abundance of digitized text affect natural language processing tasks?
\begin{choices}
\choice It hinders language understanding
\choice It speeds up text analysis tasks
\choice It has no impact on language processing
\choice It reduces the need for text preprocessing
\end{choices}
\textbf{Đáp án:} B
\textbf{Giải thích:} Sự phong phú của dữ liệu văn bản số hóa cung cấp nguồn tài nguyên lớn giúp huấn luyện mô hình hiệu quả và đẩy nhanh các tác vụ phân tích ngôn ngữ tự nhiên.
\end{question}

\begin{question}{8.73}
% [ID gốc: 5.96]
% [Nguồn: Đề thi FE/RE CSD203 (Câu #17)]
% Quiz: CSD203_SU26_FE - Question #17 (ID: 2a0da93c-b8a5-4b5c-a68b-911ad6eae999)
If Algorithm A has a time complexity of $O(n^2)$ and Algorithm B has a time complexity of $O(n \log n)$, which algorithm is more efficient for large input sizes?
\begin{choices}
\choice Algorithm B
\choice Algorithm A
\choice It depends on the specific input
\choice Both are equally efficient
\end{choices}
\textbf{Đáp án:} A
\textbf{Giải thích:} Hàm n log n tăng chậm hơn nhiều so với n^2 khi kích thước đầu vào n trở nên rất lớn, do đó Algorithm B hiệu quả hơn (tương ứng với phương án A).
\end{question}

\begin{question}{8.74}
% [ID gốc: 5.98]
% [Nguồn: Đề thi FE/RE CSD203 (Câu #28)]
% Quiz: CSD203_SU26_FE - Question #28 (ID: e5c5ad93-c04f-4f2e-86af-0d9ba40b66ce)
What will be the output of the following recursive function when n = 6?
\begin{codebox}
def func(n):
    if n <= 0:
        return 0
    if n % 2 == 0:
        return n + func(n - 2)
    return func(n - 1)

print(func(6))
\end{codebox}
\begin{choices}
\choice 12
\choice 16
\choice 18
\choice 20
\end{choices}
\textbf{Đáp án:} A
\textbf{Giải thích:} Hàm cộng dồn các số chẵn giảm dần từ 6: 6 + func(4) → 6 + 4 + func(2) → 6 + 4 + 2 + func(0) = 12.
\end{question}

\begin{question}{8.75}
% [ID gốc: 5.101]
% [Nguồn: Đề thi FE/RE CSD203 (Câu #33)]
% Quiz: CSD203_SU26_FE - Question #33 (ID: fd46f234-a029-4ebf-a7bb-4f1af078bc86)
What will be the output of the following code?
\begin{codebox}
def rec(n):
    if n <= 1:
        return n
    else:
        return rec(n - 2) + rec(n - 1)

print(rec(4))
\end{codebox}
\begin{choices}
\choice 5
\choice 4
\choice 3
\choice Error
\end{choices}
\textbf{Đáp án:} C
\textbf{Giải thích:} Đây là hàm tính số Fibonacci thứ n. Với n = 4: rec(4) = rec(2) + rec(3) = 1 + 2 = 3.
\end{question}

\begin{question}{8.76}
% [ID gốc: 5.102]
% [Nguồn: Đề thi FE/RE CSD203 (Câu #35)]
% Quiz: CSD203_SU26_FE - Question #35 (ID: b1e31dd5-f2dc-4b43-94fa-45ca1454b38f)
Which data structure is commonly used in eliminating tail recursion?
\begin{choices}
\choice Stack
\choice Queue
\choice Linked List
\choice Array
\end{choices}
\textbf{Đáp án:} A
\textbf{Giải thích:} Ngăn xếp (Stack) thường được sử dụng để khử đệ quy bằng cách lưu trữ thủ công các trạng thái hàm thay vì dùng ngăn xếp gọi của hệ thống.
\end{question}

\begin{question}{8.77}
% [ID gốc: 5.103]
% [Nguồn: Đề thi FE/RE CSD203 (Câu #36)]
% Quiz: CSD203_SU26_FE - Question #36 (ID: 88aa87e6-fbbf-44af-a0c0-39c6dbd5c90d)
What is a disadvantage of recursion?
\begin{choices}
\choice Recursive functions make the code look clean and elegant
\choice Recursive calls are inefficient as they take up a lot of memory and time
\choice A complex task can be broken down into simpler sub-problems using recursion
\choice None of the others
\end{choices}
\textbf{Đáp án:} B
\textbf{Giải thích:} Nhược điểm lớn của đệ quy là tiêu tốn nhiều bộ nhớ stack cho các lời gọi hàm và có thể gây chậm thời gian thực thi nếu không tối ưu.
\end{question}

\begin{question}{8.78}
% [ID gốc: 5.104]
% [Nguồn: Đề thi FE/RE CSD203 (Câu #38)]
% Quiz: CSD203_SU26_FE - Question #38 (ID: 15806e19-7b07-497d-b2ea-406dc0321374)
What is the condition for applying binary search in an array?
\begin{choices}
\choice The array should not be too long
\choice The array should not be too short
\choice The array should be sorted
\choice The array should not be sorted
\end{choices}
\textbf{Đáp án:} C
\textbf{Giải thích:} Điều kiện tiên quyết để áp dụng thuật toán tìm kiếm nhị phân là mảng dữ liệu phải được sắp xếp từ trước.
\end{question}

\begin{question}{8.79}
% [ID gốc: 5.105]
% [Nguồn: Đề thi FE/RE CSD203 (Câu #40)]
% Quiz: CSD203_SU26_FE - Question #40 (ID: a9584292-ed1a-4e5a-a515-63145c55bca5)
Which array operation takes $O(1)$ time complexity?
\begin{choices}
\choice Inserting at the front
\choice Removing an element
\choice Indexing an element
\choice Growing the array
\end{choices}
\textbf{Đáp án:} C
\textbf{Giải thích:} Thao tác truy xuất phần tử thông qua chỉ mục (indexing) trong mảng mất thời gian hằng số $O(1)$ vì vị trí ô nhớ được tính toán trực tiếp.
\end{question}

\begin{question}{8.80}
% [ID gốc: 5.106]
% [Nguồn: Đề thi FE/RE CSD203 (Câu #44)]
% Quiz: CSD203_SU26_FE - Question #44 (ID: 5b0ea03f-2acc-46ce-8487-3de1f5d40b81)
Given the following code, the recursion ends when the number reduces to 1. This is called?
\begin{codebox}
def f(n):
    if n == 1:
        return 1
    return n * f(n - 1)
\end{codebox}
\begin{choices}
\choice The base case
\choice The conductive case
\choice Function call
\choice Nothing
\end{choices}
\textbf{Đáp án:} A
\textbf{Giải thích:} Điều kiện dừng của hàm đệ quy để ngăn chặn việc gọi vô hạn được gọi là trường hợp cơ sở (base case).
\end{question}

\begin{question}{8.81}
% [ID gốc: 5.108]
% [Nguồn: Đề thi FE/RE CSD203 (Câu #51)]
% Quiz: CSD203_SU26_RE - Question #1 (ID: b5304cc4-45d8-41fe-9352-b828b5691cfa)
The average number of key comparisons done in a successful linear search in an array of length n is _____
\begin{choices}
\choice n
\choice n/2
\choice n * n
\choice (n + 1) / 2
\end{choices}
\textbf{Đáp án:} D
\textbf{Giải thích:} Trong tìm kiếm tuyến tính (linear search) thành công, nếu phần tử nằm ở vị trí i thì số lần so sánh là i. Trung bình số lần so sánh là (1 + 2 + ... + n) / n = [n(n + 1) / 2] / n = (n + 1) / 2.
\end{question}

\begin{question}{8.82}
% [ID gốc: 5.109]
% [Nguồn: Đề thi FE/RE CSD203 (Câu #53)]
% Quiz: CSD203_SU26_RE - Question #3 (ID: a9efc104-2f76-4bc5-bc94-6816601e3a2e)
What will be the output of the following code?
\begin{codebox}
def func(n):
    if n > 0:
        print(n % 2)
        func(int(n / 2))

func(11)
\end{codebox}
\begin{choices}
\choice 1
\choice 1011
\choice 1101
\choice 0
\end{choices}
\textbf{Đáp án:} C
\textbf{Giải thích:} Hàm in phần dư n \% 2 trước rồi mới gọi đệ quy func(int(n / 2)):
func(11) in 11 \% 2 = 1, gọi func(5)
func(5) in 5 \% 2 = 1, gọi func(2)
func(2) in 2 \% 2 = 0, gọi func(1)
func(1) in 1 \% 2 = 1, gọi func(0)
func(0) kết thúc.
Thứ tự in ra là 1, 1, 0, 1 tương ứng chuỗi 1101.
\end{question}

\begin{question}{8.83}
% [ID gốc: 5.110]
% [Nguồn: Đề thi FE/RE CSD203 (Câu #55)]
% Quiz: CSD203_SU26_RE - Question #5 (ID: e83ba0b0-aa62-4789-b021-5e144d951d23)
What will be the output of the following code?
\begin{codebox}
def func(n):
    if n > 0:
        func(int(n / 2))
        print(n % 2)

func(11)
\end{codebox}
\begin{choices}
\choice 1
\choice 1011
\choice 1101
\choice 0
\end{choices}
\textbf{Đáp án:} B
\textbf{Giải thích:} Hàm gọi đệ quy trước rồi mới in n \% 2, cho ra biểu diễn nhị phân của 11. Chuỗi gọi func(11) → func(5) → func(2) → func(1) → func(0), sau đó in ngược lên: func(1) in 1, func(2) in 0, func(5) in 1, func(11) in 1. Thứ tự in ra là 1, 0, 1, 1 tương ứng chuỗi 1011.
\end{question}

\begin{question}{8.84}
% [ID gốc: 5.112]
% [Nguồn: Đề thi FE/RE CSD203 (Câu #69)]
% Quiz: CSD203_SU26_RE - Question #19 (ID: 80125fad-37a8-45b7-b5f9-c416bf501263)
Given an array as follows: [6, 1, 5, 3, 7]. What is the number of key comparisons done to check if 9 is in the array or not by using linear search?
\begin{choices}
\choice 0
\choice 1
\choice 5
\choice 6
\end{choices}
\textbf{Đáp án:} C
\textbf{Giải thích:} Khi tìm kiếm tuyến tính giá trị 9 không tồn tại trong mảng 5 phần tử, thuật toán phải so sánh với lần lượt từng phần tử từ đầu đến cuối. Số lần so sánh đúng bằng số phần tử của mảng là 5.
\end{question}

\begin{question}{8.85}
% [ID gốc: 5.113]
% [Nguồn: Đề thi FE/RE CSD203 (Câu #70)]
% Quiz: CSD203_SU26_RE - Question #20 (ID: bf41ecbe-678c-4016-80bc-f391927c5719)
In the matrix chain multiplication problem, what determines the best order of parenthesis?
\begin{choices}
\choice Number of matrices
\choice Number of operations
\choice Size of matrices
\choice Matrix dimension matching
\end{choices}
\textbf{Đáp án:} B
\textbf{Giải thích:} Bài toán nhân chuỗi ma trận nhằm tìm cách đặt dấu ngoặc sao cho tổng số phép nhân vô hướng (number of scalar operations) là nhỏ nhất.
\end{question}

\begin{question}{8.86}
% [ID gốc: 5.114]
% [Nguồn: Đề thi FE/RE CSD203 (Câu #72)]
% Quiz: CSD203_SU26_RE - Question #22 (ID: 0aada786-910c-4900-8837-ec25ae465fad)
What is the time complexity of inserting or deleting an element at the front or rear of a Double-Ended Queue implemented with a doubly linked list?
\begin{choices}
\choice $O(1)$
\choice $O(\log n)$
\choice $O(n)$
\choice $O(n^2)$
\end{choices}
\textbf{Đáp án:} A
\textbf{Giải thích:} Khi cài đặt deque bằng danh sách liên kết đôi có con trỏ head và tail, các thao tác thêm/xóa ở hai đầu đều thực hiện trong thời gian hằng số $O(1)$.
\end{question}

\begin{question}{8.87}
% [ID gốc: 5.115]
% [Nguồn: Đề thi FE/RE CSD203 (Câu #74)]
% Quiz: CSD203_SU26_RE - Question #24 (ID: c78bb5a7-348c-4e70-b206-6b22ad4ca018)
In recursion, what is the process of breaking down a problem into smaller subproblems known as?
\begin{choices}
\choice Initialization
\choice Termination
\choice Decomposition
\choice Recombination
\end{choices}
\textbf{Đáp án:} C
\textbf{Giải thích:} Quá trình chia một bài toán lớn thành các bài toán con nhỏ hơn trong đệ quy được gọi là sự phân rã (decomposition).
\end{question}

\begin{question}{8.88}
% [ID gốc: 5.116]
% [Nguồn: Đề thi FE/RE CSD203 (Câu #76)]
% Quiz: CSD203_SU26_RE - Question #26 (ID: c0ebd313-ba78-49f3-be19-15b687af70d3)
In the recursive drawing of a ruler, how many recursive calls are made for a tick of length 3?
\begin{choices}
\choice 3
\choice 7
\choice 8
\choice 15
\end{choices}
\textbf{Đáp án:} B
\textbf{Giải thích:} Số lời gọi đệ quy để vẽ các vạch thước có độ dài n được tính theo công thức $$2^n$ - 1$. Với n = 3, số lời gọi đệ quy là 2^3 − 1 = 7.
\end{question}

\begin{question}{8.89}
% [ID gốc: 5.117]
% [Nguồn: Đề thi FE/RE CSD203 (Câu #84)]
% Quiz: CSD203_SU26_RE - Question #34 (ID: 3c8fd02d-40db-4822-affa-e82ffecc2cfb)
What is the worst-case time complexity of inserting n elements into an empty linked list, if the linked list needs to be maintained in sorted order?
\begin{choices}
\choice $O(n^2)$
\choice $O(n)$
\choice $O(2n)$
\choice $O(\log n)$
\end{choices}
\textbf{Đáp án:} A
\textbf{Giải thích:} Để chèn phần tử thứ i vào danh sách liên kết đã sắp xếp, trường hợp xấu nhất phải duyệt qua i − 1 phần tử. Tổng chi phí cho n phần tử là 1 + 2 + ... + (n − 1) = $O(n^2)$.
\end{question}

\begin{question}{8.90}
% [ID gốc: 5.118]
% [Nguồn: Đề thi FE/RE CSD203 (Câu #86)]
% Quiz: CSD203_SU26_RE - Question #36 (ID: 3fa44aaa-9a28-4615-93c5-dc1c97e250d7)
When should we use the Huffman Coding?
\begin{choices}
\choice For compression
\choice For encoding data
\choice For decoding data
\choice For searching data
\end{choices}
\textbf{Đáp án:} A
\textbf{Giải thích:} Mã hóa Huffman (Huffman coding) là thuật toán dùng cho mục đích nén dữ liệu không mất thông tin (lossless data compression).
\end{question}

\begin{question}{8.91}
% [ID gốc: 5.119]
% [Nguồn: Đề thi FE/RE CSD203 (Câu #90)]
% Quiz: CSD203_SU26_RE - Question #40 (ID: 0360a977-19ad-4fa3-8696-e7c3a66d2f30)
Which technique can be used to eliminate tail recursion?
\begin{choices}
\choice Dynamic programming
\choice Memoization
\choice Iteration
\choice Breadth-first search
\end{choices}
\textbf{Đáp án:} C
\textbf{Giải thích:} Đệ quy đuôi (tail recursion) có thể được chuyển đổi thành dạng vòng lặp (iteration) một cách dễ dàng để tối ưu bộ nhớ call stack.
\end{question}


\end{document}
